\documentclass[aps,preprint]{revtex4}%
\usepackage{amsfonts}
\usepackage{amsmath}
\usepackage{amssymb}
\usepackage{graphicx}%
\setcounter{MaxMatrixCols}{30}
%TCIDATA{OutputFilter=latex2.dll}
%TCIDATA{Version=5.50.0.2960}
%TCIDATA{CSTFile=revtex4.cst}
%TCIDATA{Created=Monday, September 20, 2004 09:58:21}
%TCIDATA{LastRevised=Friday, December 04, 2009 10:26:29}
%TCIDATA{<META NAME="GraphicsSave" CONTENT="32">}
%TCIDATA{<META NAME="SaveForMode" CONTENT="1">}
%TCIDATA{BibliographyScheme=BibTeX}
%TCIDATA{<META NAME="DocumentShell" CONTENT="Articles\SW\REVTeX 4">}
%TCIDATA{Language=American English}
%BeginMSIPreambleData
\providecommand{\U}[1]{\protect\rule{.1in}{.1in}}
%EndMSIPreambleData
\newtheorem{theorem}{Theorem}
\newtheorem{acknowledgement}[theorem]{Acknowledgement}
\newtheorem{algorithm}[theorem]{Algorithm}
\newtheorem{axiom}[theorem]{Axiom}
\newtheorem{claim}[theorem]{Claim}
\newtheorem{conclusion}[theorem]{Conclusion}
\newtheorem{condition}[theorem]{Condition}
\newtheorem{conjecture}[theorem]{Conjecture}
\newtheorem{corollary}[theorem]{Corollary}
\newtheorem{criterion}[theorem]{Criterion}
\newtheorem{definition}[theorem]{Definition}
\newtheorem{example}[theorem]{Example}
\newtheorem{exercise}[theorem]{Exercise}
\newtheorem{lemma}[theorem]{Lemma}
\newtheorem{notation}[theorem]{Notation}
\newtheorem{problem}[theorem]{Problem}
\newtheorem{proposition}[theorem]{Proposition}
\newtheorem{remark}[theorem]{Remark}
\newtheorem{solution}[theorem]{Solution}
\newtheorem{summary}[theorem]{Summary}
\newenvironment{proof}[1][Proof]{\noindent\textbf{#1.} }{\ \rule{0.5em}{0.5em}}
\begin{document}
\preprint{ }
\title[Janaks Theorem]{Janaks Theorem and Orbital Energy Functionals (in progress)}
\author{B. Weiner}
\affiliation{Department of Physics, Pennsylvania State University}
\keywords{}
\pacs{}

\begin{abstract}
Janaks theorem shows that the rate of change of the total N-electron energy
with respect to the occupation numbers of the natural orbitals of the model
Kohn-Sham (KS) independent particle state are equal to the Lagrange
multipliers associated with the normalization constraint imposed on these
orbitals. However many conceptual questions surround this theorem. In what
sense can the occupation numbers vary as the model KS state is by definition
an Independent Particle State (IPS) with integer occupation numbers? What is
the functional relationship between the eigenvalues and the occupation number?
How does one define a KS functional that is differentiable with respect to
occupation numbers? In order to examine the meaning and content of this
theorem we review the construction of the implicitly defined Hohenberg-Kohn
(HK) and KS energy functionals and obtain an expression of Janaks theorem in
the context of constrained optimization theory. It is shown that, in fact,
Janaks theorem holds for any total energy functional defined in terms of
orbitals and occupation numbers. In particular it applies to Antisymmetrized
Geminal Power energy functionals, which also solely depend on orbitals and
occupation numbers and whose functional form is explicitly known. The
properties and structure of these orbital functionals are contrasted and
compared and subtle differences are discussed..\qquad

\end{abstract}
\date{11/28/2009}
\startpage{1}
\maketitle
\tableofcontents


\section{Introduction\label{Intro}}

Janaks theorem \cite{Janak1978} is a useful auxiliary tool in the Kuhn-Sham
(KS) \cite{KS1965} approach to Density Functional Theory (DFT), that is
helpful in analyzing the structure and meaning of the eigenvalues of the KS
Fock operator. It states that these eigenvalues are equal to the derivatives
of the exact total energy with respect to the occupation numbers of the First
Order Reduced Density Operator (FORDO) of the model KS state. However many
conceptual questions surround this theorem.

\begin{itemize}
\item In what sense can the occupation numbers vary as the model KS state is
by definition an Independent Particle State (IPS)\ with integer occupation numbers?

\item What is the functional relationship between the eigenvalues and the
occupation number?

\item How does one define a KS\ functional that is differentiable with respect
to occupation numbers?
\end{itemize}

In both KS-DFT and Hartree-Fock (HF) theory these eigenvalues correspond to
Lagrange parameters associated with the normalization constraint on the
orbitals, that constitute the Independent Particle State (IPS), and are
related to the ionization energies of electrons described by\textit{ }these
orbitals in the HF case and to the highest occupied one in the KS case. This
correspondence can be generalized to Multi Configurational Self Consistent
States (MCSCF) as described, for instance, in \cite{OrtizWeiner2006}.

Constrained optimization theory \cite{Hestenes75} shows, if the Hessian of the
Lagrangian is nonsingular then the values of Lagrange parameters correspond to
the sensitivity of the object function with respect to the control parameters,
i.e. to the values of the constraints. In this case the object function is the
$N$-electron energy and the control parameters (that vary) are the norms of
the orbitals. Hence Janaks theorem implies that\emph{ the orbitals must be
normalized to their occupation numbers }and the energy is dependent on the
occupation numbers through its dependence on these orbitals. It is of course
possible to express the energy functional in terms of orthonormal orbitals and
freely varying occupation numbers, that comply to the basic constraints of
being positive with a maximum value of one, but this formulation does not
directly introduce the relationship between the variation of the energy
functional with respect to the occupation numbers and the eigenvalues of the
Fock operator.

In general orbital functionals can be classified according to whether they are

\begin{itemize}
\item invariant to local phase changes of orbitals, i.e. are gauge invariant,
thus preserve the electron density

\item invariant to global phase changes of the orbitals i.e. to the group
$U\left(  1\right)  $ $\cite{Gilmore_{b}ook}$, e.g. Density Matrix and KS\ Functionals

\item not invariant to $U\left(  1\right)  $ e.g. correlated AGP\ functionals
\cite{OrtizWeiner2006}.
\end{itemize}

In order to answer the preceding questions and analyze the form of the energy
dependence on occupation numbers all three types of orbital functionals will
be discussed.

In this paper we will review some of the basic principles and underlying
mathematical structure of Hohenburg-Kohn (HK) \cite{HK1964} and KS DFT using
the formalism of super operators \cite{Lowdin1982} (section \ref{HKDFT}). This
will allow us to express the Levy-Lieb construction of energy functionals as
an eigenvalue problem. In section \ref{GSO} we discuss orbital based
functionals and in particular we consider KSDFT and derive Janaks Theorem in
the context of a constrained variational approach and sensitivity theory again
using the super operator formalism. The KS\ density functional, density matrix
functional and the AGP functional determine the total energy in terms of a set
of occupations numbers and one-electron orbitals. The first two have only been
defined implicitly while the AGP functional has been explicitly constructed
\cite{OrtizWeiner2006} in terms of General Spin Orbitals (GSOs), which we
review in section \ref{AGP}. In this section we compare and contrast the
properties of the derivatives of this functional with respect to occupation
numbers to the derivatives obtained indirectly through sensitivity theory for
the orbital based functional introduced in section \ref{GSO}. In section
\ref{Discussion} we conclude with a discussion and summary.

\section{Hohenberg-Kohn DFT\label{HKDFT}}

In this section we review some fundamental background needed to define the HK
density functional utilizing the formalism of \emph{super operators}
\cite{Lowdin1982} applied to a constrained optimization procedure that allows
one to formulate its construction in terms of the solution of an eigenvalue
problem. Most of the important technical mathematical aspects that are needed
are examined in the appendices, but in some cases a more more detailed
analysis is deferred elsewhere.

The HK density functional is defined in terms of the kinetic energy operator,
$K,$ the coulomb operator, $V_{2},$ and the charge density $\rho$, we first
consider the properties of the charge density:

\subsection{Densities and States\label{DenStat}}

Density Functional Theory is based on the existence of a universal functional,
$\mathcal{F}\left(  \rho\right)  ,$ of the charge density, $\rho\left(
\boldsymbol{r}\right)  ,$ where
\begin{equation}
\rho\left(  \boldsymbol{r}\right)  =\int D^{1}\left(  \boldsymbol{r,\xi;r,\xi
}\right)  d\boldsymbol{\xi}%
\end{equation}
and $D^{1}\left(  \boldsymbol{r,\xi;r}^{\prime}\boldsymbol{,\xi}^{\prime
}\right)  $ is the integral kernel of the First Order Reduced Operator (FORDO)
$D^{1}$ (see Appendix \ref{Denmat} for details of density and reduced density
operators and matrices), which can be expanded as
\begin{equation}
D^{1}=\int D^{1}\left(  \boldsymbol{r,\xi;r}^{\prime}\boldsymbol{,\xi}%
^{\prime}\right)  \left\vert \boldsymbol{r,\xi}\right\rangle \left\langle
\boldsymbol{r}^{\prime}\boldsymbol{,\xi}^{\prime}\right\vert d\boldsymbol{r}%
d\boldsymbol{r}^{\prime}d\boldsymbol{\xi}d\boldsymbol{\xi}^{\prime}%
\end{equation}
in the basis $\left\{  \left\vert \boldsymbol{r,\xi}\right\rangle \right\}  $
formed by the tensor product of the spin states $\left\{  \left\vert
\boldsymbol{\xi}\right\rangle \right\}  $ and the generalized eigenvectors,
$\left\{  \left\vert \boldsymbol{r}\right\rangle \right\}  $ (see Appendix
\ref{Genbasis})$,$ of the position operator. In the proceeding we will use the
notation $\rho_{D}$ for a charge density that has been determined by a fixed
state $D$ and $\rho$ for fixed charge density that could be associated with a
set of states.

For a system that contains on average $N$-electrons the position density must
satisfy the conditions
\begin{subequations}
\label{DensityConst}%
\begin{align}
\int\rho\left(  \boldsymbol{r}\right)  d\boldsymbol{r}  &  =N\Rightarrow
\rho\in L_{1}\left(  \mathbb{R}^{3}\right) \label{number}\\
\rho\left(  \boldsymbol{r}\right)   &  \geq0\label{positivity}\\
\rho &  \in L_{3}\left(  \mathbb{R}^{3}\right)  . \label{KE}%
\end{align}
The function spaces $L_{1}\left(  \mathbb{R}^{3}\right)  $ and $L_{3}\left(
\mathbb{R}^{3}\right)  $ are defined in \cite{Lieb1983}. The constraint eq.
$\left(  \ref{number}\right)  $ does not restrict the density to be the
contraction of a $N$-electron state and one could consider incoherent and
coherent mixtures of states describing different number of electrons. However
in this article we primarily focus on $N$-electron states and mainly consider
convex sets, $\left\{  \left[  \rho\right]  _{N}\right\}  ,$ of $N$-electron
state density operators defined by
\end{subequations}
\begin{equation}
\left[  \rho\right]  _{N}=\left\{  \left.  D\right\vert \Xi_{1}\left(
D\right)  =\rho;D\in S_{N}\subset\mathcal{B}_{1}\left(  \mathcal{H}%
^{N}\right)  \right\}  ,\quad1\leq N<\infty
\end{equation}
(see Appendix \ref{Denmat} for the definition of the convex set $S_{N}$ and of
the Banach space of trace class operators $\mathcal{B}_{1}\left(
\mathcal{H}^{N}\right)  $ and eq. $\left(  \ref{LinCon}\right)  $ for the
definition of $\Xi_{1})$. The constraint eq. $\left(  \ref{KE}\right)  $
ensures that the expectation value of the kinetic energy operator is finite
\cite{Lieb1983} and the constraint eq. $\left(  \ref{positivity}\right)  $ is
a necessary and sufficient condition for representability i.e. that there
exists an electronic state that contracts to this density. The constraints
eqs. $\left(  \ref{number},\ref{KE}\text{ and }\ref{positivity}\right)  $
define a convex set $\mathcal{D}$ of densities
\begin{equation}
\rho\in\mathfrak{D}\subset L_{3+}\left(  \mathbb{R}^{3}\right)  \cap
L_{1+}\left(  \mathbb{R}^{3}\right)  , \label{Densities}%
\end{equation}
(the "+" subscript denoting positive elements) that can be contracted from
states that have a finite kinetic energy. The constraining functional $\Xi
_{1}$ can be written in terms of the set of density operators $\left\{
\left.  \Phi\left(  \boldsymbol{r}\right)  ^{\dagger}\Phi\left(
\boldsymbol{r}\right)  \right\vert \boldsymbol{r\in}\mathbb{R}^{3}\right\}  $,
(see Appendix \ref{SecQuant} for their definition) as
\begin{equation}
\left[  \Xi_{1}\left(  D\right)  \right]  \left(  \boldsymbol{r}\right)
=\operatorname{Tr}\left\{  \Phi\left(  \boldsymbol{r}\right)  ^{\dagger}%
\Phi\left(  \boldsymbol{r}\right)  D\right\}  =\rho_{D}\left(  \boldsymbol{r}%
\right)  , \label{LinCon}%
\end{equation}
the $\left[  {}\right]  $ notation emphasizes the roles of the different
arguments $D$ and $\boldsymbol{r},$ displaying that the map%

\begin{equation}
\Xi_{1}:\mathcal{B}_{1}\left(  \mathcal{H}^{N}\right)  \rightarrow
L_{1}\left(  \mathbb{R}^{3}\right)
\end{equation}
has the value $\left[  \Xi_{1}\left(  D\right)  \right]  \in$ $L_{1}\left(
\mathbb{R}^{3}\right)  $ at the point $D$ and the value $\left[  \Xi
_{1}\left(  D\right)  \right]  \left(  \boldsymbol{r}\right)
=\operatorname{Tr}\left\{  \Phi\left(  \boldsymbol{r}\right)  ^{\dagger}%
\Phi\left(  \boldsymbol{r}\right)  D\right\}  \in\mathbb{R}$ at the physical
position $\boldsymbol{r}$. In order to examine the properties of $\Xi_{1}$ it
is convenient and more useful to consider the range of $\Xi_{1}$ to be
contained in the space $L_{1}\left(  \mathbb{R}^{3}\right)  $ rather than in
$L_{3}\left(  \mathbb{R}^{3}\right)  $ or $\mathfrak{D}$ itself. The value,
$\Xi_{1}\left(  D\right)  ,$ of the function $\Xi_{1}$ at the point $D$ can be
expressed in terms of an operator valued distribution$,$ $\Phi^{\dagger}%
\Phi_{\delta},$ whose domain is restricted to the set of Dirac delta
functions, $\mathcal{D}_{\delta},$ defined on $\mathbb{R}^{3}$ with values in
the space, $\mathcal{B}\left(  \mathcal{H}^{N}\right)  ,$ of bounded operator
(see Appendices \ref{Denmat} and \ref{SecQuant}) i.e.%
\begin{align}
\Phi^{\dagger}\Phi_{\delta}  &  :\mathcal{D}_{\delta}\rightarrow
\mathcal{B}\left(  \mathcal{H}^{N}\right) \\
\Phi^{\dagger}\Phi_{\delta}\left(  \delta_{\mathbf{r}}\right)   &
=\Phi\left(  \boldsymbol{r}\right)  ^{\dagger}\Phi\left(  \boldsymbol{r}%
\right)
\end{align}
leading to
\begin{equation}
\Xi_{1}\left(  D\right)  =\operatorname{Tr}\left\{  \Phi^{\dagger}\Phi
_{\delta}D\right\}  =\rho_{D}\in L_{1}\left(  \mathbb{R}^{3}\right)  .
\end{equation}


As $D$ is a positive operator it can always be expressed as a square
\begin{equation}
D=A^{2},
\end{equation}
where
\begin{equation}
A\in\mathcal{B}_{2}\left(  \mathcal{H}^{N}\right)  \text{ and }A=A^{\dagger
}\equiv A\in\mathcal{B}_{2sa}\left(  \mathcal{H}^{N}\right)
\end{equation}
and $\mathcal{B}_{2}\left(  \mathcal{H}^{N}\right)  $ and $\mathcal{B}%
_{2sa}\left(  \mathcal{H}^{N}\right)  $ are the Hilbert spaces of
Hilbert-Schmidt and self adjoint Hilbert-Schmidt operators respectively acting
on the $N$-electron space $\mathcal{H}^{N}$ (see appendix \ref{Denmat} for the
definitions of $\mathcal{B}_{2}\left(  \mathcal{H}^{N}\right)  $ and
$\mathcal{B}_{2sa}\left(  \mathcal{H}^{N}\right)  $). The operator $A$ is the
square root of $D$, which is not unique, as can be observed from the spectral
expansion
\begin{equation}
A=D^{\frac{1}{2}}=\sum_{1\leq j\leq m}\pm\sqrt{\alpha_{j}}\mathcal{P}_{j},
\label{specA}%
\end{equation}
where
\begin{equation}
D=\sum_{1\leq j\leq m}\alpha_{j}\mathcal{P}_{j};\,\alpha_{j}\geq0,1\leq j\leq
m,
\end{equation}
$r_{N}=\sum\limits_{1\leq j\leq m}\nu_{j}=$ dimension of $\mathcal{H}^{N}$,
$\left\{  \left.  \nu_{j}\right\vert 1\leq j\leq m\right\}  $ the degeneracies
of the eigenvalues $\left\{  \left.  \alpha_{j}\right\vert 1\leq j\leq
m\right\}  $ and $\left\{  \left.  \mathcal{P}_{j}\right\vert 1\leq j\leq
m\right\}  $ are orthogonal projectors onto the eigenspaces of $D$. If one
restricts $A$ to be positive, i.e. uses the positive square roots in eq.
$\left(  \ref{specA}\right)  $, then it belongs to the cone, $\mathcal{B}%
_{2sa+}\left(  \mathcal{H}^{N}\right)  ,$ \cite{Encyclo} of positive
Hilbert-Schmidt operators and uniquely determines $D$. However, in general it
not convenient to restrict $A$ to be positive and one allows the domain of
variation of the $N$-electron energy functional to be contained in the
\emph{real} Hilbert space $\mathcal{B}_{2sa}\left(  \mathcal{H}^{N}\right)  ,$
of dimension $r_{N}^{2},$ which has an orthonormal bases
\begin{align}
\Lambda_{jk}  &  =\frac{1}{\sqrt{2}}\left\{  \left\vert \Psi_{j}\right\rangle
\left\langle \Psi_{k}\right\vert +\left\vert \Psi_{k}\right\rangle
\left\langle \Psi_{j}\right\vert \right\}  ;\quad1\leq j<k\leq r_{N}%
\nonumber\\
\Lambda_{jj}  &  =\left\vert \Psi_{j}\right\rangle \left\langle \Psi
_{j}\right\vert ;\quad1\leq j\leq r_{N}\nonumber\\
\Lambda_{jk}  &  =\frac{i}{\sqrt{2}}\left\{  \left\vert \Psi_{j}\right\rangle
\left\langle \Psi_{k}\right\vert -\left\vert \Psi_{k}\right\rangle
\left\langle \Psi_{j}\right\vert \right\}  ;\quad1\leq j>k\leq r_{N},
\end{align}
where $\left\{  \left.  \left\vert \Psi_{j}\right\rangle \right\vert 1\leq
j\leq r_{N}\right\}  $ is any complete orthonormal basis for $\mathcal{H}%
^{N},$ of course all these spaces are infinite dimensional if $r_{N}=\infty.$

\subsection{The Density Constraint\label{DenCon}}

The Levy-Leib construction \cite{Levy1979,Lieb1983} is based on defining the
value of the density functional as the minimum of an energy variation over
states constrained to have a fixed density $\rho\in\mathfrak{D}$. The density
constraint can be formulated in terms of a quadratic functional $\Xi_{2},$
based on the linear functional $\Xi_{1}$ eq. $\left(  \ref{LinCon}\right)  ,$
acting on the space of self adjoint Hilbert-Schmidt operators, $\mathcal{B}%
_{2sa}\left(  \mathcal{H}^{N}\right)  $, where
\begin{align}
\Xi_{2}  &  :\mathcal{B}_{2sa}\left(  \mathcal{H}^{N}\right)  \rightarrow
\mathfrak{D}\subset L_{3+}\left(  \mathbb{R}^{3}\right)  \cap L_{1+}\left(
\mathbb{R}^{3}\right)  \subset L_{1}\left(  \mathbb{R}^{3}\right) \nonumber\\
\Xi_{2}\left(  A\right)   &  =\Xi_{1}\left(  A^{2}\right)  =\operatorname{Tr}%
\left\{  A\Phi^{\dagger}\Phi_{\delta}A\right\}  =\left(  A\left\vert
\Phi^{\dagger}\Phi_{\delta}A\right.  \right)  =\rho, \label{QuadCon}%
\end{align}
where the real symmetric inner product $\left(  \cdot\left\vert \cdot\right.
\right)  $ in $\mathcal{B}_{2sa}\left(  \mathcal{H}^{N}\right)  $ is defined
as
\begin{equation}
\left(  X\left\vert Y\right.  \right)  =\operatorname{Tr}\left\{  XY\right\}
. \label{Opinner}%
\end{equation}
The constraint eq. $\left(  \ref{QuadCon}\right)  $ can be formulated in terms
of the \emph{super operator valued distribution} $\widehat{\Phi^{\dagger}%
\Phi_{\delta}},$ as%
\begin{equation}
\Xi_{2}\left(  A\right)  =\left(  A\left\vert \widehat{\Phi^{\dagger}%
\Phi_{\delta}}A\right.  \right)  =\rho,
\end{equation}
where
\begin{equation}
\widehat{\Phi^{\dagger}\Phi_{\delta}}:\mathcal{D}_{\mathcal{\delta}%
}\rightarrow\mathcal{B}\left(  \mathcal{B}_{2sa}\left(  \mathcal{H}%
^{N}\right)  \right)  ,
\end{equation}
and has components defined as (Appendix \ref{SecQuant})%
\begin{equation}
\widehat{\Phi^{\dagger}\Phi_{\delta}}\left(  \boldsymbol{r}\right)  \left\vert
X\right)  =\widehat{\Phi\left(  \boldsymbol{r}\right)  ^{\dagger}\Phi\left(
\boldsymbol{r}\right)  }\left\vert X\right)  =\left\vert \left[  \Phi\left(
\boldsymbol{r}\right)  ^{\dagger}\Phi\left(  \boldsymbol{r}\right)  X\right]
_{+}\right)  \text{ }\forall X\in\mathcal{B}_{2sa}\left(  \mathcal{H}%
^{N}\right)  ,\boldsymbol{r\in}\mathbb{R}^{3},
\end{equation}
where
\begin{equation}
\left[  \Phi\left(  \boldsymbol{r}\right)  ^{\dagger}\Phi\left(
\boldsymbol{r}\right)  X\right]  _{+}=\frac{1}{2}\left[  \Phi\left(
\boldsymbol{r}\right)  ^{\dagger}\Phi\left(  \boldsymbol{r}\right)  ,X\right]
_{+}.
\end{equation}
The positive commutator $\left[  ,\right]  _{+}$ is used in the definition of
the super operator in order to ensure that its action produces a self adjoint operator.

The functional $\Xi_{2}\left(  A\right)  $ is Fr\'{e}chet differentiable and
one can form the differential, $d\Xi_{2}\left(  A\right)  ,$ \cite{Encyclo},
(Appendix \ref{DerivDef})
\begin{equation}
d\Xi_{2}\left(  A\right)  :\mathcal{B}_{2sa}\left(  \mathcal{H}^{N}\right)
\rightarrow L_{1}\left(  \boldsymbol{\Delta}_{\rho}\right)  \subseteq
L_{1}\left(  \mathbb{R}^{3}\right)
\end{equation}
of $\Xi_{2}$ at any point $A\in\mathcal{B}_{2sa}\left(  \mathcal{H}%
^{N}\right)  ,$ where $\boldsymbol{\Delta}_{A}$ is the closure of the support
of the density function i.e.
\begin{equation}
\boldsymbol{\Delta}_{\rho}=\operatorname*{supp}\rho_{A}=\overline{\left\{
\left.  \mathbf{r}\right\vert \rho\left(  \mathbf{r}\right)  \neq
0;\mathbf{r\in}\mathbb{R}\right\}  },
\end{equation}
in terms of the components
\begin{equation}
\left[  d\Xi_{2}\left(  A\right)  \right]  \left(  \boldsymbol{r}\right)
=\left(  \left[  \Phi\left(  \boldsymbol{r}\right)  ^{\dagger}\Phi\left(
\boldsymbol{r}\right)  A\right]  _{+}\right\vert .
\end{equation}
The set of operators $\left\{  \left.  \left\vert \left[  \Phi\left(
\boldsymbol{r}\right)  ^{\dagger}\Phi\left(  \boldsymbol{r}\right)  A\right]
_{+}\right)  \right\vert \boldsymbol{r\in}\mathbb{R}^{3}\right\}  $ forms a
generalized basis $\left(  \text{Appendix \ref{Genbasis}}\right)  $ for the
normal spaces of the surface $\Xi_{2}\left(  A\right)  =\rho$ for all points
$A$ $\left(  \text{Appendix \ref{Regular}}\right)  ,$ which shows that
$d\Xi_{2}\left(  A\right)  $ is surjective \cite{Encyclo} onto $L_{1}\left(
\boldsymbol{\Delta}_{\rho}\right)  $ and thus the constraint $\Xi_{2}\left(
A\right)  =\rho$ is \emph{normal} \cite{Hestenes75}, (Appendix (\ref{Regular}%
)) everywhere. Hence the largrange mutltiplier function for this constraint is unique.

\subsection{Hohenberg-Kohn Density Functional Construction}

Even though the explicit form of the Hohenberg-Kohn functional, $\mathcal{F}%
_{HK}\left(  \rho\right)  ,$ is not known one can define it implicitly
\cite{Levy1979,Lieb1983} by
\begin{equation}
\mathcal{F}_{HK}\left(  \rho\right)  =\operatorname{Tr}\left\{  \left(
K+V_{2}\right)  A_{\ast}\left(  \rho\right)  ^{2}\right\}  =E_{U}\left(
A_{\ast}\left(  \rho\right)  \right)
\end{equation}
in terms of the minima, $E_{U}\left(  A_{\ast}\left(  \rho\right)  \right)  ,$
of the universal energy functional $E_{U}\left(  A\right)  ,$ that are
solutions of the constrained optimization problem
\begin{subequations}
\label{HKQ}%
\begin{align}
&  E_{U}\left(  A_{\ast}\left(  \rho\right)  \right)  =\min_{A\in
\mathcal{B}_{2sa}\left(  \mathcal{H}^{N}\right)  }\operatorname{Tr}\left\{
A\overset{}{\left(  K+V_{2}\right)  A}\right\}  =\left(  A_{\ast}\left(
\rho\right)  \left\vert \left(  \widehat{K}+\widehat{V_{2}}\right)  A_{\ast
}\left(  \rho\right)  \right.  \right)  \label{HKQa}\\
\qquad &  \qquad\qquad\qquad\text{subject to }\Xi_{2}\left(  A\right)
=\left(  A\left\vert \widehat{\Phi^{\dagger}\Phi}A\right.  \right)
=\rho,\label{HKQb}%
\end{align}
where $K,$ and $V_{2}$ are the kinetic energy and coulomb operators, the
corresponding super operators $\widehat{K\text{ }}$ and $\widehat{V_{2}}$ are
defined by
\end{subequations}
\begin{equation}
\widehat{S}\left\vert X\right)  =\left\vert \left[  SX\right]  _{+}\right)
\text{ }\forall X\in\mathcal{B}_{2sa}\left(  \mathcal{H}^{N}\right)  ,\quad
S=K,V_{2}%
\end{equation}
and the control densities $\rho\in\mathcal{D}.$ This choice of the density
$\rho$ coupled with insisting that $A$ is an $N$-electron operator, i.e.
$A\in\mathcal{B}_{2sa}\left(  \mathcal{H}^{N}\right)  ,$ ensures that $A$ is
normalized (Appendix \ref{Normalization}$),$ i.e.
\begin{equation}
\left(  A\left\vert A\right.  \right)  =1.
\end{equation}
The Lagrangian associated with this problem eq. $\left(  \ref{HKQ}\right)  $
is
\begin{equation}
\mathcal{L}\left(  A,\lambda,\rho\right)  =E_{U}\left(  A\right)
-\left\langle \left.  \lambda\right\vert \left\{  \Xi_{2}\left(  A\right)
-\rho\right\}  \right\rangle ,\label{Lar}%
\end{equation}
where the universal energy is
\begin{equation}
E_{U}\left(  A\right)  =\operatorname{Tr}\left\{  \left(  K+V_{2}\right)
A^{2}\right\}  =\left(  A\left\vert \left(  \widehat{K}+\widehat{V_{2}%
}\right)  A\right.  \right)  .
\end{equation}
$\lceil K+V_{2}$ is bounded below, hence $\left(  \widehat{K}+\widehat{V_{2}%
}\right)  $ is also bounded below$\rfloor.$ The Lagrange multiplier term can
be expanded as
\begin{equation}
\left\langle \left.  \lambda\right\vert \left\{  \Xi_{2}\left(  A\right)
-\rho\right\}  \right\rangle =\int\lambda\left(  \boldsymbol{r}\right)
\left\{  \Xi_{2}\left(  A\right)  \left(  \boldsymbol{r}\right)  -\rho\left(
\boldsymbol{r}\right)  \right\}  d\boldsymbol{r},\boldsymbol{\,}\lambda\in
L_{\infty}\left(  \mathbb{R}^{3}\right)  \oplus L_{\frac{3}{2}}\left(
\mathbb{R}^{3}\right)  ,\rho\in\mathcal{D},
\end{equation}
where $\left\langle \cdot|\cdot\right\rangle $ is the scalar product between
the dual spaces $\left\langle L_{\infty}\left(  \mathbb{R}^{3}\right)  \oplus
L_{\frac{3}{2}}\left(  \mathbb{R}^{3}\right)  ,L_{1}\left(  \mathbb{R}%
^{3}\right)  \cap L_{3}\left(  \mathbb{R}^{3}\right)  \right\rangle $ i.e.
\begin{equation}
\left\langle \cdot|\cdot\right\rangle :\left(  L_{\infty}\left(
\mathbb{R}^{3}\right)  \oplus L_{\frac{3}{2}}\left(  \mathbb{R}^{3}\right)
\right)  \times\left(  L_{1}\left(  \mathbb{R}^{3}\right)  \cap L_{3}\left(
\mathbb{R}^{3}\right)  \right)  \rightarrow\mathbb{R}.
\end{equation}
This constraint term in the Lagrangian eq. $\left(  \ref{Lar}\right)  $ can
also be expressed in the form
\begin{equation}
\left\langle \left.  \lambda\right\vert \left\{  \Xi_{2}\left(  A\right)
-\rho\right\}  \right\rangle =\left\langle \lambda\left\vert \rho\right.
\right\rangle -\left(  A\left\vert \widehat{\Omega\left(  \lambda\right)
}\right.  A\right)  =\left(  A\left\vert \left\{  \left\langle \lambda
\left\vert \rho\right.  \right\rangle \hat{I}-\widehat{\Omega\left(
\lambda\right)  }\right\}  \right.  A\right)  ,
\end{equation}
where $\hat{I}$ is the super operator identity operator, $\Omega\left(
\lambda\right)  $ is a local one body potential defined as
\begin{equation}
\Omega\left(  \lambda\right)  =\int\lambda\left(  \boldsymbol{r}\right)
\boldsymbol{\Phi}\left(  \boldsymbol{r}\right)  ^{\dagger}\boldsymbol{\Phi
}\left(  \boldsymbol{r}\right)  d\boldsymbol{r}%
\end{equation}
and the associated super operator expectation values are given by
\begin{equation}
\left(  A\left\vert \widehat{\Omega\left(  \lambda\right)  }\right.  A\right)
=\int\lambda\left(  \boldsymbol{r}\right)  \Xi_{2}\left(  A\right)  \left(
\boldsymbol{r}\right)  d\boldsymbol{r.}%
\end{equation}
The Lagrangian for a fixed $N$-particle density $\rho$ can then be recast as
\begin{equation}
\mathcal{L}\left(  A,\lambda,\rho\right)  =\left(  A\left\vert \left(
\widehat{K}+\widehat{V_{2}}-\widehat{\Omega\left(  \lambda\right)  }\right)
A\right.  \right)  +\left\langle \lambda\left\vert \rho\right.  \right\rangle
.
\end{equation}


\subsection{Optimization Conditions\label{OpCon}}

Solutions, $A_{\ast},$ of the constrained optimization eq. $\left(
\ref{HKQ}\right)  $ correspond to unconstrained local saddle points $\left(
A_{\ast},\lambda_{\ast}\right)  $ \cite{Hestenes75} of the Lagrangian
$\mathcal{L}\left(  A,\lambda,\rho\right)  $ such that%
\begin{equation}
\mathcal{L}\left(  A_{\ast},\lambda_{\ast},\rho\right)  =\max_{\lambda}%
\min_{A}\mathcal{L}\left(  A,\lambda,\rho\right)  . \label{MinProb}%
\end{equation}
These saddle points satisfy the necessary condition
\begin{equation}
d\mathcal{L}\left(  A_{\ast},\lambda_{\ast},\rho\right)  =0, \label{NecCon}%
\end{equation}
where the linear map for a fixed $\rho$
\begin{equation}
d\mathcal{L}\left(  A,\lambda,\rho\right)  :\mathcal{B}_{2sa}\left(
\mathcal{H}^{N}\right)  \times\left(  L_{\infty}\left(  \mathbb{R}^{3}\right)
\oplus L_{\frac{3}{2}}\left(  \mathbb{R}^{3}\right)  \right)  \rightarrow
\mathbb{R}%
\end{equation}
is given by%
\begin{equation}
d\mathcal{L}\left(  A,\lambda,\rho\right)  \left(  h,\varphi\right)
=\left\langle \varphi\left\vert \rho\right.  \right\rangle -\left(
A\left\vert \widehat{\Omega\left(  \varphi\right)  }\right.  A\right)
+2\left(  h\left\vert \left\{  \widehat{K}+\widehat{V_{2}}-\widehat{\Omega
\left(  \lambda\right)  }\right\}  \right.  A\right)  . \label{LagDiff}%
\end{equation}
The differential eq. $\left(  \ref{LagDiff}\right)  $ can be partitioned into
$\left(  d_{A}\mathcal{L}\left(  A,\lambda,\rho\right)  ,d_{\lambda
}\mathcal{L}\left(  A,\lambda,\rho\right)  \right)  $, where
\begin{subequations}
\label{differential}%
\begin{align}
d_{A}\mathcal{L}\left(  A,\lambda,\rho\right)  \left(  \cdot\right)   &
=2\left(  \cdot\left\vert \left\{  \widehat{K}+\widehat{V_{2}}-\widehat{\Omega
\left(  \lambda\right)  }\right\}  \right.  A\right) \\
d_{\lambda}\mathcal{L}\left(  A,\lambda,\rho\right)  \left(  \cdot\right)   &
=\left\langle \cdot\left\vert \rho\right.  \right\rangle -\left(  A\left\vert
\Omega\left(  \cdot\right)  \right.  A\right)  ,
\end{align}
thus giving the necessary stationary point conditions
\end{subequations}
\begin{subequations}
\begin{align}
\left(  h\left\vert \left\{  \widehat{K}+\widehat{V_{2}}-\widehat{\Omega
\left(  \lambda_{\ast}\right)  }\right\}  \right.  A_{\ast}\right)   &
=0\,\forall h\in\mathcal{B}_{2sa}\left(  \mathcal{H}^{N}\right) \label{stata}%
\\
\left\langle \varphi\left\vert \rho\right.  \right\rangle -\left(  A_{\ast
}\left\vert \widehat{\Omega\left(  \varphi\right)  }\right.  A_{\ast}\right)
&  =0\,\forall\varphi\in L_{\infty}\left(  \mathbb{R}^{3}\right)  \oplus
L_{\frac{3}{2}}\left(  \mathbb{R}^{3}\right)  . \label{statb}%
\end{align}
Necessary and sufficent conditions for $A_{\ast}$ to be a constrained minimum
of eq. $\left(  \ref{HKO}\right)  $ are that
\end{subequations}
\begin{equation}
d\mathcal{L}\left(  A_{\ast},\lambda_{\ast},\rho\right)  =0,
\end{equation}
and
\begin{equation}
d_{AA}^{2}\mathcal{L}\left(  A_{\ast},\lambda_{\ast},\rho\right)  \geq
0\forall\delta A\in\ker\left\{  d\Xi_{2}\left(  A_{\ast}\right)  \right\}  ,
\end{equation}
where
\begin{equation}
d_{AA}^{2}\mathcal{L}\left(  A_{\ast},\lambda_{\ast},\rho\right)  =2\left(
\cdot\left\vert \left\{  \widehat{K}+\widehat{V_{2}}-\widehat{\Omega\left(
\lambda\right)  }\right\}  \right.  \cdot\right)
\end{equation}
i.e.
\begin{equation}
\widehat{P}\left(  \widehat{K}+\widehat{V_{2}}-\widehat{\Omega\left(
\lambda\right)  }\right)  \widehat{P}\geq0, \label{PHess}%
\end{equation}
where $\widehat{P}$ is the orthogonal projector onto $\ker\left\{  d\Xi
_{2}\left(  A_{\ast}\right)  \right\}  .$

Harrimans Theorem \cite{Harriman81} ensures that there are solutions to eq.
$\left(  \ref{statb}\right)  .$ Further if there exist $\lambda_{\ast}$ such
that there exist simultaneous solutions, $A_{\ast},$ to both eq. $\left(
\ref{statb}\right)  $ and eq. $\left(  \ref{stata}\right)  $ then these
$\lambda_{\ast}$ are unique, as the constraint is normal \cite{Hestenes75},
but there may be many solutions, $A_{\ast}$'s, for a given $\lambda_{\ast}$
i.e. degenerate $A_{\ast}$'s. In such cases the restriction of the Hessian to
$\ker\left\{  d\Xi_{2}\left(  A_{\ast}\right)  \right\}  $ shown in eq.
$\left(  \ref{PHess}\right)  $ is not positive definite i.e. has some zero
eigenvalues. ********

If eq. $\left(  \ref{stata}\right)  $ and eq. $\left(  \ref{statb}\right)  $
are true for all operators $h$ bq.elonging to $\mathcal{B}_{2sa}\left(
\mathcal{H}^{N}\right)  $ and all functions $\varphi$ belonging to $L_{\infty
}\left(  \mathbb{R}^{3}\right)  \oplus L_{\frac{3}{2}}\left(  \mathbb{R}%
^{3}\right)  $ these conditions are equivalent to
\begin{equation}
\left\{  \widehat{K}+\widehat{V_{2}}-\widehat{\Omega\left(  \lambda_{\ast
}\right)  }\right\}  \left\vert A_{\ast}\right)  =0 \label{Stateig}%
\end{equation}
and the constraining condition
\begin{equation}
\left(  A_{\ast}\left\vert \widehat{\Phi^{\dagger}\Phi}\right.  A_{\ast
}\right)  =\rho,
\end{equation}
which imply that $\widehat{K}+\widehat{V_{2}}-\widehat{\Omega\left(
\lambda_{\ast}\right)  }$ has a non empty point spectra. If no $A_{\ast}$
exists (i.e. no $\lambda_{\ast}$ exists) for a given $\rho$ then $\rho$ is not V-representable.

$\lceil$%
\[
\left\{  \widehat{K}+\widehat{V_{2}}-\widehat{\Omega\left(  \lambda_{\ast
}\right)  }\right\}  \left\vert A_{\ast}\right)  =0\Rightarrow
\operatorname{Tr}\left\{  h\left\{  K+V_{2}-\Omega\left(  \lambda_{\ast
}\right)  \right\}  A_{\ast}\right\}  =0\forall h
\]
$\rfloor$

The optimal value of the Lagrangian satisfies the following identities:
\begin{align}
\mathcal{L}\left(  A_{\ast},\lambda_{\ast},\rho\right)   &  =\left(  A_{\ast
}\left(  \rho\right)  \left\vert \left\{  \left\{  \widehat{K}+\widehat{V_{2}%
}-\widehat{\Omega\left(  \lambda_{\ast}\right)  }\right\}  +\left\langle
\left.  \lambda_{\ast}\left(  \rho\right)  \right\vert \rho\right\rangle
\hat{I}\right\}  \right.  A_{\ast}\left(  \rho\right)  \right)  =\left\langle
\left.  \lambda_{\ast}\left(  \rho\right)  \right\vert \rho\right\rangle
\nonumber\\
&  =\left(  A_{\ast}\left(  \rho\right)  \left\vert \left\{  \widehat{K}%
+\widehat{V_{2}}\right\}  \right.  A_{\ast}\left(  \rho\right)  \right)
=E_{U}\left(  A_{\ast}\left(  \rho\right)  \right)  .
\end{align}
This shows that the optimal values of the universal energy functional, the
value of the Lagrangian and the value of the HK energy functional are all
equal at the density $\rho$ i.e.
\begin{equation}
\mathcal{F}_{HK}\left(  \rho\right)  =E_{U}\left(  A_{\ast}\left(
\rho\right)  \right)  =\left\langle \left.  \lambda_{\ast}\left(  \rho\right)
\right\vert \rho\right\rangle .
\end{equation}
Moreover $\rho$ is the density of the ground state $D_{\ast}\left(
\rho\right)  =A_{\ast}\left(  \rho\right)  ^{2}$ of the Hamiltonian
\begin{equation}
K+V_{2}-\Omega\left(  \lambda_{\ast}\left(  \rho\right)  \right)
\end{equation}
i.e. of an electronic system subject to an external potential $\lambda_{\ast
}\left(  \rho\right)  $. Note that the ground state might be a ensemble state.

\subsection{Sensitivity Analysis\label{SenAnal}}

Sensitivity analysis \cite{Hestenes75} examines how the optimal solution
changes with variation of the value of the constraining functions, which can
be thought of as control parameters. In this case the $N$-electron state
energy with variations of $\rho.$ If the optimal points $\left\{  A_{\ast
},\lambda_{\ast},\rho\right\}  $ are regular (Appendix \ref{Regular}) for a
set of $\rho^{\prime}s$ then $\left\{  A_{\ast},\lambda_{\ast}\right\}  $ are
in fact functions of $\rho,$ i.e.
\begin{align}
A_{\ast}  &  =A_{\ast}\left(  \rho\right) \nonumber\\
\lambda_{\ast}  &  =\lambda_{\ast}\left(  \rho\right)  ,
\end{align}
and thus single valued i.e. the ground state is non degenerate. The
sensitivity of the optimal $N$-electron energy to changes in the density
$\rho$ is then given by
\begin{equation}
\frac{\delta E_{U}\left(  A_{\ast}\left(  \rho\right)  \right)  }{\delta\rho
}=\lambda_{\ast}\left(  \rho\right)  .
\end{equation}
By construction
\begin{equation}
\mathcal{F}_{HK}\left(  \rho\right)  =E_{U}\left(  A_{\ast}\left(
\rho\right)  \right)  ,
\end{equation}
which shows
\begin{equation}
\frac{\delta\mathcal{F}_{HK}\left(  \rho\right)  }{\delta\rho}=\lambda_{\ast
}\left(  \rho\right)
\end{equation}
i.e. that the external potential is equal to the functional derivative of the
HK functional with respect to the density.

\subsection{Decomposition of the Hohenburg-Kohn Functional\label{DecompHK}}

If a minimizer $D_{\ast}\left(  \rho\right)  $ exists and is unique and then
one can define kinetic energy, coulomb and exchange-correlation functionals
as
\begin{align}
\mathcal{F}_{K}\left(  \rho\right)   &  =\operatorname{Tr}\left\{  KD_{\ast
}\left(  \rho\right)  \right\}  ,\nonumber\\
\mathcal{F}_{C}\left(  \rho\right)   &  =\int\frac{\rho\left(
\boldsymbol{r,\xi}\right)  \rho\left(  \boldsymbol{r}^{\prime}\boldsymbol{,\xi
}\right)  }{\left\Vert \boldsymbol{r-r}^{\prime}\right\Vert }d\boldsymbol{r}%
d\boldsymbol{r}^{\prime}d\boldsymbol{\xi},\nonumber\\
\mathcal{F}_{XC}\left(  \rho\right)   &  =\operatorname{Tr}\left\{
V_{2}D_{\ast}\left(  \rho\right)  \right\}  -\mathcal{F}_{C}\left(
\rho\right)  ,
\end{align}
producing the decomposition
\begin{equation}
\mathcal{F}_{HK}\left(  \rho\right)  =\mathcal{F}_{K}\left(  \rho\right)
+\mathcal{F}_{C}\left(  \rho\right)  +\mathcal{F}_{XC}\left(  \rho\right)  .
\end{equation}
The concept of a unique exchange-correlation functional is central to the
theory, but the existence of non isolated degenerate $D_{\ast}\left(
\rho\right)  ^{\prime}$s could give different $\mathcal{F}_{XC}\left(
\rho\right)  ^{\prime}$s and $\mathcal{F}_{K}\left(  \rho\right)  ^{\prime}$s.
However if degeneracies do occur one can still define an exchange-correlation
functional through the imposition of further constrained minimizations.

\subsection{System Energy\label{MolSysEn}}

The total energy functional of a system that contains $N$ electrons and
experiences an environment described by a \emph{local} potential $\mathrm{v}$
is
\begin{equation}
E\left(  \rho\boldsymbol{,}\mathrm{v}\right)  =\mathcal{F}_{HK}\left(
\rho\right)  +f_{\rho}\left(  \mathrm{v}\right)  ,
\end{equation}
where the explicit effect of the environment is described by the functional
\begin{equation}
f_{\rho}\left(  \mathrm{v}\right)  =\int\mathrm{v}\left(  \boldsymbol{r}%
\right)  \rho\left(  \boldsymbol{r}\right)  d\boldsymbol{r.}%
\end{equation}
The ground state energy $E_{GS}\left(  \mathrm{v}\right)  $ and the
corresponding density $\rho_{GS}$ $\left(  \mathrm{v}\right)  $ are then
determined by the solution to the constrained optimization problem
\begin{equation}
E_{GS}\left(  \mathrm{v}\right)  =E\left(  \rho_{GS}\boldsymbol{,}%
\mathrm{v}\right)  =\min_{\rho\in\mathcal{D}}E\left(  \rho\boldsymbol{,}%
\mathrm{v}\right)  , \label{GS}%
\end{equation}
which one can make unconstrained in terms of real functions $f\in H_{1}\left(
\mathbb{R}^{3}\right)  $ \cite{Lieb1983} by setting
\begin{equation}
\rho\left(  \boldsymbol{r}\right)  =\frac{Nf\left(  \boldsymbol{r}\right)
^{2}}{\left\langle f|f\right\rangle }=\frac{N\left\langle f\left\vert
\Phi\left(  \boldsymbol{r}\right)  ^{\dagger}\Phi\left(  \boldsymbol{r}%
\right)  f\right.  \right\rangle }{\left\langle f|f\right\rangle },\,f\in
H_{1}\left(  \mathbb{R}^{3}\right)  . \label{rho_f}%
\end{equation}
The minimizer $f_{GS}$ of $E\left(  f\boldsymbol{,}\mathrm{v}\right)  $ then
satisfies
\begin{equation}
\frac{\delta E\left(  f_{GS}\boldsymbol{,}\mathrm{v}\right)  }{\delta f\left(
\boldsymbol{r}\right)  }=2f_{GS}\left(  \boldsymbol{r}\right)  \left\{
\frac{N}{\left\langle f_{GS}|f_{GS}\right\rangle }\frac{\delta E\left(
f_{GS}\boldsymbol{,}\mathrm{v}\right)  }{\delta\rho\left(  \boldsymbol{r}%
\right)  }-\int\frac{\delta E\left(  f_{GS}\boldsymbol{,}\mathrm{v}\right)
}{\delta\rho\left(  \boldsymbol{r}^{\prime}\right)  }\rho_{GS}\left(
\boldsymbol{r}^{\prime}\right)  d\boldsymbol{r}^{\prime}\right\}
\overset{a.e.}{=}0
\end{equation}
which implies that
\begin{equation}
f_{GS}\left(  \boldsymbol{r}\right)  \overset{a.e.}{=}0\Leftrightarrow
\rho_{GS}\left(  \boldsymbol{r}\right)  \overset{a.e.}{=}0
\end{equation}
or
\begin{equation}
\frac{\delta E\left(  f_{GS}\boldsymbol{,}\mathrm{v}\right)  }{\delta
\rho\left(  \boldsymbol{r}\right)  }\overset{a.e.}{=}\alpha\left(  \rho
_{GS}\right)  , \label{fmin}%
\end{equation}
where%
\begin{equation}
\alpha\left(  \rho_{GS}\right)  =\frac{\left\langle f_{GS}|f_{GS}\right\rangle
}{N}\int\frac{\delta E\left(  f_{GS}\boldsymbol{,}\mathrm{v}\right)  }%
{\delta\rho\left(  \boldsymbol{r}^{\prime}\right)  }\rho_{GS}\left(
\boldsymbol{r}^{\prime}\right)  d\boldsymbol{r}^{\prime}. \label{alpha}%
\end{equation}
Hence
\begin{align}
\alpha\left(  \rho_{GS}\right)   &  =\frac{\left\langle f_{GS}|f_{GS}%
\right\rangle }{N}\alpha\left(  \rho_{GS}\right)  \int\rho_{GS}\left(
\boldsymbol{r}^{\prime}\right)  d\boldsymbol{r}^{\prime}=\frac{\left\langle
f_{GS}|f_{GS}\right\rangle }{N}\alpha\left(  \rho_{GS}\right) \\
&  \Rightarrow\alpha\left(  \rho_{GS}\right)  =0\text{ or }f_{GS}\left(
\boldsymbol{r}\right)  \overset{a.e.}{=}0.
\end{align}
So as $f_{GS}\left(  \boldsymbol{r}\right)  \overset{a.e.}{\neq}0$ we obtain
the well known result
\begin{equation}
\frac{\delta E\left(  f_{GS}\boldsymbol{,}\mathrm{v}\right)  }{\delta
\rho\left(  \boldsymbol{r}\right)  }\overset{a.e.}{=}0\Leftrightarrow
\frac{\delta\mathcal{F}_{HK}\left(  \rho_{GS}\right)  }{\delta\rho\left(
\boldsymbol{r}\right)  }\overset{a.e.}{=}-\mathrm{v}\left(  \boldsymbol{r}%
\right)  .
\end{equation}


\section{Orbital Variables \label{GSO}}

There are two standard approaches to formulating the total energy in terms of
orbitals and occupation numbers, one is to consider functionals of FORDO's,
which leads to DMFT, the other is to factor the dependence through the
density, which leads to KS-DFT and is based on the HK functional. However as
the explicit forms of these functionals are not known one must rely on the
implicit definitions via constrained optimizations, though in practice
approximate explicit forms are utilized. We have recently
\cite{OrtizWeiner2006} introduced a third approach which is defined on states
generated by the submanifold of $2M$-electron AGP states. The total energy
functional, which is explicitly given, \emph{is more general }than a density
matrix functional while still expressing the total energy in terms of orbitals
and occupation numbers.

The central focus of this article is on functionals of the density that
contain no explicit reference to the spin properties of the system. However
one could widen this focus by considering spin and general spin densities that
are described by vectors $\boldsymbol{\rho}\left(  \boldsymbol{r}\right)
=\left(  \rho_{\alpha}\left(  \boldsymbol{r}\right)  ,\rho_{\beta}\left(
\boldsymbol{r}\right)  \right)  $ and $\boldsymbol{\rho}\left(  \boldsymbol{r}%
\right)  =\left(  \rho_{\alpha}\left(  \boldsymbol{r}\right)  ,\rho
_{\alpha\beta}\left(  \boldsymbol{r}\right)  ,\rho_{\beta}\left(
\boldsymbol{r}\right)  \right)  $ respectively and modify the preceding
analysis in the appropriate fashion (Appendix \ref{Spin}).

\subsection{Correspondence between Orbitals and Densities\label{Orb_Den}}

Any density $\rho$ can expressed in terms of the Natural Orbitals (NO's)
$\left\{  \psi_{j}\right\}  $ and occupation numbers $\left\{  N_{j}\right\}
$ of the FORDO $D_{\ast}^{1}\left(  \rho\right)  $ (Appendix \ref{Denmat}) as%

\begin{equation}
\rho\left(  \boldsymbol{r}\right)  =\left[  D_{\ast}^{1}\left(  \rho\right)
\right]  \left(  \boldsymbol{r,r}\right)  =\sum_{1\leq j\leq r}N_{j}\left\vert
\psi_{j}\left(  \boldsymbol{r}\right)  \right\vert ^{2}, \label{NGSOexpan}%
\end{equation}
but many other sets of orbitals also produce the same density, in particular
Harimanns \cite{Harriman81} theorem ensures that there are always expansions
of the form
\begin{equation}
\rho\left(  \boldsymbol{r}\right)  =\left[  D_{\ast}^{1}\left(  \rho\right)
\right]  \left(  \boldsymbol{r,r}\right)  =\sum_{1\leq j\leq N}\chi_{j}\left(
\boldsymbol{r}\right)  ^{2} \label{Genexp}%
\end{equation}
in terms of real orthonormal square summable functions $\left\{  \left.
\chi_{j}\right\vert 1\leq j\leq N\right\}  $.

In general one can express the density in terms of orthogonal square
integrable functions as
\begin{equation}
\rho\left(  \boldsymbol{r}\right)  =\sum_{1\leq j\leq\nu}\left\vert
\varphi_{j}\left(  \boldsymbol{r}\right)  \right\vert ^{2}, \label{OrthogOrb}%
\end{equation}
where
\begin{equation}
\left\langle \left.  \varphi_{j}\right\vert \varphi_{k}\right\rangle
=N_{j}\delta_{jk};0<N_{j}\leq N,1\leq j\leq\nu;N_{j}=0,\nu+1\leq j\leq r
\end{equation}
and $\nu\equiv\nu\left(  \boldsymbol{\varphi}\right)  $ is called the rank of
the expansion. If $\nu\geq N$ then the functions $\boldsymbol{\varphi}=\left(
\varphi_{1},\cdots,\varphi_{r}\right)  $ can be interpretted as orbitals. One
can also expand the density $\rho$ in terms of orthonormal functions
$\boldsymbol{\tilde{\varphi}}$ and weights $\boldsymbol{N}=\left(
N_{1},\cdots,N_{r}\right)  $ as
\begin{equation}
\rho\left(  \boldsymbol{r}\right)  =\sum_{1\leq j\leq\nu}N_{j}\left\vert
\tilde{\varphi}_{j}\left(  \boldsymbol{r}\right)  \right\vert ^{2}.
\label{orthonormal}%
\end{equation}
Again if $\nu\geq N$ then $\boldsymbol{\tilde{\varphi}}$ and $\boldsymbol{N}$
can be interpretted as normalized orbitals and occupation numbers
respectively. One can make the relationship expressed eq. $\left(
\ref{OrthogOrb}\right)  $ and eq. $\left(  \ref{orthonormal}\right)  $
explicit through the introduction of two functionals $\Gamma_{1}$ and
$\Gamma_{2}:$

\begin{itemize}
\item The first functional $\Gamma_{1}$ is defined by
\begin{equation}
\Gamma_{1}:\mathfrak{S}_{1}\subset\prod\limits_{1\leq j\leq r}\mathcal{H}%
^{1}\rightarrow L_{1+}\left(  \mathbb{R}^{3}\right)  \subset L_{1}\left(
\mathbb{R}^{3}\right)  ,
\end{equation}
where
\begin{subequations}
\label{GammaAll}%
\begin{align}
\rho &  =\Gamma_{1}\left(  \boldsymbol{\varphi}\right)  ,\\
\left[  \Gamma_{1}\left(  \boldsymbol{\varphi}\right)  \right]  \left(
\boldsymbol{r}\right)   &  =\sum_{1\leq j\leq r}\left\vert \varphi_{j}\left(
\boldsymbol{r}\right)  \right\vert ^{2}=\rho\left(  \boldsymbol{r}\right)
\label{Gamma}%
\end{align}
and%
\end{subequations}
\begin{equation}
\mathfrak{S}_{1}=\left\{  \boldsymbol{\varphi}\left\vert \varphi_{j}%
,\varphi_{k}\in\mathcal{H}^{1},\left\langle \left.  \varphi_{j}\right\vert
\varphi_{k}\right\rangle =N_{j}\delta_{jk},\,\boldsymbol{N}\in\mathfrak{N}%
\left(  N\right)  \right.  \right\}  , \label{set1}%
\end{equation}
where
\begin{equation}
\mathfrak{N}\left(  N\right)  =\left\{  \boldsymbol{N}\left\vert 0\leq
N_{j}\leq N,1\leq j\leq r;\sum_{1\leq j\leq r}N_{j}=N\right.  \right\}  .
\end{equation}
and as many as $r-1$ of the component functions, $\varphi_{j},$ may be
identically zero. The set $\mathfrak{S}_{1}$ includes all possible expansions
of $\rho$ even the one term case considered in eq. $\left(  \ref{rho_f}%
\right)  $.

\item The second functional $\Gamma_{2}$ is defined by
\begin{equation}
\Gamma_{2}:\mathfrak{S}_{2}\subset\prod\limits_{1\leq j\leq r}\mathcal{H}%
^{1}\times\mathbb{R}^{r}\rightarrow L_{1+}\left(  \mathbb{R}^{3}\right)
\subset L_{1}\left(  \mathbb{R}^{3}\right)  ,
\end{equation}
where
\begin{equation}
\mathfrak{S}_{2}=\left\{  \left(  \boldsymbol{\tilde{\varphi},N}\right)
\left\vert \left\vert \tilde{\varphi}_{j}\right\rangle ,\left\vert
\tilde{\varphi}_{k}\right\rangle \in\mathcal{H}^{1},\left\langle \left.
\varphi_{j}\right\vert \varphi_{k}\right\rangle =\delta_{jk},\,1\leq j\leq
r;\quad\boldsymbol{N}\in\mathfrak{N}\left(  N\right)  \right.  \right\}  .
\label{set2}%
\end{equation}


Clearly it is always possible to find $\boldsymbol{\varphi}$ and $\left(
\boldsymbol{\tilde{\varphi},N}\right)  $ such that
\begin{equation}
\rho=\Gamma_{1}\left(  \boldsymbol{\varphi}\right)  =\Gamma_{2}\left(
\boldsymbol{\tilde{\varphi},N}\right)
\end{equation}
and if the number of non zero components in $\boldsymbol{\varphi}$ and
$\boldsymbol{N}$ are greater than or equal to $N$ then $\boldsymbol{\varphi}$
and $\left\{  \boldsymbol{\tilde{\varphi},N}\right\}  $ determine the same
FORDO $D^{1}$ i.e.
\begin{equation}
D^{1}\left(  \boldsymbol{\varphi}\right)  =D^{1}\left(  \boldsymbol{\tilde
{\varphi},N}\right)  =\sum_{1\leq j\leq r}\left\vert \varphi_{j}\right\rangle
\left\langle \varphi_{j}\right\vert =\sum_{1\leq j\leq r}\left\vert
\tilde{\varphi}_{j}\right\rangle \left\langle \tilde{\varphi}_{j}\right\vert
N_{j}.
\end{equation}
The optimization wrt $\rho$ can thus be expressed in terms of
$\boldsymbol{\varphi}$ via $\Gamma_{1}$ or $\left(  \boldsymbol{\tilde
{\varphi},N}\right)  $ via $\Gamma_{2}$, However in order to use sensitivity
theory to relate the variation of the energy with respect to occupation
numbers to the eigenvalues of the Kohn-Sham operator we focus on $\Gamma_{1}.$
\end{itemize}

The functional $\Gamma_{1}$ is not one to one and is constant on the
equivalence class $\mathfrak{S}\left(  \rho\right)  $ of functions%
\begin{equation}
\mathfrak{S}\left(  \rho\right)  =\bigcup\limits_{1\leq\nu\leq r}%
\mathfrak{S}\left(  \nu,\rho\right)  ,
\end{equation}
where the equivalence classes of constant rank, $\mathfrak{S}\left(  \nu
,\rho\right)  ,$ are defined by
\begin{equation}
\mathfrak{S}\left(  \nu,\rho\right)  =\left\{  \left.  \boldsymbol{\varphi
}\right\vert \nu=\nu\left(  \boldsymbol{\varphi}\right)  ,\Gamma_{1}\left(
\boldsymbol{\varphi}\right)  =\rho\right\}  .
\end{equation}
In the following we will limit consideration to the case where $\nu\left(
\boldsymbol{\varphi}\right)  \geq N$ i.e. to orbitals. Each preimage set
$\Gamma_{1}^{-1}\left(  \rho\right)  =\mathfrak{S}\left(  \rho\right)
\subset\mathfrak{S}_{1}$ contains, in general, both single sets of orthogonal
orbitals normalized to different $\boldsymbol{N}\in\mathfrak{N}\left(
N\right)  $ and different sets of orthogonal functions normalized to the same
$\boldsymbol{N}\in\mathfrak{N}\left(  N\right)  ,$ which corresponds to the
fact that a given density can be expanded in many ways in terms of occupation
numbers and orbitals.

\subsection{Orbital functionals\label{GSOFunc}}

One can express the total energy via a universal orbital functional,
$\mathcal{F}_{O}$, whose value is defined by the HK energy functional, as
\begin{equation}
\mathcal{F}_{O}\left(  \boldsymbol{\varphi}\right)  =\mathcal{F}_{HK}\left(
\Gamma_{1}\left(  \boldsymbol{\varphi}\right)  \right)  =\mathcal{F}%
_{HK}\left(  \rho\right)  . \label{FO}%
\end{equation}
The orbital energy functional $\mathcal{F}_{O}$ can be decomposed as
\begin{equation}
\mathcal{F}_{O}\left(  \boldsymbol{\varphi}\right)  =\mathcal{F}_{T}\left(
\boldsymbol{\varphi}\right)  +\mathcal{F}_{C}\left(  \boldsymbol{\varphi
}\right)  +\mathcal{F}_{OXC}\left(  \boldsymbol{\varphi}\right)  , \label{KS1}%
\end{equation}
where the orbital kinetic energy functional $\mathcal{F}_{T}$ is given by
\begin{equation}
\mathcal{F}_{T}\left(  \boldsymbol{\varphi}\right)  =\sum_{1\leq j\leq
r}\left\langle \varphi_{j}\left\vert K\varphi_{j}\right.  \right\rangle ,
\end{equation}
the orbital exchange-corellation functional $\mathcal{F}_{OXC}$ (which depends
on $\mathcal{F}_{T}\boldsymbol{)}$, is defined as
\begin{equation}
\mathcal{F}_{OXC}\left(  \boldsymbol{\varphi}\right)  =\mathcal{F}_{XC}\left(
\Gamma_{1}\left(  \boldsymbol{\varphi}\right)  \right)  +\mathcal{F}%
_{K}\left(  \Gamma_{1}\left(  \boldsymbol{\varphi}\right)  \right)
-\mathcal{F}_{T}\left(  \boldsymbol{\varphi}\right)  , \label{KS2}%
\end{equation}
and one recalls the definitions HK exchange-corellation functional
\begin{equation}
\mathcal{F}_{XC}\left(  \Gamma_{1}\left(  \boldsymbol{\varphi}\right)
\right)  =\operatorname{Tr}\left\{  V_{2}D_{\ast}\left(  \Gamma_{1}\left(
\boldsymbol{\varphi}\right)  \right)  \right\}  -\mathcal{F}_{C}\left(
\Gamma_{1}\left(  \boldsymbol{\varphi}\right)  \right)  \label{KS3}%
\end{equation}
and the HK kinetic functional
\begin{equation}
\mathcal{F}_{K}\left(  \Gamma_{1}\left(  \boldsymbol{\varphi}\right)  \right)
=\operatorname{Tr}\left\{  KD_{\ast}\left(  \Gamma_{1}\left(
\boldsymbol{\varphi}\right)  \right)  \right\}  \label{KS4}%
\end{equation}
that only depend on the density $\rho$.

One should note that:

\begin{itemize}
\item for fixed values of $\rho$ the HK functionals $\mathcal{F}_{XC}\left(
\Gamma_{1}\left(  \boldsymbol{\varphi}\right)  \right)  $ and $\mathcal{F}%
_{K}\left(  \Gamma_{1}\left(  \boldsymbol{\varphi}\right)  \right)  $ in eqs.
$\left(  \ref{KS3}\right)  $ and $\left(  \ref{KS4}\right)  $ are
\emph{constant} for all $\boldsymbol{\varphi}$ that satisfy $\rho=\Gamma
_{1}\left(  \boldsymbol{\varphi}\right)  $, while the orbital functionals
$\mathcal{F}_{T}\left(  \boldsymbol{\varphi}\right)  $ and $\mathcal{F}%
_{OXC}\left(  \boldsymbol{\varphi}\right)  $ can have \emph{different} values
for values of $\boldsymbol{\varphi}$ that correspond to a fixed values of
$\rho$

\item the total energy functionals $\mathcal{F}_{O}\left(  \boldsymbol{\varphi
}\right)  ,\mathcal{\ }$and $\mathcal{F}_{HK}\left(  \rho\right)  $ have the
same value when $\rho=$ $\Gamma_{1}\left(  \boldsymbol{\varphi}\right)
\equiv\rho\left(  \boldsymbol{\varphi}\right)  $

\item the composition in eq. $\left(  \ref{FO}\right)  $ defines a class of
orbital functionals that have a particular form. One could define orbital
functionals without reference to densities (as done in section \ref{AGPSect})
and their form and properties would be quite different.
\end{itemize}

One can recast the minimization problem that determines the ground state
energy and density of a system characterized by a local potential $\mathrm{v}$
in terms of orbitals as
\begin{equation}
E_{GS}\left(  \mathrm{v}\right)  =\mathcal{E}\left(  \Gamma_{1}\left(
\boldsymbol{\varphi}_{GS}\right)  \boldsymbol{,}\mathrm{v}\right)
=\min_{\boldsymbol{\varphi\in}\mathfrak{S}_{1}}\mathcal{E}_{O}\left(
\boldsymbol{\varphi,}\mathrm{v}\right)  , \label{KSmin}%
\end{equation}
where
\begin{equation}
\mathcal{E}_{O}\left(  \boldsymbol{\varphi,}\mathrm{v}\right)  =\mathcal{F}%
_{O}\left(  \boldsymbol{\varphi}\right)  +f_{\Gamma_{1}\left(
\boldsymbol{\varphi}\right)  }\left(  \mathrm{v}\right)
\end{equation}
and the feasible set $\mathfrak{S}_{1}$ is defined in eq. $\left(
\ref{set1}\right)  $.

The orbital minimization problem eq. $\left(  \ref{KSmin}\right)  $ can be
decomposed into two sequential optimizations.

\subsubsection{Varying orbitals with fixed normalization}

\ One first varies the orbitals keeping their normalization fixed and
determines the saddle point $\mathcal{L}_{\boldsymbol{\varphi}}\left(
\boldsymbol{\varphi}\left(  \boldsymbol{N}\right)  ,\boldsymbol{\varepsilon
}\left(  \boldsymbol{N}\right)  ,\mu\left(  N\right)  ,\mathrm{v}\right)  $ of
the Lagrangian
\begin{equation}
\mathcal{L}_{\boldsymbol{\varphi}}\left(  \boldsymbol{\varphi}%
,\boldsymbol{\varepsilon,N,}\mu\boldsymbol{,}\mathrm{v}\right)  =\mathcal{E}%
_{O}\left(  \boldsymbol{\varphi,}\mathrm{v}\right)  -\sum_{1\leq j,k\leq
r}\varepsilon_{jk}\left\{  \overset{\,}{\left\langle \left.  \varphi
_{j}\right\vert \varphi_{k}\right\rangle -N_{j}\delta_{jk}}\right\}
+\mu\left\{  \sum_{\,1\leq j\leq r}\left\langle \left.  \varphi_{j}\right\vert
\varphi_{j}\right\rangle -N\right\}  \label{OrbLag}%
\end{equation}
i.e.
\begin{equation}
\mathcal{L}_{\boldsymbol{\varphi}}\left(  \boldsymbol{\varphi}\left(
\boldsymbol{N}\right)  ,\boldsymbol{\varepsilon}\left(  \boldsymbol{N}\right)
,\mu\left(  N\right)  ,\mathrm{v}\right)  =\max_{\boldsymbol{\varepsilon,}\mu
}\min_{\boldsymbol{\varphi}}\mathcal{L}_{\boldsymbol{\varphi}}\left(
\boldsymbol{\varphi},\boldsymbol{\varepsilon},\mu\boldsymbol{,N,}%
\mathrm{v}\right)  , \label{KSM1}%
\end{equation}
which leads to optimal values $\mathcal{E}_{O}\left(  \boldsymbol{\varphi
}\left(  \boldsymbol{N}\right)  \boldsymbol{,}\mathrm{v}\right)  =$
$\mathcal{L}_{\boldsymbol{\varphi}}\left(  \boldsymbol{\varphi}\left(
\boldsymbol{N}\right)  ,\boldsymbol{\varepsilon}\left(  \boldsymbol{N}\right)
,\mu\left(  N\right)  \boldsymbol{,}\mathrm{v}\right)  $ of the orbital energy
functional for fixed values of the control parameters $\boldsymbol{N}$. The
particle number constraint $\sum_{\,1\leq j\leq r}\left\langle \left.
\varphi_{j}\right\vert \varphi_{j}\right\rangle =N$ is equivalent to the
normalization constraint imposed on the density so that $\rho\in\mathcal{D}$.

The saddle points of $\mathcal{L}_{\boldsymbol{\varphi}}\left(
\boldsymbol{\varphi},\boldsymbol{\varepsilon,N,}\mu\boldsymbol{,}%
\mathrm{v}\right)  $ are characterized by the necessary conditions
\begin{subequations}
\label{KSC}%
\begin{align}
\frac{\delta\mathcal{L}_{\boldsymbol{\varphi}}\left(  \boldsymbol{\varphi
},\boldsymbol{\varepsilon,}\mu,\boldsymbol{N,}\mathrm{v}\right)  }%
{\delta\varphi_{j}^{\ast}}  &  =0;\,1\leq j\leq r,\label{KSC1}\\
\frac{\partial\mathcal{L}_{\boldsymbol{\varphi}}\left(  \boldsymbol{\varphi
},\boldsymbol{\varepsilon,}\mu,\boldsymbol{N,}\mathrm{v}\right)  }%
{\partial\varepsilon_{jk}}  &  =0;\,1\leq j,k\leq r\label{KSC2}\\
\frac{\partial\mathcal{L}_{\boldsymbol{\varphi}}\left(  \boldsymbol{\varphi
},\boldsymbol{\varepsilon,}\mu,\boldsymbol{N,}\mathrm{v}\right)  }{\partial
\mu}  &  =0 \label{KSC3}%
\end{align}
The condition eq. $\left(  \ref{KSC1}\right)  $ implies that at an optimal
point the matrix of Lagrange multipliers $\boldsymbol{\varepsilon}\left(
\boldsymbol{N}\right)  $ is hermitian i.e.
\end{subequations}
\begin{equation}
\varepsilon\left(  \boldsymbol{N}\right)  _{jk}=\varepsilon\left(
\boldsymbol{N}\right)  _{kj}^{\ast},\quad1\leq j,k\leq r
\end{equation}
and
\begin{equation}
\frac{\delta\mathcal{E}_{O}\left(  \boldsymbol{\varphi}\left(  \boldsymbol{N}%
\right)  ,\mathrm{v}\right)  }{\delta\varphi_{j}^{\ast}}-\sum_{1\leq k\leq
r}\left(  \varepsilon_{jk}\left(  \boldsymbol{N}\right)  -\mu\left(  N\right)
\delta_{jk}\right)  \left\vert \varphi_{k}\left(  \boldsymbol{N}\right)
\right\rangle =0;\,1\leq j\leq r. \label{LagPar}%
\end{equation}
Noting that as $\rho=\Gamma_{1}\left(  \boldsymbol{\varphi}\right)  $ and
$\mathcal{E}_{O}\left(  \boldsymbol{\varphi},\mathrm{v}\right)  \equiv$
$\mathcal{E}_{O}\left(  \rho\left(  \boldsymbol{\varphi}\right)
,\mathrm{v}\right)  $ one can obtain
\begin{align}
\frac{\delta\mathcal{E}_{O}\left(  \rho\left(  \boldsymbol{\varphi}\right)
\boldsymbol{,}\mathrm{v}\right)  }{\delta\varphi_{j}^{\ast}\left(
\boldsymbol{r}\right)  }  &  =\int\frac{\delta\mathcal{E}_{O}\left(
\rho\left(  \boldsymbol{\varphi}\right)  \boldsymbol{,}\mathrm{v}\right)
}{\delta\rho\left(  \boldsymbol{r}^{\prime}\right)  }\frac{\delta\rho\left(
\boldsymbol{r}^{\prime}\right)  }{\delta\varphi_{j}^{\ast}\left(
\boldsymbol{r}\right)  }d\boldsymbol{r}^{\prime}\nonumber\\
&  =\int\frac{\delta\mathcal{E}_{O}\left(  \rho\left(  \boldsymbol{\varphi
}\right)  \boldsymbol{,}\mathrm{v}\right)  }{\delta\rho\left(  \boldsymbol{r}%
\right)  }\varphi_{j}\left(  \boldsymbol{r}^{\prime}\right)  \delta\left(
\boldsymbol{r-r}^{\prime}\right)  d\boldsymbol{r}^{\prime}=\frac
{\delta\mathcal{E}_{O}\left(  \rho\left(  \boldsymbol{\varphi}\right)
\boldsymbol{,}\mathrm{v}\right)  }{\delta\rho\left(  \boldsymbol{r}\right)
}\varphi_{j}\left(  \boldsymbol{r}\right)
\end{align}
which gives
\begin{equation}
\frac{\delta\mathcal{E}_{O}\left(  \rho\left(  \boldsymbol{\varphi}\right)
\boldsymbol{,}\mathrm{v}\right)  }{\delta\rho}\left\vert \varphi
_{j}\right\rangle =\left\vert \frac{\delta\mathcal{E}_{O}\left(  \rho\left(
\boldsymbol{\varphi}\right)  \boldsymbol{,}\mathrm{v}\right)  }{\delta\rho
}\varphi_{j}\right\rangle ,
\end{equation}
where the operator $\frac{\delta\mathcal{E}_{O}\left(  \rho\left(
\boldsymbol{\varphi}\right)  \boldsymbol{,}\mathrm{v}\right)  }{\delta\rho}$
is a multiplication operator in the coordinate representation defined by
\begin{equation}
\left[  \frac{\delta\mathcal{E}_{O}\left(  \rho\left(  \boldsymbol{\varphi
}\right)  \boldsymbol{,}\mathrm{v}\right)  }{\delta\rho}\varphi_{j}\right]
\left(  \boldsymbol{r}\right)  =\left(  \frac{\delta\mathcal{E}_{O}\left(
\rho\left(  \boldsymbol{\varphi}\right)  \boldsymbol{,}\mathrm{v}\right)
}{\delta\rho\left(  \boldsymbol{r}\right)  }\right)  \left(  \varphi
_{j}\left(  \boldsymbol{r}\right)  \right)  .
\end{equation}
Substituting this identity back into eq. $\left(  \ref{LagPar}\right)  $ one
obtains for the special case of optimized orbitals, $\boldsymbol{\varphi
\left(  \boldsymbol{N}\right)  },$ that%
\begin{equation}
\frac{\delta\mathcal{E}_{O}\left(  \rho\left(  \boldsymbol{\varphi\left(
\boldsymbol{N}\right)  }\right)  ,\mathrm{v}\right)  }{\delta\rho}\left\vert
\varphi_{j}\left(  \boldsymbol{N}\right)  \right\rangle =\sum_{1\leq k\leq
r}\left(  \varepsilon_{jk}\left(  \boldsymbol{N}\right)  -\mu\left(  N\right)
\delta_{jk}\right)  \left\vert \varphi_{k}\left(  \boldsymbol{N}\right)
\right\rangle ,\,\varepsilon_{jk}\left(  \boldsymbol{N}\right)  =\varepsilon
_{kj}^{\ast}\left(  \boldsymbol{N}\right)  , \label{FockEq}%
\end{equation}
thus obtaining orbital equations corresponding to the condition eq. $\left(
\ref{KSC1}\right)  .$

The eqs $\left(  \ref{KSC2}\right)  $ and $\left(  \ref{KSC3}\right)  $
restate the constraint conditions and as $\left\Vert \left\vert \varphi
_{k}\left(  \boldsymbol{N}\right)  \right\rangle \right\Vert =N_{k}%
=0,\nu+1\leq k\leq r$ eqs $\left(  \ref{OrbLag},\ref{KSM1},\ref{KSC}%
,\ref{LagPar},\ref{FockEq}\right)  $ \emph{only} refer to the variables
$\left(  \varphi_{1},\cdots,\varphi_{\nu}\right)  $ i.e. to the occupied
orbitals. One should note that eq. $\left(  \ref{FockEq}\right)  $ is
\emph{not }an eigenvalue equation unless $\mathcal{E}_{O}\left(
\boldsymbol{\varphi}\left(  \boldsymbol{N}\right)  ,\mathrm{v}\right)  $ is
invariant to unitary mixing of $\boldsymbol{\varphi}\left(  \boldsymbol{N}%
\right)  $ that produce the same density. However to this end one can invoke
Harrimans theorem \cite{Harriman81} which shows that there is always a class
of IPS's, $\mathfrak{S}\left(  N,\rho\right)  $ (the Harriman set for
densities associated to states that contain on the average $N$ electrons) that
produce a given density, so without loss of generality we can always restrict
the domain of $\Gamma_{1}$ to $\mathfrak{S}\left(  N,\rho\right)  $ and
replace eq. $\left(  \ref{FockEq}\right)  $ by
\begin{equation}
\frac{\delta\mathcal{E}_{O}\left(  \rho\left(  \boldsymbol{\varphi\left(
\boldsymbol{N}\right)  }\right)  \boldsymbol{,}\mathrm{v}\right)  }{\delta
\rho}\left\vert \varphi_{j}\left(  \boldsymbol{N}\right)  \right\rangle
=\left(  \varepsilon_{jj}\left(  \boldsymbol{N}\right)  -\mu\left(  N\right)
\right)  \left\vert \varphi_{j}\left(  \boldsymbol{N}\right)  \right\rangle
;\quad1\leq j\leq N \label{FockHarri}%
\end{equation}
as $\mathcal{E}_{O}\left(  \boldsymbol{\varphi}\left(  \boldsymbol{\iota}%
_{N}\right)  ,\mathrm{v}\right)  $, where
\begin{equation}
\boldsymbol{\iota}_{N}=\left(  \overset{N}{\overbrace{1,\cdots,1}%
},\overset{r-N}{\overbrace{0,\cdots,0}}\right)  ,
\end{equation}
is inavariant to unitary mixing of $\boldsymbol{\varphi}\left(
\boldsymbol{\iota}_{N}\right)  \in\mathfrak{S}\left(  N,\rho\right)  .$

From sensitivity theory the rate of change of the orbital energy functional
with respect to changes in occupation numbers, which in this case are the
control parameters $\boldsymbol{N}$, are
\begin{equation}
\frac{\partial\mathcal{E}_{O}\left(  \boldsymbol{\varphi\left(  \boldsymbol{N}%
\right)  },\mathrm{v}\right)  }{\partial N_{j}}=\varepsilon_{jj}\left(
\boldsymbol{N}\right)  -\mu\left(  N\right)  ;1\leq j\leq\nu, \label{Janak}%
\end{equation}
which is Janaks Theorem for \emph{any }fixed set $\boldsymbol{\boldsymbol{N}}$
of occupation numbers and optimized orbitals $\boldsymbol{\varphi\left(
\boldsymbol{N}\right)  }$ normalized to these values. If we take expansion eq.
$\left(  \ref{FockEq}\right)  $ into account we can obtain these relationships
as expectation values of the operator $\frac{\delta\mathcal{E}_{O}\left(
\rho\left(  \boldsymbol{\varphi\left(  \boldsymbol{N}\right)  }\right)
\boldsymbol{,}\mathrm{v}\right)  }{\delta\rho}-\mu I,$ i.e.
\begin{equation}
\frac{\partial\mathcal{E}_{O}\left(  \rho\left(  \boldsymbol{\varphi\left(
\boldsymbol{N}\right)  }\right)  ,\mathrm{v}\right)  }{\partial N_{j}%
}=\left\langle \tilde{\varphi}_{j}\left(  \boldsymbol{N}\right)  \left\vert
\left\{  \frac{\delta\mathcal{E}_{O}\left(  \rho\left(  \boldsymbol{\varphi
\left(  \boldsymbol{N}\right)  }\right)  \boldsymbol{,}\mathrm{v}\right)
}{\delta\rho}-\mu\left(  N\right)  I\right\}  \right.  \tilde{\varphi}%
_{j}\left(  \boldsymbol{N}\right)  \right\rangle =\varepsilon_{jj}\left(
\boldsymbol{N}\right)  -\mu\left(  N\right)  ,
\end{equation}
where
\begin{equation}
\left\vert \tilde{\varphi}_{j}\left(  \boldsymbol{N}\right)  \right\rangle
=\frac{\left\vert \varphi_{j}\left(  \boldsymbol{N}\right)  \right\rangle
}{\left(  \left\langle \left.  \varphi_{j}\left(  \boldsymbol{N}\right)
\right\vert \varphi_{j}\left(  \boldsymbol{N}\right)  \right\rangle \right)
^{\frac{1}{2}}}.
\end{equation}
By considering the sensitivity of the energy to changes of the average number
of electrons in the system one obtains
\begin{equation}
\frac{\partial\mathcal{E}_{O}\left(  \boldsymbol{\varphi}\left(
\boldsymbol{N}\right)  \boldsymbol{,}\mathrm{v}\right)  }{\partial N}%
=\mu\left(  N\right)
\end{equation}
and $\mu\left(  N\right)  $ can be interpreted as the chemical potential of
the interacting molecular system.

\subsubsection{Varying the normalization}

In the second optimization we can restrict attention to control parameters,
$\boldsymbol{N},$ such that $\nu\left(  \boldsymbol{N}\right)  \geq N$ and
$0\leq N_{j}\leq1,1\leq j\leq r$, i.e. that can be interpretted as occupation
numbers. There one finds minima of $\mathcal{E}_{O}\left(  \boldsymbol{\varphi
\left(  \boldsymbol{\theta}\right)  ,}\mathrm{v}\right)  \equiv\mathcal{E}%
_{O}\left(  \boldsymbol{\varphi\left(  N\left(  \boldsymbol{\theta}\right)
\right)  ,}\mathrm{v}\right)  $ with respect to a constrained variation of the
normalization control parameters $\boldsymbol{N}=N\left(  \boldsymbol{\theta
}\right)  $ parameterized by $\boldsymbol{\theta}$, where
\[
0\leq\theta_{j}\leq\frac{\pi}{2};\quad1\leq j\leq r
\]
for a fixed $N$ that are constrained by $\sum_{1\leq j\leq r}N_{j}=\sum_{1\leq
j\leq r}\sin\theta_{j}=N$. Stationary points $\boldsymbol{\theta}_{\#}\left(
N\right)  $ are determined by
\begin{equation}
\mathcal{L}_{\boldsymbol{\theta}}\left(  \boldsymbol{\varphi\left(
\boldsymbol{\theta}_{\#}\right)  ,}N\boldsymbol{,}\eta\left(  N\right)
\boldsymbol{,}\mathrm{v}\right)  =\max_{\boldsymbol{\eta}}\min
_{\boldsymbol{\theta}}\mathcal{L}_{\boldsymbol{\theta}}\left(
\boldsymbol{\varphi,}N\boldsymbol{,}\eta\boldsymbol{,}\mathrm{v}\right)
\end{equation}
the Lagrangian
\begin{equation}
\mathcal{L}_{\boldsymbol{\theta}}\left(  \boldsymbol{\varphi,}N\boldsymbol{,}%
\eta\boldsymbol{,}\mathrm{v}\right)  =\mathcal{E}_{O}\left(
\boldsymbol{\varphi}\left(  \boldsymbol{\theta}\right)  \boldsymbol{,}%
\mathrm{v}\right)  +\eta\left\{  \sum_{1\leq j\leq r}\sin\theta_{j}-N\right\}
\end{equation}
where
\begin{align}
\frac{\partial\mathcal{L}_{\boldsymbol{\theta}}\left(  \boldsymbol{\varphi
,}N\boldsymbol{,}\eta\boldsymbol{,}\mathrm{v}\right)  }{\partial\theta_{j}}
&  =0;1\leq j\leq r\\
\frac{\partial\mathcal{L}_{\boldsymbol{\theta}}\left(  \boldsymbol{\varphi
,}N\boldsymbol{,}\eta\boldsymbol{,}\mathrm{v}\right)  }{\partial\eta}  &
=0.\nonumber
\end{align}
The derivatives with respect to $\boldsymbol{\theta}$ are given by the chain
rule
\begin{equation}
\frac{\partial\mathcal{E}_{O}\left(  \boldsymbol{\varphi}\left(
\boldsymbol{\theta}\left(  N\right)  \right)  \boldsymbol{,}\mathrm{v}\right)
}{\partial\theta_{j}}=\sum\limits_{1\leq k\leq r}\frac{\partial\mathcal{E}%
_{O}\left(  \boldsymbol{\varphi}\left(  \boldsymbol{N}\right)  \boldsymbol{,}%
\mathrm{v}\right)  }{\partial N_{k}}\frac{\partial N_{k}}{\partial\theta_{j}%
}=\frac{\partial\mathcal{E}_{O}\left(  \boldsymbol{\varphi}\left(
\boldsymbol{N}\right)  \boldsymbol{,}\mathrm{v}\right)  }{\partial N_{j}}%
\cos\theta_{j};\quad1\leq j\leq r \label{ThetaDer}%
\end{equation}
and
\begin{equation}
\frac{\partial\mathcal{E}_{O}\left(  \boldsymbol{\varphi}\left(
\boldsymbol{N}\right)  \boldsymbol{,}\mathrm{v}\right)  }{\partial N_{j}%
}=\varepsilon_{jj}\left(  \boldsymbol{N}\left(  \boldsymbol{\theta}\right)
\right)  -\mu\left(  N\right)  .
\end{equation}
Giving
\begin{equation}
\frac{\partial\mathcal{E}_{O}\left(  \boldsymbol{\varphi}\left(
\boldsymbol{\theta}_{\#}\right)  \boldsymbol{,}\mathrm{v}\right)  }%
{\partial\theta_{j}}=\left\{  \varepsilon_{jj}\left(  \boldsymbol{N}\left(
\boldsymbol{\theta}_{\#}\right)  \right)  -\mu\left(  N\right)  \right\}
\cos\theta_{\#j}=0. \label{theta}%
\end{equation}
Thus
\begin{equation}
\varepsilon_{jj}\left(  \boldsymbol{N}\left(  \boldsymbol{\theta}_{\#}\right)
\right)  -\mu\left(  N\right)  =0\text{ and/or }\left\{  \cos\theta
_{\#j}=0\Leftrightarrow\theta_{\#j}=\frac{\pi}{2}\right\}  ;\quad1\leq j\leq
r. \label{Con1}%
\end{equation}
If $\ \theta_{j}=\frac{\pi}{2}$ ($N_{j}\left(  \boldsymbol{\theta}%
_{\#}\right)  =1)$ then $\varepsilon_{jj}$ can have any value (even
$\mu\left(  N\right)  $), while if $\theta_{\#j}\neq\frac{\pi}{2}$ (which
corresponds to non integer occupation) then $\varepsilon_{jj}=$ $\mu\left(
N\right)  $ i.e. the Lagrange parameters $\left\{  \varepsilon_{jj}\right\}  $
are degenerate. If $\varepsilon_{jj}=\mu\left(  N\right)  $ then $\theta
_{\#j}$ does not have to have the value $\frac{\pi}{2}$ though it can. Since
the inception of the KS approach to DFT\ there has been an extensive
discussion of the physical meaning of the eigenvalues of the KS operator and
their relationship to the Fermi energy and chemical potential.

The ground state energy $E\left(  \rho_{GS}\boldsymbol{,}\mathrm{v}\right)
=\mathcal{E}_{O}\left(  \boldsymbol{\varphi}\left(  \boldsymbol{\theta}%
_{\#}\right)  ,\mathrm{v}\right)  $ and the ground state density can then be
expanded as
\begin{equation}
\rho_{GS}\left(  \boldsymbol{r}\right)  =\sum_{1\leq j\leq r}\left\vert
\left[  \varphi_{j}\left(  \boldsymbol{\theta}_{\#}\right)  \right]  \left(
\boldsymbol{r}\right)  \right\vert ^{2}. \label{GSen}%
\end{equation}
The expansion eq. $\left(  \ref{GSen}\right)  $ is very \emph{non} unique in
that any set $\boldsymbol{\varphi}\left(  \boldsymbol{\theta}_{\#}^{\prime
}\right)  \in\mathfrak{S}\left(  \rho_{GS}\right)  $, where
$\boldsymbol{\theta}_{\#}^{\prime}$ satisfies eq. $\left(  \ref{theta}\right)
$ produces $\rho_{GS}$ and the same value of the energy functional i.e.
$\mathcal{E}_{O}\left(  \boldsymbol{\varphi}\left(  \boldsymbol{\theta}%
_{\#}^{\prime}\right)  ,\mathrm{v}\right)  =\mathcal{E}_{O}\left(
\boldsymbol{\varphi}\left(  \boldsymbol{\theta}_{\#}\right)  ,\mathrm{v}%
\right)  $.

\subsection{The Kohn-Sham Density Functional\label{KSDF}}

As pointed out in the last subsection the minima of $\mathcal{E}_{O}\left(
\boldsymbol{\varphi}\left(  \boldsymbol{\theta}_{\#}\right)  ,\mathrm{v}%
\right)  ,$
\begin{equation}
\mathcal{E}_{O}\left(  \boldsymbol{\varphi}\left(  \boldsymbol{\theta}%
_{\#}\right)  \boldsymbol{,}\mathrm{v}\right)  =\min_{\boldsymbol{\varphi\in
}\mathfrak{S}_{1}}\mathcal{E}_{O}\left(  \boldsymbol{\varphi,}\mathrm{v}%
\right)
\end{equation}
are highly degenerate e.g.
\begin{equation}
\mathcal{E}_{O}\left(  \boldsymbol{\varphi}\left(  \boldsymbol{\theta}%
_{\#}\right)  \boldsymbol{,}\mathrm{v}\right)  =const.\,\forall
\boldsymbol{\varphi\in}\mathfrak{S}\left(  \rho_{GS}\right)  .
\end{equation}
This degeneracy is due to the singularity of the map $\Gamma_{1}$ which can be
removed by constraining $\boldsymbol{\varphi}$ to be a specific representative
of the equivalence classes $\mathfrak{S}\left(  \rho\right)  .$ One way to
obtain such a representative is to consider the orbitals formed by the NOsand
occupation numbers of the FORDO of the $N$-electron ground state $D_{0}\left(
\rho\right)  =A_{0}\left(  \rho\right)  ^{2}$ of a model independent particle
$N$-electron system$,$ that has a fixed density $\rho.$ This corresponds to
finding the minimizers of the universal kinetic energy functional $K_{U}$
\begin{equation}
K_{U}\left(  A\right)  =\operatorname{Tr}\left\{  KA^{2}\right\}  ,
\end{equation}
subject to the constraint
\begin{equation}
\Xi_{2}\left(  A\right)  =\left(  A\left\vert \widehat{\Phi^{\dagger}\Phi
}A\right.  \right)  =\rho,\,A\in\mathcal{B}_{2sa}\left(  \mathcal{H}%
^{N}\right)  .
\end{equation}
The constrained minima $K_{U}\left(  D_{0}\left(  \rho\right)  \right)
=\operatorname{Tr}\left\{  KD_{0}\left(  \rho\right)  \right\}
=\operatorname{Tr}\left\{  KA_{0}\left(  \rho\right)  ^{2}\right\}  $ is given
as an unconstrained saddle point $\left(  A_{0}\left(  \rho\right)
,\lambda_{0}\left(  \rho\right)  \right)  ,$
\begin{equation}
\mathcal{L}_{K}\left(  A_{0}\left(  \rho\right)  ,\lambda_{0}\left(
\rho\right)  ,\rho\right)  =\max_{\lambda}\min_{A}\mathcal{L}_{K}\left(
A,\lambda,\rho\right)
\end{equation}
of the Lagrangian
\begin{equation}
\mathcal{L}_{K}\left(  A,\lambda,\rho\right)  =K_{U}\left(  A\right)
-\Omega_{A}\left(  \lambda\right)  +\left\langle \left.  \lambda\right\vert
\rho\right\rangle =\left(  A\left\vert \left\{  \left\{  \widehat{K}%
-\widehat{\Omega\left(  \lambda\right)  }\right\}  +\left\langle \left.
\lambda\right\vert \rho\right\rangle \hat{I}\right\}  \right.  A\right)  .
\end{equation}
If a bound state minimum exists for some local one body potential $\lambda
_{0}\left(  \rho\right)  $ (which is also the Lagrange parameter function) one
has
\begin{equation}
\left\{  \widehat{K}-\widehat{\Omega\left(  \lambda_{0}\left(  \rho\right)
\right)  }\right\}  \left\vert A_{0}\left(  \rho\right)  \right)  =0,
\end{equation}
where $A_{0}\left(  \rho\right)  \in\mathcal{B}_{2sa}\left(  \mathcal{H}%
^{N}\right)  $. This gives a stationary state $D_{0}\left(  \rho\right)
=A_{0}\left(  \rho\right)  ^{2}$ of the independent particle system with
Hamiltonian $\left\{  K-\Omega\left(  \lambda_{0}\left(  \rho\right)  \right)
\right\}  $ i.e.
\begin{equation}
\left[  \left\{  K-\Omega\left(  \lambda_{0}\left(  \rho\right)  \right)
\right\}  ,D_{0}\left(  \rho\right)  \right]  =0,
\end{equation}
which in turn defines a FORDO $D_{0}^{1}\left(  \rho\right)  $ given by (see
Appendix \ref{Denmat})
\begin{equation}
D_{0}^{1}\left(  \rho\right)  =NL_{N}^{1}\left(  D_{0}\left(  \rho\right)
\right)  .
\end{equation}
As $K-\Omega\left(  \lambda_{0}\left(  \rho\right)  \right)  $ is an
one-electron operator the state $D_{0}\left(  \rho\right)  $ must be a mixture
of IPSsor a pure IPS and thus determined by the eigenvalues $\left\{
N_{0j}\left(  \rho\right)  \right\}  $ and eigenvectors $\left\{  \left\vert
\widetilde{\varphi_{0j}\left(  \rho\right)  }\right\rangle \right\}  $ of
$D_{0}^{1}\left(  \rho\right)  $. If $D_{0}\left(  \rho\right)  $ is a pure
IPS\ then the occupation numbers $\left\{  N_{0j}\left(  \rho\right)
\right\}  $ are equal to either $1$ or $0$ (i.e. $N$ ones and $r-N$ zeroes),
while if it is a ensemble state than there will be some fractional values and
the number of zero values will be less than $r-N$. Expanding $D_{0}^{1}\left(
\rho\right)  $ in terms of normalized NOs$\left\{  \left.  \left\vert
\widetilde{\varphi_{0j}}\right\rangle \right\vert 1\leq j\leq r\right\}  $
gives
\begin{equation}
D_{0}^{1}\left(  \rho\right)  =\sum_{1\leq j\leq r}N_{0j}\left(  \rho\right)
\left\vert \widetilde{\varphi_{0j}\left(  \rho\right)  }\right\rangle
\left\langle \widetilde{\varphi_{0j}\left(  \rho\right)  }\right\vert ,
\end{equation}
(where in contrast to eq. $\left(  \ref{Gamma}\right)  $ $r\geq N$ ), which
after renormalizing the NOsto the occupation numbers i.e.
\begin{equation}
\left\vert \varphi_{0j}\left(  \rho\right)  \right\rangle =N_{0j}\left(
\rho\right)  ^{\frac{1}{2}}\left\vert \widetilde{\varphi_{0j}\left(
\rho\right)  }\right\rangle ,1\leq j\leq r
\end{equation}
results in
\begin{equation}
\rho\left(  \boldsymbol{r}\right)  =\sum_{1\leq j\leq r}\left\vert \left[
\varphi_{0j}\left(  \rho\right)  \right]  \left(  \boldsymbol{r}\right)
\right\vert ^{2}. \label{gammaKS}%
\end{equation}
If the constrained minimum of $K_{U}$ is non degenerate then there is a unique
correspondence, up to unitary mixing of degenerate eigenvectors of $D_{0}%
^{1}\left(  \rho\right)  ,$ between $\rho$ and $\left\{  \boldsymbol{\varphi
}_{0}\right\}  ,$ thus modulo this unitary mixing one can define an inverse
functional $\Gamma_{0}^{-1}$ as
\begin{equation}
\boldsymbol{\varphi}_{0}=\Gamma_{0}^{-1}\left(  \rho\right)  .
\end{equation}
If the constrained minimum of $K_{U}$ is degenerate, e.g. corresponding to
FORDOs$\left\{  \left.  D_{0\kappa}^{1}\left(  \rho\right)  \right\vert
1\leq\kappa\leq m\right\}  $ of the stationary states $\left\{  \left.
D_{0\kappa}\left(  \rho\right)  \right\vert 1\leq\kappa\leq m\right\}  $ of
$K-\Omega\left(  \lambda_{0}\left(  \rho\right)  \right)  $ then as
\begin{equation}
\left[  \left\{  K-\Omega\left(  \lambda_{0}\left(  \rho\right)  \right)
\right\}  ,D_{0\kappa}\left(  \rho\right)  \right]  =\left[  \left\{
K-\Omega\left(  \lambda_{0}\left(  \rho\right)  \right)  \right\}
,D_{0\kappa}^{1}\left(  \rho\right)  \right]  =0,
\end{equation}
the space of degenerate states corresponds to the same set of orbitals
$\left\{  \left\vert \boldsymbol{\varphi}_{0}\left(  \rho\right)
\right\rangle \right\}  $ but normalized to different occupation numbers
$\left\{  \boldsymbol{N}_{0}\left(  \rho\right)  \right\}  ,$ where
$\boldsymbol{N}_{0}\left(  \rho\right)  $ are convex combinations of $\left\{
\left.  \boldsymbol{N}_{0\kappa}\left(  \rho\right)  \right\vert 1\leq
\kappa\leq m\right\}  $. Hence if degeneracies occur these are reflected in
multiple optimal values for $\boldsymbol{\boldsymbol{N}}_{0}$ for fixed
$\left\{  \boldsymbol{\varphi}_{0}\left(  \rho\right)  \right\}  $, some of
the $\left\{  \left.  N_{0j}\right\vert 1\leq j\leq r\right\}  $ can have non
integer values even for degenerate pure states$\boldsymbol{.}$

The minimal values of the universal kinetic functional $K_{U}\left(
D_{0}\left(  \rho\right)  \right)  $ produce a kinetic energy functional
$\mathcal{F}_{T}$ defined on the representatives $\left\{  \boldsymbol{\varphi
}_{0}\right\}  $ of the equivalence classes $\mathfrak{S}\left(  \rho\right)
$ whose values are given by
\begin{equation}
\mathcal{F}_{T}\left(  \boldsymbol{\varphi}_{0}\right)  =K_{U}\left(
D_{0}\left(  \Gamma_{1}\left(  \boldsymbol{\varphi}_{0}\right)  \right)
\right)  . \label{KEFunc}%
\end{equation}
The value of $\mathcal{F}_{T}$ is constant on the set of orbitals formed from
the unitary mixing and the convex combinations of degenerate states just
described. This allows one to define the Kohn-Sham functional $\mathcal{F}%
_{KS}\left(  \boldsymbol{\varphi}_{0}\right)  $ as a specific decomposition of
the HK energy functional $\mathcal{F}_{HK}\left(  \rho\right)  $ in terms of
orbitals as
\begin{equation}
\mathcal{F}_{KS}\left(  \boldsymbol{\varphi}_{0}\right)  =\mathcal{F}%
_{T}\left(  \boldsymbol{\varphi}_{0}\right)  +\mathcal{F}_{C}\left(
\boldsymbol{\varphi}_{0}\right)  +\mathcal{F}_{KSXC}\left(
\boldsymbol{\varphi}_{0}\right)  ,
\end{equation}
where
\begin{equation}
\mathcal{F}_{KS}\left(  \boldsymbol{\varphi}_{0}\right)  =\mathcal{F}%
_{HK}\left(  \Gamma_{1}\left(  \boldsymbol{\varphi}_{0}\right)  \right)  .
\end{equation}


The relationships described in subsection \ref{GSOFunc} derived for stationary
points $\boldsymbol{\varphi}\left(  \boldsymbol{N}\right)  $ are valid for any
points contained in the equivalence classes that these points belong to, in
particular $\boldsymbol{\varphi}_{0}\left(  \boldsymbol{\boldsymbol{N}}%
_{0}\right)  $. Eq. $\left(  \ref{FockEq}\right)  $, is now an equation for
determining the KS orbitals $\boldsymbol{\varphi}_{0}\left(  \boldsymbol{N}%
\right)  $ and the operator introduced there can be identified as the KS
operator. If the state $D_{0}\left(  \rho\right)  $ is non degenerate and pure
(therefore has only occupation numbers equal to one or zero) then one can
choose a set of orbitals that diagonalize this operator as the energy
$\mathcal{E}_{O}\left(  \boldsymbol{\varphi}_{0}\left(  \boldsymbol{N}%
_{0}\right)  ,\mathrm{v}\right)  $ is invariant to unitary mixing of occupied
orbitals, which is the standard way of producing a canonical set of occupied orbitals.

The reference $N$-electron IPS, $D_{0}\left(  \boldsymbol{\boldsymbol{N}}%
_{0}\right)  $ formed from the NOs$\boldsymbol{\varphi}_{0}$ has two energies
associated with it. One, $E_{0}\left(  \boldsymbol{\varphi}_{0}\left(
\boldsymbol{N}_{0}\right)  \right)  $, is the ground state energy of the
Hamiltonian%
\begin{equation}
H_{0}=K+\mathrm{v},
\end{equation}
which can be expressed as
\begin{equation}
E_{0}\left(  \boldsymbol{\varphi}_{0}\left(  \boldsymbol{N}_{0}\right)
\right)  =\sum_{1\leq j\leq r}\left\langle \tilde{\varphi}_{0j}\left\vert
\left(  K+\mathrm{v}\right)  \right.  \tilde{\varphi}_{0j}\right\rangle
N_{0j}.
\end{equation}
The other is the energy of the Hamiltonian
\begin{equation}
H=H_{0}+V_{2}%
\end{equation}
wrt the IPS $D_{0}\left(  \boldsymbol{N}_{0}\right)  $ given by
\begin{equation}
E_{IP}\left(  D_{0}\left(  \boldsymbol{N}_{0}\right)  \right)
=\operatorname{Tr}\left\{  HD_{0}\left(  \boldsymbol{N}_{0}\right)  \right\}
.
\end{equation}
It should be noted that this value of the energy functional $E_{IP}$ is
\emph{not }optimized in any way and is \emph{not }related in any simple
fashion to the energies $\mathcal{E}_{O}\left(  \varphi\boldsymbol{\left(
\boldsymbol{N}_{0}\right)  },\mathrm{v}\right)  .$

The preceding relationships are equally for orbitals $\boldsymbol{\varphi}%
_{0}\left(  \boldsymbol{N}_{0}\left(  \boldsymbol{\theta}_{\#}\left(
N\right)  \right)  \right)  $ that produce the ground state density $\rho
_{GS}$ and energy $\mathcal{E}_{O}\left(  \boldsymbol{\varphi}_{0}\left(
\boldsymbol{N}\left(  \boldsymbol{\theta}_{\#}\left(  N\right)  \right)
\right)  ,\mathrm{v}\right)  $ and the classical Janaks Theorem is a special
case of eq. $\left(  \ref{Janak}\right)  $
\begin{equation}
\frac{\partial\mathcal{E}_{O}\left(  \boldsymbol{\varphi}_{0}\left(
\boldsymbol{N}_{0}\left(  \boldsymbol{\theta}_{\#}\left(  N\right)  \right)
\right)  ,\mathrm{v}\right)  }{\partial N_{j}}\equiv\frac{\partial
\mathcal{E}_{KS}\left(  \boldsymbol{\varphi}_{0}\left(  \boldsymbol{N}%
_{0}\left(  \boldsymbol{\theta}_{\#}\left(  N\right)  \right)  \right)
,\mathrm{v}\right)  }{\partial N_{j}}=\varepsilon_{jj}\left(
\boldsymbol{\varphi}_{0}\left(  \boldsymbol{N}_{0}\left(  \boldsymbol{\theta
}_{\#}\left(  N\right)  \right)  \right)  \right)  , \label{ClassJanak}%
\end{equation}
where we have constrained the orbitals and occupation numbers to be the ones
that minimize the universal kinetic energy, $K_{U},$ thus producing a model
IPS. One must, however, note that the Lagrange parameters $\left\{
\varepsilon_{jj}\left(  \boldsymbol{\varphi}_{0}\left(  \boldsymbol{N}%
_{0}\left(  \boldsymbol{\theta}_{\#}\left(  N\right)  \right)  \right)
\right)  \right\}  $ in this equation \emph{are not }associated with the
minimization of the universal kinetic energy.

The KS energy functional is defined in terms of orbitals and occupation
numbers and is implicitly defined by the preceding constrained optimization
procedures, that factor the functional through the density. If one did not
consider this factorization and directly optimized the orbitals and occupation
numbers one would determine the same ground state, but not need to introduce a
model IPS\ and orbitals $\boldsymbol{\varphi}_{0}\left(  \boldsymbol{N}%
_{0}\right)  $. In the next section we compare and contrast the relationships
obtained for these functionals with ones obtained for the explicitly defined
AGP functional. These functionals are defined on the space of correlated
states formed by the non linear manifold of AGP\ states, where the energy has
been expressed as an explicit functional of GSOsand parameters that determine
occupation numbers \cite{OrtizWeiner2006} The minimization of this explicit
functional can then be compared to the minimization of the implicit
KS\ functional restricted to the AGP\ manifold.

\section{AGP Energy Functionals \label{AGP}}

We briefly review some of the results found in \cite{OrtizWeiner2006}
describing AGP states and their associated energy functionals. In
\cite{OrtizWeiner2006} orbital and occupation number optimization equations
were derived and it was displayed how the $N$-electron total energy could be
expressed in terms of ionization potentials, kinetic energies and eigenvalues
of a generalized Fock operator associated with GSO's. In contrast to the KS
approach, the occupation numbers were not constrained to refer to IPS's, but
allowed to vary over the much larger domain of states with doubly degenerate
state occupation numbers. It has been argued while this set does not contain
all possible states it contains a large percentage of states important for
molecular systems \cite{Kramers Deg}. The density matrix energy functional in
DMFT allows unrestricted variation of occupation numbers (and thus all
states), but unlike the AGP functional (but in common with the HK\ and KS
density functionals) it is invariant to the phases of the GSOsi.e. the group
$U\left(  1\right)  .$

\subsection{Canonical expansion of an AGP state}

The $2M$-electron AGP state, $\left\vert g^{M}\right\rangle $, is the $M^{th}$
antisymmetric power of the $2$-electron state described by a normalized
geminal
\begin{equation}
\left\vert g\right\rangle =\sum_{1\leq j\leq s}\left\vert \varphi_{j}%
\varphi_{j+s}\right\rangle =\sum_{1\leq j\leq s}n_{j}^{\frac{1}{2}}\left\vert
\tilde{\varphi}_{j}\tilde{\varphi}_{j+s}\right\rangle , \label{geminal}%
\end{equation}
given by
\begin{equation}
\left\vert g^{M}\right\rangle =\left(  S_{M}\right)  ^{-\frac{1}{2}}%
\sum_{1\leq j_{1}<\cdots<j_{M}\leq s}n_{j_{1}}^{\frac{1}{2}}\cdots n_{j_{s}%
}^{\frac{1}{2}}\left\vert \tilde{\varphi}_{j_{1}}\tilde{\varphi}_{j_{1}%
+s}\cdots\tilde{\varphi}_{j_{M}}\tilde{\varphi}_{j_{M}+s}\right\rangle
=\left(  S_{M}\right)  ^{-\frac{1}{2}}\sum_{1\leq j_{1}<\cdots<j_{M}\leq
s}\left\vert \varphi_{j_{1}}\varphi_{j_{1}+s}\cdots\varphi_{j_{M}}%
\varphi_{j_{M}+s}\right\rangle , \label{agp}%
\end{equation}
wherer $\left\{  \left.  n_{j}^{\frac{1}{2}}\right\vert 1\leq j\leq s\right\}
$ are the positive square roots of the occupation numbers%
\begin{equation}
\Omega=\left\{  n_{j}\left\vert 0\leq n_{j}\leq1;\sum_{1\leq j\leq s}%
n_{j}=1;1\leq j\leq s\right.  \right\}
\end{equation}
of the FORDO of the two electron AGP\ state density operator, $\left\{
\left.  \left\vert \tilde{\varphi}_{j}\right\rangle \right\vert 1\leq
j\leq2s\right\}  $ are the \emph{orthonormal} Canonical General Spin Orbitals
(CGSO's) of $\left\vert g\right\rangle $ and $\left\{  \left.  \left\vert
\varphi_{j}\right\rangle \right\vert 1\leq j\leq2s\right\}  $ are the
\emph{orthogonal} CGSOsnormalized to $\boldsymbol{n}$ i.e.
\begin{equation}
\left\vert \varphi_{j}\right\rangle =n_{j}^{\frac{1}{2}}\left\vert
\tilde{\varphi}_{j}\right\rangle ,1\leq j\leq2s.
\end{equation}
The normalization factor, $\left(  S_{M}\right)  ^{-\frac{1}{2}}$, is given in
terms of the elementary symmetric function
\begin{equation}
S_{M}=\sum_{1\leq j_{1}<\cdots<j_{M}\leq s}n_{j_{1}}\cdots n_{j_{M}}.
\end{equation}
It is evident from eq.\ $\left(  \ref{geminal}\right)  $ that the geminal and
geminal power state are at least invariant to the group $SU\left(  2\right)
,$ that is formed by $2\times2$ \emph{special} unitary transformations of the
paired CGSOs$\left\{  \left.  \left\vert \tilde{\varphi}_{j}\right\rangle
,\left\vert \tilde{\varphi}_{j+s}\right\rangle \right\vert 1\leq j\leq
s\right\}  ,$ while if degeneracies exist amongst the $\left\{  \left.
n_{j}\right\vert 1\leq j\leq s\right\}  $ this invariance group is larger
\cite{AGPCoherent}. The canonical CGSOsare unique up to transformations
belonging this invariance group.

\subsection{AGP Energy Functionals\label{AGPSect}}

The energy of the AGP state can be expressed as a functional $\mathcal{E}%
_{AGP}$ of $\boldsymbol{\varphi}\left(  \boldsymbol{n}\right)  $ by
\begin{align}
\mathcal{E}_{AGP}\left(  \boldsymbol{\varphi}\left(  \boldsymbol{n}\right)
\right)   &  =\frac{\mathcal{H}\left(  \boldsymbol{\varphi}\left(
\boldsymbol{n}\right)  \right)  }{S_{M}\left(  \boldsymbol{n}\right)
}\nonumber\\
&  =\frac{1}{S_{M}\left(  \boldsymbol{n}\right)  }\left\{  \sum_{1\leq j\leq
s}n_{j}S_{M-1}\left(  \jmath\right)  \left\{  \left\langle \tilde{\varphi}%
_{j}\left\vert h_{1}\right.  \tilde{\varphi}_{j}\right\rangle +\left\langle
\tilde{\varphi}_{j+s}\left\vert h\right.  \tilde{\varphi}_{j+s}\right\rangle
+\left\langle \tilde{\varphi}_{j}\tilde{\varphi}_{j+s}||\tilde{\varphi}%
_{j}\tilde{\varphi}_{j+s}\right\rangle \right\}  \right. \nonumber\\
&  +\sum_{1\leq j,k\leq s}n_{j}^{\frac{1}{2}}n_{k}^{\frac{1}{2}}S_{M-1}\left(
\jmath k\right)  \left\langle \tilde{\varphi}_{j}\tilde{\varphi}_{j+s}%
||\tilde{\varphi}_{k}\tilde{\varphi}_{k+s}\right\rangle \left.  +\sum_{1\leq
j<k\leq s}n_{i}n_{j}S_{M-2}\left(  \jmath k\right)  \left\langle
\tilde{\varphi}_{j}\tilde{\varphi}_{k}|||\tilde{\varphi}_{j}\tilde{\varphi
}_{k}\right\rangle \right\}  . \label{energy}%
\end{align}


\subsection{AGP Lagrangians\label{Lagran}}

Using the same philosophy for optimization as discussed for orbital
functionals in the preceding sections one can minimize the energy of the
$2M$-electron state over the manifold of AGP\ states using a nested two step process:

\begin{enumerate}
\item Keeping the normalization $\boldsymbol{n}$ of the
GSOs$\boldsymbol{\varphi}$ fixed one finds the constrained minimum
$\mathcal{E}_{AGP}\left(  \boldsymbol{\varphi}\left(  \boldsymbol{n}\right)
\right)  $ of the functional $\mathcal{E}_{AGP}\left(  \boldsymbol{\varphi
}\right)  $ for the specified values of the control parameters $\boldsymbol{n}%
$ i.e.
\begin{equation}
\mathcal{E}_{AGP}\left(  \boldsymbol{\varphi}\left(  \boldsymbol{n}\right)
\right)  =\min_{\boldsymbol{\varphi\in}\Lambda\left(  \boldsymbol{n}\right)
}\mathcal{E}_{AGP}\left(  \boldsymbol{\varphi}\right)  ,
\end{equation}
where
\[
\Lambda\left(  \boldsymbol{n}\right)  =\left\{  \boldsymbol{\varphi}\left\vert
\left\langle \left.  \varphi_{j}\right\vert \varphi_{k}\right\rangle
=n_{j}\delta_{jk};\sum_{1\leq j\leq s}n_{j}=1\right.  \right\}  .
\]
This is accomplished by finding the unconstrained saddle points $\mathcal{L}%
\left(  \boldsymbol{\varphi\left(  \boldsymbol{n}\right)  ,\varepsilon\left(
\boldsymbol{n}\right)  },\eta\right)  $ of the Lagrangian
\begin{equation}
\mathcal{L}\left(  \boldsymbol{\varphi,\varepsilon}\right)  =\mathcal{E}%
_{AGP}\left(  \boldsymbol{\varphi}\right)  -\sum_{1\leq j,k\leq2s}%
\varepsilon_{kj}\left(  \left\langle \left.  \varphi_{j}\right\vert
\varphi_{k}\right\rangle -n_{j}\delta_{jk}\right)  +\eta\left\{  \sum_{\,1\leq
j\leq2s}\left\langle \left.  \varphi_{j}\right\vert \varphi_{j}\right\rangle
-2\right\}  ,
\end{equation}
i.e.
\begin{equation}
\mathcal{L}\left(  \boldsymbol{\varphi\left(  \boldsymbol{n}\right)
,\varepsilon\left(  \boldsymbol{n}\right)  },\eta\right)  =\max
_{\boldsymbol{\varepsilon},\eta.}\min_{\boldsymbol{\varphi}}\mathcal{L}\left(
\boldsymbol{\varphi,\varepsilon},\eta\right)
\end{equation}
The corresponding stationarity conditions are
\begin{subequations}
\label{AGPstat}%
\begin{align}
\frac{\delta\mathcal{L}\left(  \boldsymbol{\varphi\left(  \boldsymbol{n}%
\right)  ,\varepsilon\left(  \boldsymbol{n}\right)  },\eta\right)  }%
{\delta\boldsymbol{\varphi}}  &  =\boldsymbol{0}\label{agpstat1}\\
\frac{\partial\mathcal{L}\left(  \boldsymbol{\varphi\left(  \boldsymbol{n}%
\right)  ,\varepsilon\left(  \boldsymbol{n}\right)  },\eta\right)  }%
{\partial\boldsymbol{\varepsilon}}  &  =\boldsymbol{0}\label{agpstat2}\\
\frac{\partial\mathcal{L}\left(  \boldsymbol{\varphi\left(  \boldsymbol{n}%
\right)  ,\varepsilon\left(  \boldsymbol{n}\right)  },\eta\right)  }%
{\partial\eta}  &  =\boldsymbol{0} \label{agpstat3}%
\end{align}
and eq. $\left(  \ref{agpstat1}\right)  $ gives
\end{subequations}
\begin{equation}
\frac{\delta\mathcal{E}_{AGP}\left(  \boldsymbol{\varphi}\left(
\boldsymbol{n}\right)  \right)  }{\delta\varphi_{j}}=\sum_{1\leq k\leq
2s}\varepsilon\boldsymbol{\left(  \boldsymbol{n}\right)  }_{jk}\left\vert
\varphi_{k}\boldsymbol{\left(  \boldsymbol{n}\right)  }\right\rangle ,
\end{equation}
which can be expressed \cite{OrtizWeiner2006} in terms of a generalized Fock
operator as
\begin{equation}
F\left(  \boldsymbol{\varphi}\left(  \boldsymbol{n}\right)  \right)
\left\vert \varphi_{j}\left(  \boldsymbol{n}\right)  \right\rangle
=\sum_{1\leq k\leq2s}\varepsilon\boldsymbol{\left(  \boldsymbol{n}\right)
}_{jk}\left\vert \varphi_{k}\boldsymbol{\left(  \boldsymbol{n}\right)
}\right\rangle .
\end{equation}
Sensitivity theory gives%
\[
\frac{\partial\mathcal{E}_{AGP}\left(  \boldsymbol{\varphi}\left(
\boldsymbol{n}\right)  \right)  }{\partial n_{j}}=\varepsilon_{jj}\left(
\boldsymbol{n}\right)  .
\]


\item Varying $\mathcal{E}_{AGP}\left(  \boldsymbol{\varphi}\left(
\boldsymbol{n}\right)  \right)  $ over the normalization control parameters
$\boldsymbol{n}\in\Omega$ to obtain the minimum $\mathcal{E}_{AGP}\left(
\boldsymbol{\varphi}\left(  \boldsymbol{n}_{\#}\right)  \right)  $ i.e.
\begin{equation}
\mathcal{E}_{AGP}\left(  \boldsymbol{\varphi}\left(  \boldsymbol{n}%
_{\#}\right)  \right)  =\min_{\boldsymbol{n}\in\Omega}\mathcal{E}_{AGP}\left(
\boldsymbol{\varphi}\left(  \boldsymbol{n}\right)  \right)  .
\label{agp_control}%
\end{equation}
The optimization eq. $\left(  \ref{agp_control}\right)  $ can be made
unconstrained by expressing $\boldsymbol{n}$ in terms of the variable
$\boldsymbol{\theta}$ as
\begin{equation}
n_{j}\left(  \boldsymbol{\theta}\right)  =\frac{\cos^{2}\theta_{j}}%
{\sum\limits_{1\leq j\leq s}\cos^{2}\theta_{j}}\equiv\frac{\cos^{2}\theta_{j}%
}{\mathcal{N}_{n}\left(  \boldsymbol{\theta}\right)  };0\leq\theta_{j}%
\leq\frac{\pi}{2} \label{ntheta}%
\end{equation}
$\lceil$%
\begin{equation}
\sum\limits_{1\leq j\leq s}\frac{\partial\mathcal{E}_{AGP}\left(
\boldsymbol{\varphi}\left(  \boldsymbol{n}\right)  \right)  }{\partial n_{j}%
}\frac{\partial n_{j}}{\partial\theta_{k}}=0;1\leq k\leq s.
\end{equation}
Using eq.$\left(  \ref{ntheta}\right)  $ and observing that
\begin{equation}
\frac{\partial n_{j}}{\partial\theta_{k}}=\frac{\cos^{2}\theta_{j}\sin
2\theta_{k}}{\mathcal{N}\left(  \boldsymbol{\theta}\right)  ^{2}}-\frac
{\sin2\theta_{k}\delta_{jk}}{\mathcal{N}\left(  \boldsymbol{\theta}\right)
}=\frac{\sin2\theta_{k}}{\mathcal{N}\left(  \boldsymbol{\theta}\right)
}\left\{  \frac{\cos^{2}\theta_{j}}{\mathcal{N}\left(  \boldsymbol{\theta
}\right)  }-\delta_{jk}\right\}
\end{equation}
and
\begin{equation}
\frac{\partial\mathcal{E}_{AGP}\left(  \boldsymbol{\varphi}\left(
\boldsymbol{n}\left(  \boldsymbol{\theta}\right)  \right)  \right)  }{\partial
n_{j}}=\varepsilon_{jj}\left(  \boldsymbol{n}\left(  \boldsymbol{\theta
}\right)  \right)
\end{equation}
this can be developed to give
\begin{align}
&  \frac{\partial\mathcal{E}_{AGP}\left(  \boldsymbol{\varphi}\left(
\boldsymbol{\theta}\right)  \right)  }{\partial\theta_{k}}=\sum\limits_{1\leq
j\leq s}\varepsilon_{jj}\left(  \boldsymbol{\theta}\right)  \frac{\sin
2\theta_{k}}{\mathcal{N}\left(  \boldsymbol{\theta}\right)  }\left\{
\frac{\cos^{2}\theta_{j}}{\mathcal{N}\left(  \boldsymbol{\theta}\right)
}-\delta_{jk}\right\}  =0\nonumber\\
&  \Rightarrow\sum\limits_{1\leq j\leq s}\varepsilon_{jj}\left(
\boldsymbol{\theta}\right)  \frac{\cos^{2}\theta_{j}\sin2\theta_{k}%
}{\mathcal{N}\left(  \boldsymbol{\theta}\right)  ^{2}}=\varepsilon_{kk}\left(
\boldsymbol{\theta}\right)  \sin2\theta_{k}%
\end{align}
$\rfloor$
\end{enumerate}

We are interested in the dependence on the occupation numbers of the
$2M$-electron state which we can obtain by using the chain rule
\begin{align}
&  \frac{\partial\mathcal{E}_{AGP}\left(  \boldsymbol{\varphi}\left(
\boldsymbol{n}\right)  \right)  }{\partial N_{j}}=\sum_{1\leq k\leq2s}%
\frac{\partial\mathcal{E}_{AGP}\left(  \boldsymbol{\varphi}\left(
\boldsymbol{n}\right)  \right)  }{\partial n_{k}}\frac{\partial n_{k}%
}{\partial N_{j}}\nonumber\\
&  \Rightarrow\frac{\partial\mathcal{E}_{AGP}\left(  \boldsymbol{\varphi
}\left(  \boldsymbol{n}\right)  \right)  }{\partial N_{j}}=\sum_{1\leq
k\leq2s}\frac{\partial\mathcal{E}_{AGP}\left(  \boldsymbol{\varphi}\left(
\boldsymbol{n}\right)  \right)  }{\partial n_{k}}\frac{\partial n_{k}%
}{\partial N_{j}}=\sum_{1\leq k\leq2s}\frac{\partial\mathcal{E}_{AGP}\left(
\boldsymbol{\varphi}\left(  \boldsymbol{n}\right)  \right)  }{\partial n_{k}%
}\left(  \frac{\partial N_{k}}{\partial n_{j}}\right)  ^{-1}%
\end{align}
i.e.%
\begin{equation}
\frac{\partial\mathcal{E}_{AGP}\left(  \boldsymbol{\varphi}\left(
\boldsymbol{n}\right)  \right)  }{\partial\boldsymbol{N}}=\frac{\partial
\mathcal{E}_{AGP}\left(  \boldsymbol{\varphi}\left(  \boldsymbol{n}\right)
\right)  }{\partial\boldsymbol{n}}\left(  \frac{\partial\boldsymbol{N}%
}{\partial\boldsymbol{n}}\right)  ^{-1}.
\end{equation}
The Jacobian $\frac{\partial\boldsymbol{N}}{\partial\boldsymbol{n}}$ is
nonsingular as the occupation numbers of the AGP\ state are explicitly given
in a $1-1$ fashion by $\left[  {}\right]  $
\begin{equation}
\boldsymbol{N}=\mathcal{N}\left(  \boldsymbol{n}\right)  ,
\end{equation}
where
\begin{equation}
N_{j}=\mathcal{N}_{j}\left(  \boldsymbol{n}\right)  =n_{j}\frac{S_{M-1}\left(
\jmath\right)  }{S_{M}}=n_{j}\frac{\partial\ln S_{M}}{\partial n_{j}}%
;\quad1\leq j\leq s.
\end{equation}
The matrix elements of the Jacobian $\left(  \frac{\partial\boldsymbol{N}%
}{\partial\boldsymbol{n}}\right)  $ are
\begin{align}
\frac{\partial N_{k}}{\partial n_{j}}  &  =\frac{\partial}{\partial n_{j}%
}\left\{  n_{k}\frac{\partial\ln S_{M}}{\partial n_{k}}\right\}  =\delta
_{jk}\frac{\partial\ln S_{M}}{\partial n_{k}}+n_{k}\frac{\partial^{2}\ln
S_{M}}{\partial n_{j}\partial n_{k}}\nonumber\\
&  =\delta_{jk}\frac{S_{M-1}\left(  \jmath\right)  }{S_{M}}-n_{k}\frac
{S_{M-1}\left(  \jmath\right)  S_{M-1}\left(  k\right)  }{S_{M}^{2}}%
+n_{k}\frac{S_{M-2}\left(  jk\right)  }{S_{M}}.
\end{align}


\section{Discussion\label{Discussion}}

\appendix


\section{State Density, Reduced Density Operators and Operator Spaces
\label{Denmat}}

\subsection{State Density Operators and Operator Spaces}

The convex set $S_{N}$ of normalized $N$-electron density operators $D$ is
defined as
\begin{equation}
S_{N}=\left\{  D\left\vert D\in\mathcal{B}_{1+}\left(  \mathcal{H}^{N}\right)
;\operatorname{Tr}D=1\right.  \right\}  ,
\end{equation}
where $\mathcal{B}_{1+}\left(  \mathcal{H}^{N}\right)  $ is the cone (see
below) of positive trace class operators, defined as
\begin{equation}
\mathcal{B}_{1+}\left(  \mathcal{H}^{N}\right)  =\left\{  D\in\mathcal{B}%
_{1}\left(  \mathcal{H}^{N}\right)  \left\vert \overset{}{\left\langle
\Psi\left\vert D\Psi\right.  \right\rangle }\geq0\quad\forall\Psi
\in\mathcal{H}^{N}\right.  \right\}  ,
\end{equation}
where $\mathcal{B}_{1}\left(  \mathcal{H}^{N}\right)  $ is the space of trace
class operators defined as
\begin{equation}
\mathcal{B}_{1}\left(  \mathcal{H}^{N}\right)  =\left\{  \left.  X\right\vert
\operatorname{Tr}\left\{  \left(  X^{\dagger}X\right)  ^{\frac{1}{2}}\right\}
<\infty\right\}  ,
\end{equation}
which is contained in the space of bounded $N$-electron operators
$\mathcal{B}\left(  \mathcal{H}^{N}\right)  $ defined as
\begin{align}
\mathcal{B}\left(  \mathcal{H}^{N}\right)   &  =\left\{  X\left\vert
\sup_{\substack{\Psi\in\mathcal{H}^{N}\\\left\Vert \Psi\right\Vert
=1}}\left\{  \left\langle \Psi\left\vert X^{\dagger}X\Psi\right.
\right\rangle \right\}  <\infty=\sup_{\substack{\Psi\in\mathcal{H}%
^{N}\\\left\Vert \Psi\right\Vert =1}}\left\{  \left\Vert X\Psi\right\Vert
\right\}  <\infty\right.  \right. \\
&  =\left.  \sup_{\substack{\Psi,\Phi\in\mathcal{H}^{N}\\\left\Vert
\Psi\right\Vert =1,\left\Vert \Phi\right\Vert =1}}\left\{  \left\langle
\Psi\left\vert X^{\dagger}X\Phi\right.  \right\rangle \right\}  <\infty
\right\}  .
\end{align}
It can be shown that if one considers bounded self-adjoint operators then
\begin{equation}
\left\Vert X\right\Vert =\sup_{\substack{\Psi\in\mathcal{H}^{N}\\\left\Vert
\Psi\right\Vert =1}}\left\{  \left\vert \left\langle \Psi\left\vert
X\Psi\right.  \right\rangle \right\vert \right\}  =r\left(  X\right)  <\infty,
\end{equation}
where $r\left(  X\right)  $ is the spectral radius of $X$.

The convex set property of $S_{N}$ is explicitly displayed by
\begin{equation}
S_{N}=\left\{  \left.  \sum_{1\leq j\leq m}\alpha_{j}D_{j}\right\vert
\sum_{1\leq j\leq m}\alpha_{j}=1;0\leq\alpha_{j}\leq1,D_{j}\in S_{N},1\leq
j\leq m\right\}  .
\end{equation}
and a set $C$ is a cone if $X\in C\Leftrightarrow\alpha X\in C$ $\forall
\alpha>0,$ which is satisfied by $\mathcal{B}_{1+}\left(  \mathcal{H}%
^{N}\right)  .$

The spaces $\mathcal{B}\left(  \mathcal{H}^{N}\right)  $ and $\mathcal{B}%
_{1}\left(  \mathcal{H}^{N}\right)  $ are Banach spaces with norms defined by
\begin{equation}
\left\Vert X\right\Vert =\left[  \sup_{\substack{\Psi\in\mathcal{H}^{N}%
\\\Psi\neq0}}\left\{  \frac{\left\langle \Psi\left\vert X^{\dagger}%
X\Psi\right.  \right\rangle }{\left\langle \Psi\left\vert \Psi\right.
\right\rangle }\right\}  \right]  ^{\frac{1}{2}}=\sup_{_{\substack{\Phi
,\Psi\in\mathcal{H}^{N}\\\Phi,\Psi\neq0}}}\left\{  \frac{\left\vert
\left\langle \Phi\left\vert X\Psi\right.  \right\rangle \right\vert
}{\left\Vert \Phi\right\Vert \left\Vert \Psi\right\Vert }\right\}
\end{equation}
and
\begin{equation}
\left\Vert X\right\Vert _{1}=\operatorname{Tr}\left\{  \left(  X^{\dagger
}X\right)  ^{\frac{1}{2}}\right\}
\end{equation}
respectively. A closely related space that we have utilized is the space of
Hilbert Schimdt operators $\mathcal{B}_{2}\left(  \mathcal{H}^{N}\right)  $
defined as
\begin{equation}
\mathcal{B}_{2}\left(  \mathcal{H}^{N}\right)  =\left\{  X\left\vert
\operatorname{Tr}\left\{  X^{\dagger}X\right\}  <\infty\right.  \right\}  ,
\end{equation}
which is a Hilbert space with inner product given by
\begin{equation}
\left\langle X\left\vert Y\right.  \right\rangle _{\mathcal{B}_{2}\left(
\mathcal{H}^{N}\right)  }=\operatorname{Tr}\left\{  X^{\dagger}Y\right\}
\end{equation}
and norm by
\begin{equation}
\left\Vert X\right\Vert _{2}=\left(  \operatorname{Tr}\left\{  X^{\dagger
}X\right\}  \right)  ^{\frac{1}{2}}=\left\langle X\left\vert X\right.
\right\rangle _{\mathcal{B}_{2}\left(  \mathcal{H}^{N}\right)  }^{\frac{1}{2}%
}.
\end{equation}


As $\left\Vert X\right\Vert \leq\left\Vert X\right\Vert _{2}\leq\left\Vert
X\right\Vert _{1}$ \cite{KummerI} on has
\begin{equation}
\mathcal{B}\left(  \mathcal{H}^{N}\right)  \subseteq\mathcal{B}_{2}\left(
\mathcal{H}^{N}\right)  \subseteq\mathcal{B}_{1}\left(  \mathcal{H}%
^{N}\right)  .
\end{equation}
One can also note that the spaces of compact of operators $\mathcal{C}\left(
\mathcal{H}^{N}\right)  ,$Hilbert Schmidt operators and trace class operators
are the completions of the space of finite rank operators, $\mathcal{B}%
_{0}\left(  \mathcal{H}^{N}\right)  ,$ in the operator norm, $\left\Vert
{}\right\Vert ,$ the Hilbert-Schmidt norm $\left\Vert {}\right\Vert _{2}$ and
the trace class norm $\left\Vert X\right\Vert _{1}$ respectively, so that one
obtains
\begin{equation}
\mathcal{B}\left(  \mathcal{H}^{N}\right)  \subseteq\mathcal{C}\left(
\mathcal{H}^{N}\right)  \subseteq\mathcal{B}_{2}\left(  \mathcal{H}%
^{N}\right)  \subseteq\mathcal{B}_{1}\left(  \mathcal{H}^{N}\right)
\subseteq\mathcal{B}_{0}\left(  \mathcal{H}^{N}\right)  .
\end{equation}


\subsection{Reduced Density Operators \label{RDO}}

\qquad The space-spin integral kernel, $D^{p}\left(  \boldsymbol{\kappa}%
_{p}\boldsymbol{,\kappa}_{p}^{\prime}\right)  $ of the $p^{th}$ order reduced
density operator of the $N$-electron state, $D,$ is defined as%
\begin{equation}
D^{p}\left(  \boldsymbol{\kappa}_{p};\boldsymbol{\kappa}_{p}^{\prime}\right)
=\operatorname{Tr}\left\{  \Phi\left(  \boldsymbol{\kappa}_{p}^{\prime
}\right)  ^{\dagger}\Phi\left(  \boldsymbol{\kappa}_{p}\right)  D\right\}  ,
\label{DenMatDef}%
\end{equation}
where
\begin{align}
\boldsymbol{\kappa}_{p}  &  \equiv\boldsymbol{x}_{1},\cdots,\boldsymbol{x}%
_{p}\\
\boldsymbol{x}_{j}  &  \equiv\boldsymbol{r}_{j}\boldsymbol{,\xi}_{j};1\leq
j\leq p\nonumber
\end{align}
and the field operators are defined in Appendix $\left(  \ref{SecQuant}%
\right)  .$ The multi-electron operators $\Phi\left(  \boldsymbol{\kappa}%
_{p}\right)  ^{\dagger}$ are defined as the product of the one-electrons
operators $\Phi\left(  \boldsymbol{\kappa}_{j}\right)  ^{\dagger}$ i.e.%
\begin{equation}
\Phi\left(  \boldsymbol{\kappa}_{p}\right)  ^{\dagger}=\Phi\left(
\boldsymbol{x}_{1}\right)  ^{\dagger}\cdots\Phi\left(  \boldsymbol{x}%
_{p}\right)  ^{\dagger}.
\end{equation}
The space integral kernels are also given by eq.$\left(  \ref{DenMatDef}%
\right)  $ but with
\begin{equation}
\boldsymbol{\kappa}_{p}\equiv\boldsymbol{r}_{1},\cdots,\boldsymbol{r}_{p}.
\end{equation}
The $p^{th}$-Order RDO is then given by
\begin{equation}
D^{p}=\frac{1}{\left(  p!\right)  ^{2}}\int D^{p}\left(  \boldsymbol{\kappa
}_{p};\boldsymbol{\kappa}_{p}^{\prime}\right)  \left\vert \boldsymbol{\kappa
}_{p}\right\rangle \left\langle \boldsymbol{\kappa}_{p}^{\prime}\right\vert
d\boldsymbol{\kappa}_{p}d\boldsymbol{\kappa}_{p}^{\prime},
\end{equation}
where $\left\{  \left\vert \boldsymbol{\kappa}_{p}\right\rangle ,\left\langle
\boldsymbol{\kappa}_{p}^{\prime}\right\vert \right\}  $ are to be understood
as generalized bases (Appendix $\left(  \ref{Genbasis}\right)  ).$

The RDO kernel for a state, $D$, that does not have a fixed number of
electrons can be given in terms of the kernels for pure number states as
\begin{equation}
D\left(  \boldsymbol{\kappa};\boldsymbol{\kappa}^{\prime}\right)  =\sum_{0\leq
p\leq\infty}\alpha_{p}D^{p}\left(  \boldsymbol{\kappa}_{p};\boldsymbol{\kappa
}_{p}^{\prime}\right)  ;\sum_{1\leq p\leq\infty}\alpha_{j}=1;0\leq\alpha
_{j}\leq1
\end{equation}


\subsection{Reduced Density Matrices (RDM) \label{RDM}}

The matrix elements $D_{j_{1}\cdots j_{p}k_{1}\cdots k_{p}}^{p}$ of the
$p^{th}$-Order RDM $\mathbb{D}^{p}$ with respect to the complete orthonormal
basis $\left\{  \left.  \left\vert \varphi_{j}\right\rangle \right\vert 1\leq
j_{1}\leq\cdots\leq j_{p}\leq\infty\right\}  $ of $\mathcal{H}^{1}$ are given
by%
\begin{align}
D_{j_{1}\cdots j_{p}k_{1}\cdots k_{p}}^{p}  &  =\int\varphi_{j_{1}}\left(
\boldsymbol{x}_{1}\right)  ^{\ast}\cdots\varphi_{j_{p}}\left(  \boldsymbol{x}%
_{p}\right)  ^{\ast}D^{p}\left(  \boldsymbol{\kappa}_{p};\boldsymbol{\kappa
}_{p}^{\prime}\right)  \varphi_{k_{1}}\left(  \boldsymbol{x}_{1}^{\prime
}\right)  \cdots\varphi_{k_{p}}\left(  \boldsymbol{x}_{p}^{\prime}\right)
d\boldsymbol{\kappa}_{p}d\boldsymbol{\kappa}_{p}^{\prime}\\
1  &  \leq j_{1}\leq\cdots\leq j_{p}\leq\infty;1\leq k_{1}\leq\cdots\leq
k_{p}\leq\infty.
\end{align}
Thus
\begin{equation}
D^{p}=\sum_{\substack{1\leq j_{1}\leq\cdots\leq j_{p}\leq\infty\\1\leq
k_{1}\leq\cdots\leq k_{p}\leq\infty}}D_{j_{1}\cdots j_{p}k_{1}\cdots k_{p}%
}^{p}\left\vert \varphi_{j_{1}}\cdots\varphi_{j_{p}}\right\rangle \left\langle
\varphi_{k_{1}}\cdots\varphi_{k_{p}}\right\vert .
\end{equation}
$\lceil\left\vert \varphi_{j_{1}}\cdots\varphi_{j_{p}}\right\rangle =\frac
{1}{\sqrt{p!}}\varphi_{j_{1}}\wedge\cdots\wedge\varphi_{j_{p}}=\mathcal{A}%
_{p}\varphi_{j_{1}}\otimes\cdots\otimes\varphi_{j_{p}}$ thus the matrix
elements are given by expectation value with respect to the simple tensor
product i.e.
\begin{align}
&  \left\langle \left.  \varphi_{j_{1}}\cdots\varphi_{j_{p}}\right\vert
D^{p}\varphi_{k_{1}}\cdots\varphi_{k_{p}}\right\rangle =\frac{1}%
{p!}\left\langle \left.  \varphi_{j_{1}}\wedge\cdots\wedge\varphi_{j_{p}%
}\right\vert D^{p}\varphi_{k_{1}}\wedge\cdots\wedge\varphi_{k_{p}%
}\right\rangle \\
&  =\frac{1}{p!}\left\langle \left.  \mathcal{A}_{p}\varphi_{j_{1}}%
\otimes\cdots\otimes\varphi_{j_{p}}\right\vert D^{p}\mathcal{A}_{p}%
\varphi_{k_{1}}\otimes\cdots\otimes\varphi_{k_{p}}\right\rangle =\frac{1}%
{p!}\left\langle \left.  \varphi_{j_{1}}\otimes\cdots\otimes\varphi_{j_{p}%
}\right\vert \mathcal{A}_{p}D^{p}\mathcal{A}_{p}\varphi_{k_{1}}\otimes
\cdots\otimes\varphi_{k_{p}}\right\rangle \\
&  =\frac{1}{p!}\left\langle \left.  \varphi_{j_{1}}\otimes\cdots
\otimes\varphi_{j_{p}}\right\vert D^{p}\varphi_{k_{1}}\otimes\cdots
\otimes\varphi_{k_{p}}\right\rangle \rfloor
\end{align}


\subsection{Densities \label{Gen. Den}}

\qquad In general we define the density, $\eta_{A}$, of an operator
$A\in\mathcal{B}_{1}\left(  \mathcal{H}^{N}\right)  $ to be
\begin{equation}
\eta_{A}\left(  \boldsymbol{r}\right)  =\operatorname{Tr}\left\{  \Phi\left(
\boldsymbol{r}\right)  ^{\dagger}\Phi\left(  \boldsymbol{r}\right)  A\right\}
;\text{ }\boldsymbol{r\in}\mathbb{R}^{3},
\end{equation}
which can be complex valued, $A=A^{\dagger}$ then $\eta_{A}$ is real valued
but can have both positive and negative values, while if $A\geq0$ then
$\eta_{A}\left(  \boldsymbol{r}\right)  \geq0$ and we denote it by the
positive density $\rho_{A}$ i.e.
\begin{equation}
\rho_{A}\left(  \boldsymbol{r}\right)  =\eta_{A}\left(  \boldsymbol{r}\right)
=\operatorname{Tr}\left\{  \Phi\left(  \boldsymbol{r}\right)  ^{\dagger}%
\Phi\left(  \boldsymbol{r}\right)  A\right\}  \geq0;\text{ }A\geq
0,\boldsymbol{r\in}\mathbb{R}^{3}.
\end{equation}
If $A$ is self adjoint then $A$ has the generalized expansion (see Appendix
\ref{Genbasis})
\begin{equation}
A=\int\left\vert a\right\rangle \left\langle a\right\vert da
\end{equation}
and
\begin{equation}
\eta_{A}\left(  \boldsymbol{r}\right)  =0\Leftrightarrow\Phi\left(
\boldsymbol{r}\right)  ^{\dagger}\Phi\left(  \boldsymbol{r}\right)  \left\vert
a\right\rangle =0
\end{equation}
as%
\begin{equation}
\left(  \Phi\left(  \boldsymbol{r}\right)  ^{\dagger}\Phi\left(
\boldsymbol{r}\right)  \right)  ^{2}=\Phi\left(  \boldsymbol{r}\right)
^{\dagger}\Phi\left(  \boldsymbol{r}\right)  ,
\end{equation}
leading to
\begin{equation}
\left\langle a\left\vert \Phi\left(  \boldsymbol{r}\right)  ^{\dagger}%
\Phi\left(  \boldsymbol{r}\right)  a\right.  \right\rangle =\left\langle
a\left\vert \left(  \Phi\left(  \boldsymbol{r}\right)  ^{\dagger}\Phi\left(
\boldsymbol{r}\right)  \right)  ^{2}a\right.  \right\rangle =\left\langle
\Phi\left(  \boldsymbol{r}\right)  ^{\dagger}\Phi\left(  \boldsymbol{r}%
\right)  a\left\vert \Phi\left(  \boldsymbol{r}\right)  ^{\dagger}\Phi\left(
\boldsymbol{r}\right)  a\right.  \right\rangle =0.
\end{equation}


\section{Generalized Bases \label{Genbasis}}

Generalized bases make use of the "rigged Hilbert space" construction
\cite{Bohm}. A rigging of a Hilbert space $\mathcal{H}$ is produced by a dense
(in the Hilbert space norm) Topological Vector Space (TVS) space
$\mathcal{R\subset H},$ which has a dual $\mathcal{R}^{d}$ such that
$\mathcal{H=H}^{d}$ can be embedded in $\mathcal{R}^{d}$ i.e.
$\mathcal{H\subset}$ $\mathcal{R}^{d}$ produing a Gelfand Triple
$\mathcal{R\subset H\subset R}^{d}.$ Many topologies may be put on
$\mathcal{R}$ (that are finer than the Hilbert space topology) producing
different $\mathcal{R}^{d}$'s, the only restriction is that the embedding
$\mathit{i}$ $:\mathcal{H\rightarrow}$ $\mathcal{R}^{d}$ is antilinear,
injective and continuous. Different $\mathcal{R}$'s produce different
riggings. A space that is often chosen for $\mathcal{R}$ is $\mathcal{S}%
\left(  \mathbb{R}^{n}\right)  $ the space of rapidly decreasing antisymmetric
test functions on $\mathbb{R}^{n},$ its dual space $\mathcal{S}\left(
\mathbb{R}^{n}\right)  ^{d}$ is the space of antisymmetric tempered
distributions, which are generated by the one-electron space that contains
Dirac delta functions, $\delta_{\mathbf{r}}\left(  \mathbf{r}^{\prime}\right)
=\delta\left(  \mathbf{r-r}^{\prime}\right)  ,$ and plane waves $w_{\mathbf{k}%
}\left(  \mathbf{r}^{\prime}\right)  =e^{i\mathbf{k\cdot r}^{\prime}}$
\cite{Bogolubov}, where $\mathbf{r}^{\prime}\in$ $\mathbb{R}^{3}$ is the
argument of the distributions indexed by $\mathbf{r}$ and $\mathbf{k}$.

As $\mathcal{R}^{d}$ is the dual of $\mathcal{R}$ one denotes the elements of
$\mathcal{R}^{d}$ as $\left\langle a^{\prime}\right\vert $ and those of
$\mathcal{R}$ as $\left\vert a\right\rangle .\mathcal{R}$ and $\mathcal{H}$
can be embedded in $\mathcal{R}^{d}$ thus one can always consider
$\left\langle a\right\vert ;$ $a\in\mathcal{R}$ and $\left\langle h\right\vert
;$ $h\in\mathcal{H}$ with the understanding that $\left\langle a\right\vert $
are defined on $\mathcal{R}$ and $\left\langle h\right\vert $ on $\mathcal{H}$.

The Nuclear Spectral Theorem \cite{Bohm} states that any self-adjoint operator
$A$ (in general not bounded) has an expansion
\begin{equation}
A=\int_{\sigma\left(  A\right)  }a\left\vert a\right\rangle \left\langle
a\right\vert da,
\end{equation}
$\lceil$where the measure $da$ can be expanded as
\begin{equation}
da=da_{ac}+da_{pp}+da_{sc}=da_{ac}+da_{sc}+\sum_{1\leq j\leq r}\delta\left(
\alpha_{j}-a\right)  da.
\end{equation}
One should observe that if there are sets/points in the spectrum that have
measure zero with respect to the measure $da_{ac}$ that is absolutely
continuous with respect to Lebesgue measure then they have non zero measure
with respect $da_{sc},$ the measure that is singularly continuous with respect
to Lebesgue measure.$\rfloor$

in the sense that
\begin{align}
\left\langle \varphi\left\vert A\psi\right.  \right\rangle  &  =\int%
_{\sigma\left(  A\right)  }a\left\langle \left.  \varphi\right\vert
a\right\rangle \left\langle \left.  a\right\vert \psi\right\rangle da\\
\left\langle \left.  \varphi\right\vert a\right\rangle  &  =\overline
{\left\langle \left.  a\right\vert \varphi\right\rangle },
\end{align}
where $\left\{  \left\vert \varphi\right\rangle ,\left\langle \psi\right\vert
\right\}  \in\mathcal{R},\left\{  \left\vert a\right\rangle ,\left\langle
a\right\vert \right\}  \in\mathcal{R}^{d}$, the spectrum of $\sigma\left(
A\right)  \subseteq\mathbb{R}$ and $da$ is a point measure. All of the vectors
$\left\vert a\right\rangle $ are eigenvectors of $A$ i.e.
\begin{equation}
A\left\vert a\right\rangle =a\left\vert a\right\rangle ,
\end{equation}
again in the sense%
\begin{equation}
\left\langle \varphi\left\vert Aa\right.  \right\rangle =a\left\langle
\varphi\left\vert a\right.  \right\rangle =a\overline{\left\langle
\varphi\left\vert a\right.  \right\rangle },\left\vert \varphi\right\rangle
\in\mathcal{R}.
\end{equation}
The set of vectors $\left\{  \left.  \left\vert a\right\rangle \right\vert
a\in\sigma\left(  A\right)  \right\}  $ forms a generalized basis of
$\mathcal{H}$ and any vector $\left\vert v\right\rangle \in\mathcal{R}$ has
the expansion
\begin{equation}
\left\vert v\right\rangle =\int_{\sigma\left(  A\right)  }\left\langle \left.
a\right\vert v\right\rangle \left\vert a\right\rangle da;\left\vert
a\right\rangle \in\mathcal{R}^{d}%
\end{equation}
The basis $\left\{  \left.  \left\vert a\right\rangle \right\vert a\in
\sigma\left(  A\right)  \right\}  $ is orthonormal in the sense
\begin{equation}
\left\langle \left.  a_{1}\right\vert \nu\right\rangle =\int_{\sigma\left(
A\right)  }\left\langle \left.  a_{1}\right\vert a_{2}\right\rangle
\left\langle \left.  a_{2}\right\vert v\right\rangle da_{2}\forall\left\vert
v\right\rangle \in\mathcal{R}%
\end{equation}
so that $\left\langle \left.  a_{1}\right\vert a_{2}\right\rangle $ can be
identified as
\begin{equation}
\left\langle \left.  a_{1}\right\vert a_{2}\right\rangle =\delta_{a_{1}-a_{2}%
}.
\end{equation}
The operators $\left\vert a\right\rangle \left\langle a^{\prime}\right\vert $
can thus be interpretted as
\begin{equation}
\left\vert a\right\rangle \left\langle a^{\prime}\right\vert :\mathcal{R}%
^{d}\rightarrow\mathcal{R}^{d}.
\end{equation}
The vectors $\left\vert v\right\rangle \in\mathcal{R}$ have unique expansion
coefficients $\left\langle \left.  a\right\vert v\right\rangle $ (if the
spectral measure has no continuously singular parts) i.e.
\begin{align}
\left\vert v_{1}\right\rangle  &  =\left\vert v_{2}\right\rangle
\Leftrightarrow\left\langle \left.  a\right\vert v_{1}\right\rangle
=\left\langle \left.  a\right\vert v_{2}\right\rangle \forall a\nonumber\\
\left\langle a^{\prime}\left\vert \int_{\sigma\left(  A\right)  }\left\langle
\left.  a\right\vert v_{1}\right\rangle \left\vert a\right\rangle da\right.
\right\rangle  &  =\left\langle a^{\prime}\left\vert \int_{\sigma\left(
A\right)  }\left\langle \left.  a\right\vert v_{2}\right\rangle \left\vert
a\right\rangle da\right.  \right\rangle \nonumber\\
\int_{\sigma\left(  A\right)  }\left\langle \left.  a\right\vert
v_{1}\right\rangle \delta\left(  a^{\prime}-a\right)  da  &  =\int%
_{\sigma\left(  A\right)  }\left\langle \left.  a\right\vert v_{2}%
\right\rangle \delta\left(  a^{\prime}-a\right)  da\nonumber\\
\left\langle \left.  a^{\prime}\right\vert v_{1}\right\rangle  &
=\left\langle \left.  a^{\prime}\right\vert v_{2}\right\rangle .
\end{align}
Thus%
\begin{equation}
\left\vert v_{1}\right\rangle -\left\vert v_{2}\right\rangle =0\Rightarrow
\int_{\sigma\left(  A\right)  }\left\{  \left\langle \left.  a\right\vert
v_{1}\right\rangle -\left\langle \left.  a\right\vert v_{2}\right\rangle
\right\}  \delta\left(  a^{\prime}-a\right)  da=0
\end{equation}
in contrast to
\begin{equation}
\int_{\sigma\left(  A\right)  }\left\vert v_{1}\left(  a\right)  -v_{2}\left(
a\right)  \right\vert ^{2}da=0\Leftrightarrow v_{1}\left(  a\right)
\overset{a.e}{=}v_{2}\left(  a\right)  \text{ wrt }da
\end{equation}
as $\left\vert v_{1}\right\rangle $ and $\left\vert v_{2}\right\rangle $ do
not have to be in $\mathcal{R}$.

If the operator $A$ is a compact \cite{Encyclo} and self-adjoint (thus
bounded) then the basis $\left\{  \left.  \left\vert a\right\rangle
\right\vert a\in\sigma\left(  A\right)  ;\right\}  $ is a conventional
orthonormal basis contained in $\mathcal{H}$ i.e. $\left\vert a\right\rangle
\in\mathcal{H\,}\forall a\in\sigma\left(  A\right)  .$

In particular consider the position operator $\mathbf{Q},($which is not
bounded),
\begin{equation}
\mathbf{Q}=\int_{\mathbb{R}^{3}}\mathbf{r}\left\vert \mathbf{r}\right\rangle
\left\langle \mathbf{r}\right\vert d\mathbf{r,}%
\end{equation}
then any state, $\left\vert \varphi\right\rangle ,$ that has a wavefunction,
$\left\langle \left.  \mathbf{r}\right\vert \varphi\right\rangle ,$ belonging
to the test function space $\mathcal{R}$ $=\mathcal{S}\left(  \mathbb{R}%
^{n}\right)  ,$ can be expanded as
\begin{equation}
\left\vert \varphi\right\rangle =\int_{\mathbb{R}^{3}}\left\langle \left.
\mathbf{r}\right\vert \varphi\right\rangle \left\vert \mathbf{r}\right\rangle
d\mathbf{r,}%
\end{equation}
where
\begin{equation}
\left\langle \left.  \mathbf{r}\right\vert \mathbf{r}^{\prime}\right\rangle
=\delta_{\mathbf{r-r}^{\prime}}.
\end{equation}
(Note that $\mathcal{R}$ $\neq\mathcal{H}$ as there are functions in
$\mathcal{H}$ that are not bounded on compact sets and hence $\left\langle
\left.  \mathbf{r}\right\vert \varphi\right\rangle $ is not defined for all
$\left\vert \varphi\right\rangle \in\mathcal{H}$).

$\lceil$%
\begin{equation}
\left\vert \varphi\right\rangle =\left\vert \psi\right\rangle \text{
}\Leftrightarrow\left\{  \left\langle \left.  \varphi\right\vert
\chi\right\rangle =\left\langle \left.  \psi\right\vert \chi\right\rangle
\equiv\left\langle \left.  \varphi\right\vert \chi\right\rangle -\left\langle
\left.  \psi\right\vert \chi\right\rangle =0\right\}  \forall\chi\mathbf{\in
}\text{ }\mathcal{R}%
\end{equation}
$\rfloor$

\section{Second Quantized Density Operators \label{SecQuant}}

Fermion second quantized operators $\left\{  \left.  a_{j}^{\dagger
}\right\vert 1\leq j\leq\infty\right\}  $ are defined by the action on the
vacuum state $\left\vert \phi\right\rangle $ by
\begin{equation}
a_{j}^{\dagger}\left\vert \phi\right\rangle =\left\vert \varphi_{j}%
\right\rangle ,1\leq j\leq\infty,
\end{equation}
where $\left\{  \left.  \left\vert \varphi_{j}\right\rangle \right\vert 1\leq
j\leq\infty\right\}  $ is a complete orthonormal basis of one-electron space
$\mathcal{H}^{1}$ and $\left\{  \left.  a_{j}^{\dagger},a_{k}\right\vert 1\leq
j,k\leq\infty\right\}  $ satisfy the canonical anti-commutation relationships
\begin{equation}
\left[  a_{j}^{\dagger},a_{k}\right]  _{+}=a_{j}^{\dagger}a_{k}+a_{k}%
a_{j}^{\dagger}=\delta_{jk};1\leq j,k\leq\infty.
\end{equation}
The operators $\left\{  \left.  a_{j,}a_{j}^{\dagger}\right\vert 1\leq
j\leq\infty\right\}  $ can be considered as the values of the maps
\begin{equation}
\Phi,\Phi^{\dagger}:\mathcal{H}^{1}\rightarrow\mathcal{B}\left(
\mathcal{H}_{\mathcal{F}}\right)
\end{equation}
given by
\begin{align}
a_{j}  &  =\Phi\left(  \varphi_{j}\right)  ,\nonumber\\
a_{j}^{\dagger}  &  =\Phi\left(  \varphi_{j}\right)  ^{\dagger},
\end{align}
where the Fock space, $\mathcal{H}_{\mathcal{F}},$ is defined by
\begin{equation}
\mathcal{H}_{\mathcal{F}}=\underset{0\leq N\leq\infty}{\oplus}\wedge
^{N}\mathcal{H}^{1}.
\end{equation}
The point field operators $\left\{  \left.  \Phi\left(  \boldsymbol{r}\right)
\right\vert \boldsymbol{r\in}\mathbb{R}^{3}\right\}  $ are defined by
\cite{Thirring}
\begin{equation}
\Phi\left(  \boldsymbol{r}\right)  =\sum_{1\leq j\leq\infty}\varphi_{j}\left(
\boldsymbol{r}\right)  ^{\ast}a_{j}%
\end{equation}
and the density operators, $\Phi\left(  \boldsymbol{r}\right)  ^{\dagger}%
\Phi\left(  \boldsymbol{r}\right)  $, formally by
\begin{equation}
\Phi\left(  \boldsymbol{r}\right)  ^{\dagger}\Phi\left(  \boldsymbol{r}%
\right)  =\sum_{1\leq j,k\leq\infty}\varphi_{j}\left(  \boldsymbol{r}\right)
\varphi_{k}\left(  \boldsymbol{r}\right)  ^{\ast}a_{j}^{\dagger}a_{k}.
\end{equation}
The unbounded operators $\Phi\left(  \boldsymbol{r}\right)  $ are densely
defined (but not closeable) on the Fock space, $\mathcal{H}_{\mathcal{F}},$
nor the subspaces $\mathcal{H}^{N}=\wedge^{N}\mathcal{H}^{1},$ while the
domains of their adjoints $\Phi\left(  \boldsymbol{r}\right)  ^{\dagger}$ are
empty i.e. do not contain any vectors belonging to $\mathcal{H}_{\mathcal{F}}$
\cite{Thirring}. However, matrix elements, $\operatorname{Tr}\left\{
A^{\dagger}\Phi\left(  \boldsymbol{r}\right)  ^{\dagger}\Phi\left(
\boldsymbol{r}\right)  B\right\}  ,$ of the density super operators
$\widehat{\Phi\left(  \boldsymbol{r}\right)  ^{\dagger}\Phi\left(
\boldsymbol{r}\right)  }$ can be defined as the distributional limits of
$\operatorname{Tr}\left\{  A^{\dagger}\Phi\left(  f\right)  ^{\dagger}%
\Phi\left(  f\right)  B\right\}  $ as $f\rightarrow\delta_{\boldsymbol{r}}$
where
\begin{equation}
\Phi\left(  \boldsymbol{r}\right)  ^{\dagger}\Phi\left(  \boldsymbol{r}%
\right)  \equiv\Phi\left(  \delta_{\boldsymbol{r}}\right)  ^{\dagger}%
\Phi\left(  \delta_{\boldsymbol{r}}\right)  .
\end{equation}
Even though $\Phi\left(  \boldsymbol{r}\right)  $ are not closeable the
operators $\Phi\left(  \boldsymbol{r}\right)  ^{\dagger}\Phi\left(
\boldsymbol{r}\right)  $ are bounded as
\begin{equation}
\left\Vert \Phi\left(  \boldsymbol{r}\right)  ^{\dagger}\Phi\left(
\boldsymbol{r}\right)  \right\Vert =\left[  \sup_{\substack{\Psi\in
\mathcal{H}_{\mathcal{F}}\\\left\Vert \Psi\right\Vert =1}}\left\{
\left\langle \Psi\left\vert \Phi\left(  \boldsymbol{r}\right)  ^{\dagger}%
\Phi\left(  \boldsymbol{r}\right)  \Psi\right.  \right\rangle \right\}
\right]  ^{\frac{1}{2}}=\sup_{\substack{\Psi\in\mathcal{H}_{\mathcal{F}%
}\\\left\Vert \Psi\right\Vert =1}}\left\{  \rho_{\Psi}\left(  \boldsymbol{r}%
\right)  \right\}  ^{\frac{1}{2}}<\infty
\end{equation}
as all states $\left\vert \Psi\right\rangle \left\langle \Psi\right\vert
\in\mathcal{B}_{1sa+}\left(  \mathcal{H}_{\mathcal{F}}\right)  $ thus
$\rho_{\Psi}\left(  \boldsymbol{r}\right)  \leq\int\rho_{\Psi}\left(
\boldsymbol{r}\right)  d\boldsymbol{r}<\infty.$

\section{Super Operators \label{super}}

The set, $\mathcal{B}_{2sa}\left(  \mathcal{H}^{N}\right)  ,$ of self adjoint
Hilbert-Schmidt operators form a real Hilbert space equipped with inner
product
\begin{equation}
\left(  X\left\vert Y\right.  \right)  =\left\langle X\left\vert Y\right.
\right\rangle _{\mathcal{B}_{2}\left(  \mathcal{H}^{N}\right)  };X=X^{\dagger
},Y=Y^{\dagger}.
\end{equation}
Super-operators (which could be unbounded), $\widehat{T},$ acting on the
Hilbert space $\mathcal{B}_{2sa}\left(  \mathcal{H}^{N}\right)  ,$ i.e.
\begin{equation}
\widehat{T}:\mathcal{B}_{2sa}\left(  \mathcal{H}^{N}\right)  \rightarrow
\mathcal{B}_{2sa}\left(  \mathcal{H}^{N}\right)  ,
\end{equation}
can be generated by self adjoint operators belonging to $\mathcal{B}%
_{2sa}\left(  \mathcal{H}^{N}\right)  $
\begin{equation}
T:\mathcal{H}^{N}\rightarrow\mathcal{H}^{N}%
\end{equation}
by%
\begin{equation}
\widehat{T}\left(  Y\right)  =\frac{1}{2}\left[  T,Y\right]  _{+}\equiv
\frac{1}{2}\left(  TY+YT\right)  .
\end{equation}
Noting that
\begin{equation}
\operatorname{Tr}\left\{  TYX\right\}  =\operatorname{Tr}\left\{  XTY\right\}
=\operatorname{Tr}\left\{  YTX\right\}  =\operatorname{Tr}\left\{
XYT\right\}  \text{ as }X,Y,T\in\mathcal{B}_{2sa}\left(  \mathcal{H}%
^{N}\right)
\end{equation}
one can observe that
\begin{align}
\frac{1}{2}\left(  \operatorname{Tr}\left\{  TYX\right\}  +\operatorname{Tr}%
\left\{  YTX\right\}  \right)   &  =\operatorname{Tr}\left\{  TYX\right\}
\nonumber\\
\frac{1}{2}\left(  \operatorname{Tr}\left\{  XTY\right\}  +\operatorname{Tr}%
\left\{  XYT\right\}  \right)   &  =\operatorname{Tr}\left\{  TYX\right\}
\end{align}
leading to
\begin{equation}
\left(  \left.  \widehat{T}\left(  Y\right)  \right\vert X\right)  =\left(
\left.  TY\right\vert X\right)  .
\end{equation}
The operator $\widehat{T}$ is self adjoint as
\begin{align}
\left(  \widehat{T}X\left\vert Y\right.  \right)   &  =\frac{1}{2}%
\operatorname{Tr}\left\{  \left[  T,X\right]  _{+}Y\right\}  =\frac{1}%
{2}\operatorname{Tr}\left\{  TXY+XTY\right\} \nonumber\\
&  =\frac{1}{2}\operatorname{Tr}\left\{  YTX+TYX\right\}  =\frac{1}%
{2}\operatorname{Tr}\left\{  XTY+XYT\right\} \nonumber\\
&  =\left(  X\left\vert \widehat{T}Y\right.  \right)  .
\end{align}
The domain, $\operatorname{Dom}\left\{  \widehat{T}\right\}  $, of
$\widehat{T}$
\begin{equation}
\operatorname{Dom}\left\{  \widehat{T}\right\}  =\left\{  Y\left\vert
\operatorname{Tr}\left\{  \left(  \widehat{T}\left(  Y\right)  \right)
^{2}\right\}  <\infty\,\right.  \right\}  \equiv\frac{1}{4}\left\{  \left.
Y\right\vert \operatorname{Tr}\left\{  \left(  TY+YT\right)  \left(
TY+YT\right)  \right\}  <\infty\,\right\}
\end{equation}
can be generated by projectors, $\left\vert \Phi\right\rangle \left\langle
\Phi\right\vert $ onto normalized states $\left\vert \Phi\right\rangle $ that
lie in the domain, $\operatorname{Dom}\left\{  T\right\}  $, of $T.$This
domain is dense in $\mathcal{B}_{2sa}\left(  \mathcal{H}^{N}\right)  $ if
$\operatorname{Dom}\left\{  T\right\}  $ is dense in $\mathcal{H}^{N}$ as
\[
\left\{  \left\vert \Phi\right\rangle \left\langle \Phi\right\vert \left\vert
\operatorname{Tr}\left\{  \left(  \widehat{T}\left(  \left\vert \Phi
\right\rangle \left\langle \Phi\right\vert \right)  \right)  ^{2}\right\}
=\frac{1}{2}\left\{  \left\langle \left.  \Phi\right\vert T\Phi\right\rangle
^{2}+\left\langle \left.  \Phi\right\vert T^{2}\Phi\right\rangle \right\}
=\left\Vert T\right\Vert ^{2}<\infty\,\right.  ;\left\vert \Phi\right\rangle
\in\operatorname{Dom}\left\{  T\right\}  \right\}
\]
and the domain of $T$ is generated by the set of HS\ convergent real linear
combinations%
\begin{equation}
Y=\sum_{1\leq j\leq\infty}Y_{j}\left\vert \Phi_{j}\right\rangle \left\langle
\Phi_{j}\right\vert
\end{equation}


The operator $\widehat{T}$ could also be defined in terms of the quadratic
form
\begin{align}
\left(  X\left\vert \widehat{T}Y\right.  \right)   &  =\frac{1}{2}%
\operatorname{Tr}\left\{  X\left[  T,Y\right]  _{+}\right\}  =\frac{1}%
{2}\operatorname{Tr}\left\{  XTY+XYT\right\}  =\,\frac{1}{2}\operatorname{Tr}%
\left\{  XTY+YTX\right\} \nonumber\\
&  =\frac{1}{2}\operatorname{Tr}\left\{  XTY+XTY\right\}  =\operatorname{Tr}%
\left\{  XTY\right\}  \forall X,\,Y\in\operatorname{Dom}\left\{
\widehat{T}\right\}  .
\end{align}


\subsection{Density Super Operators \label{Den Sup Ops}}

The action of the super operator $\widehat{\Phi\left(  \boldsymbol{r}\right)
^{\dagger}\Phi\left(  \boldsymbol{r}\right)  }:\mathcal{B}_{2sa}\left(
\mathcal{H}^{N}\right)  \rightarrow\mathcal{B}_{2sa}\left(  \mathcal{H}%
^{N}\right)  $ is given by
\begin{equation}
\widehat{\Phi\left(  \boldsymbol{r}\right)  ^{\dagger}\Phi\left(
\boldsymbol{r}\right)  }\left\vert A\right)  =\left\vert \left[  \Phi\left(
\boldsymbol{r}\right)  ^{\dagger}\Phi\left(  \boldsymbol{r}\right)  A\right]
_{+}\right)  =\frac{1}{2}\left\vert \left[  \Phi\left(  \boldsymbol{r}\right)
^{\dagger}\Phi\left(  \boldsymbol{r}\right)  ,A\right]  _{+}\right)  .
\label{supproj}%
\end{equation}


\begin{lemma}
If $A=A^{\dagger}$ and $A\in\mathcal{B}_{2sa}\left(  \mathcal{H}^{N}\right)  $
then
\begin{align}
\operatorname*{Tr}\left\{  \Phi\left(  \boldsymbol{r}\right)  ^{\dagger}%
\Phi\left(  \boldsymbol{r}\right)  A^{2}\right\}   &  =0\Leftrightarrow
\Phi\left(  \boldsymbol{r}\right)  ^{\dagger}\Phi\left(  \boldsymbol{r}%
\right)  A\Leftrightarrow A\Phi\left(  \boldsymbol{r}\right)  ^{\dagger}%
\Phi\left(  \boldsymbol{r}\right)  =0\nonumber\\
&  \Leftrightarrow\widehat{\Phi\left(  \boldsymbol{r}\right)  ^{\dagger}%
\Phi\left(  \boldsymbol{r}\right)  }\left\vert A\right)  =0.
\end{align}

\end{lemma}

\begin{proof}
$\Phi\left(  \boldsymbol{r}\right)  ^{\dagger}\Phi\left(  \boldsymbol{r}%
\right)  $ is a projection operator and $A^{2}$ is a positive trace class
operator as $A\in\mathcal{B}_{2sa}\left(  \mathcal{H}^{N}\right)  $ thus
\begin{align}
\operatorname*{Tr}\left\{  \Phi\left(  \boldsymbol{r}\right)  ^{\dagger}%
\Phi\left(  \boldsymbol{r}\right)  A^{2}\right\}   &  =\operatorname*{Tr}%
\left\{  A\Phi\left(  \boldsymbol{r}\right)  ^{\dagger}\Phi\left(
\boldsymbol{r}\right)  \Phi\left(  \boldsymbol{r}\right)  ^{\dagger}%
\Phi\left(  \boldsymbol{r}\right)  A\right\} \nonumber\\
&  =\left\langle A\Phi\left(  \boldsymbol{r}\right)  ^{\dagger}\Phi\left(
\boldsymbol{r}\right)  \left\vert A\Phi\left(  \boldsymbol{r}\right)
^{\dagger}\Phi\left(  \boldsymbol{r}\right)  \right.  \right\rangle
_{\mathcal{B}_{2}\left(  \mathcal{H}^{N}\right)  }=0,
\end{align}
which implies that
\begin{equation}
A\Phi\left(  \boldsymbol{r}\right)  ^{\dagger}\Phi\left(  \boldsymbol{r}%
\right)  =0
\end{equation}
and hence, by taking the adjoint, that
\begin{equation}
\Phi\left(  \boldsymbol{r}\right)  ^{\dagger}\Phi\left(  \boldsymbol{r}%
\right)  A=0.
\end{equation}
Eq. $\left(  \ref{supproj}\right)  $ thus shows that
\begin{equation}
\widehat{\Phi\left(  \boldsymbol{r}\right)  ^{\dagger}\Phi\left(
\boldsymbol{r}\right)  }\left\vert A\right)  =\left\vert \left[  \Phi\left(
\boldsymbol{r}\right)  ^{\dagger}\Phi\left(  \boldsymbol{r}\right)  A\right]
_{+}\right)  =0.
\end{equation}

\end{proof}

Consider a set $\Delta\subseteq\mathbb{R}^{3}$ then one can define a subspace,
$\mathcal{S}_{1\Delta^{c}},$ of $\mathcal{B}_{1sa}\left(  \mathcal{H}%
^{N}\right)  $ by
\begin{equation}
\mathcal{S}_{1\Delta^{c}}=\text{linear span}\left\{  Y\left\vert
\operatorname*{Tr}\left\{  \Phi\left(  \boldsymbol{r}\right)  ^{\dagger}%
\Phi\left(  \boldsymbol{r}\right)  Y\right\}  =0\forall\boldsymbol{r}\in
\Delta^{c}\right.  ,Y\in\mathcal{B}_{1sa}\left(  \mathcal{H}^{N}\right)
\right\}  .
\end{equation}
The space, $\left(  \mathcal{S}_{1\Delta^{c}}\right)  ^{\bot}\subset
\mathcal{B}_{sa}\left(  \mathcal{H}^{N}\right)  ,$ perpendicular to
$\mathcal{S}_{1\Delta^{c}}$ is given by
\begin{equation}
\left(  \mathcal{S}_{1\Delta^{c}}\right)  ^{\bot}=\left\{  X\left\vert
\operatorname*{Tr}\left\{  XY\right\}  =0;\forall Y\in\mathcal{S}_{1\Delta
^{c}}\right.  \right\}  .
\end{equation}
In particular $\left.  \Phi\left(  \boldsymbol{r}\right)  ^{\dagger}%
\Phi\left(  \boldsymbol{r}\right)  \right\vert _{\mathcal{H}^{N}}$
$,\boldsymbol{r}\in\Delta^{c}$ generate a subspace, $\mathcal{C}_{1\Delta^{c}%
},$ of commuting operators of $\left(  \mathcal{S}_{1\Delta^{c}}\right)
^{\bot}$ given by
\begin{equation}
\left.  \mathfrak{f}_{\Delta^{c}}\right\vert _{\mathcal{H}^{N}}=\int%
_{\Delta^{c}}f\left(  \boldsymbol{r}\right)  \left.  \Phi\left(
\boldsymbol{r}\right)  ^{\dagger}\Phi\left(  \boldsymbol{r}\right)
\right\vert _{\mathcal{H}^{N}}d\boldsymbol{r},
\end{equation}
where $f\in L_{\infty}\left(  \Delta\right)  $ and $\mathfrak{f}$ is given by
\begin{equation}
\mathfrak{f}_{\Delta^{c}}=\text{ }\int_{\Delta^{c}}f\left(  \boldsymbol{r}%
\right)  \Phi\left(  \boldsymbol{r}\right)  ^{\dagger}\Phi\left(
\boldsymbol{r}\right)  d\boldsymbol{r}.
\end{equation}%
\begin{equation}
\left\vert \Psi_{1}\right\rangle \left\langle \Psi_{2}\right\vert
\in\mathcal{S}_{1\Delta^{c}}\text{ if }\Phi\left(  \boldsymbol{r}\right)
^{\dagger}\Phi\left(  \boldsymbol{r}\right)  \left\vert \Psi_{1}\right\rangle
=0\text{ or }\Phi\left(  \boldsymbol{r}\right)  ^{\dagger}\Phi\left(
\boldsymbol{r}\right)  \left\vert \Psi_{2}\right\rangle =0
\end{equation}%
\begin{equation}
\mathfrak{f}_{\Delta^{c}}\left\vert \Psi_{1}\right\rangle \left\langle
\Psi_{2}\right\vert =0\forall\left\vert \Psi_{1}\right\rangle \left\langle
\Psi_{2}\right\vert \in\mathcal{S}_{1\Delta^{c}}%
\end{equation}%
\begin{equation}
\Phi\left(  \boldsymbol{r}\right)  ^{\dagger}\Phi\left(  \boldsymbol{r}%
\right)  \left\vert \Psi\right\rangle \left\langle \Psi\right\vert
\neq0\forall\left\vert \Psi\right\rangle \left\langle \Psi\right\vert
\in\mathcal{S}_{1\Delta^{c}},\boldsymbol{r}\in\Delta
\end{equation}%
\begin{equation}
\mathfrak{f}_{\Delta}=\left\langle \text{ }\Psi\left\vert \int_{\Delta
}f\left(  \boldsymbol{r}\right)  \Phi\left(  \boldsymbol{r}\right)  ^{\dagger
}\Phi\left(  \boldsymbol{r}\right)  d\boldsymbol{r}\right.  \Psi\right\rangle
=
\end{equation}
The complimentary operator $\mathfrak{f}_{\Delta}$ is given by
\begin{equation}
\mathfrak{f}_{\Delta}=\text{ }\int_{\Delta}f\left(  \boldsymbol{r}\right)
\Phi\left(  \boldsymbol{r}\right)  ^{\dagger}\Phi\left(  \boldsymbol{r}%
\right)  d\boldsymbol{r}.
\end{equation}
If $Y\geq0$ i.e. $Y=A^{2}$ for some $A\in\mathcal{B}_{2sa}\left(
\mathcal{H}^{N}\right)  ,$ this leads one to consider the subspace,
$\mathcal{S}_{2\Delta^{c}},$ of $\mathcal{B}_{2sa}\left(  \mathcal{H}%
^{N}\right)  $ defined by
\begin{subequations}
\begin{equation}
\mathcal{S}_{2\Delta^{c}}=\text{ linear span}\left\{  A\left\vert
\widehat{\Phi\left(  \boldsymbol{r}\right)  ^{\dagger}\Phi\left(
\boldsymbol{r}\right)  }\left\vert A\right)  =0\forall\boldsymbol{r}\in
\Delta^{c},A\in\mathcal{B}_{2sa}\left(  \mathcal{H}^{N}\right)  \right.
\right\}  ,
\end{equation}
that by the preceeding lemma corresponds to
\end{subequations}
\begin{align}
\mathcal{S}_{2\Delta^{c}}  &  =\text{ linear span}\left\{  A\left\vert
\operatorname*{Tr}\left\{  \Phi\left(  \boldsymbol{r}\right)  ^{\dagger}%
\Phi\left(  \boldsymbol{r}\right)  A^{2}\right\}  =0\forall\boldsymbol{r}%
\in\Delta^{c},A\in\mathcal{B}_{2sa}\left(  \mathcal{H}^{N}\right)  \right.
\right\} \nonumber\\
&  =\text{ linear span}\left\{  A\left\vert \Phi\left(  \boldsymbol{r}\right)
^{\dagger}\Phi\left(  \boldsymbol{r}\right)  A=0\forall\boldsymbol{r}\in
\Delta^{c},A\in\mathcal{B}_{2sa}\left(  \mathcal{H}^{N}\right)  \right.
\right\}  .
\end{align}
One should note that $Y\in\mathcal{S}_{1\Delta^{c}}$ or $A\in\mathcal{S}%
_{2\Delta^{c}}$ does not preclude that
\begin{equation}
\operatorname*{Tr}\left\{  \Phi\left(  \boldsymbol{r}\right)  ^{\dagger}%
\Phi\left(  \boldsymbol{r}\right)  Y\right\}  =0,\text{for some }%
\boldsymbol{r}\in\Delta\text{ and }Y\in\mathcal{S}_{1\Delta^{c}},
\end{equation}
nor
\begin{equation}
\widehat{\Phi\left(  \boldsymbol{r}\right)  ^{\dagger}\Phi\left(
\boldsymbol{r}\right)  }\left\vert A\right)  =0,\text{ for some }%
\boldsymbol{r}\in\Delta\text{ and }A\in\mathcal{S}_{2\Delta^{c}}.
\end{equation}


As $\mathcal{B}_{2sa}\left(  \mathcal{H}^{N}\right)  $ is a Hilbert space one
can define the orthogonal compliment, $\left(  \mathcal{S}_{2\Delta^{c}%
}\right)  ^{\bot},$ as
\begin{equation}
\left(  \mathcal{S}_{2\Delta^{c}}\right)  ^{\bot}=\left\{  B\left\vert \left(
B\left\vert A\right.  \right)  =0\forall A\in\mathcal{S}_{2\Delta^{c}}\right.
\right\}  ,\forall A\in\mathcal{S}_{2\Delta^{c}},B\in\mathcal{B}_{2sa}\left(
\mathcal{H}^{N}\right) \nonumber
\end{equation}
and thus
\begin{equation}
\mathcal{B}_{2sa}\left(  \mathcal{H}^{N}\right)  =\mathcal{S}_{2\Delta^{c}%
}\oplus\left(  \mathcal{S}_{2\Delta^{c}}\right)  ^{\bot}.
\end{equation}
This is equivalent to
\begin{equation}
\int_{\Delta}\left(  \sum_{1\leq j\leq N}f\left(  \boldsymbol{r}_{j}\right)
\right)  ^{2}d\boldsymbol{r}_{1}\boldsymbol{\cdots}d\boldsymbol{r}_{N}%
<\infty;1\leq N\leq\infty.
\end{equation}


\section{Regular Points \label{Regular}}

\subsection{Derivatives \label{DerivDef}}

\bigskip A Banach space mapping, $F,$between spaces $\mathcal{B}_{1}$ and
$\mathcal{B}_{2}$ i.e.
\begin{equation}
F:\mathcal{B}_{1}\rightarrow\mathcal{B}_{2},
\end{equation}
is G\^{a}teaux differentiable at a point $A\in\mathcal{B}_{1}$ contained in an
open subset $U$ of $\mathcal{B}_{1}$ if
\begin{equation}
\nabla_{B}F\left(  A\right)  =\lim_{t\rightarrow0}\frac{1}{t}\left\{  F\left(
A+tB\right)  -F\left(  A\right)  \right\}  =\left.  \frac{d}{dt}F\left(
A+tB\right)  \right\vert _{t=0},
\end{equation}
exists for any $B\in\mathcal{B}_{1}.$ If $\nabla F\left(  A\right)  $ exists
for all $B\in\mathcal{B}_{1}$ it is called the G\^{a}teaux derivative of $F$
at $A$ and $\nabla_{B}F\left(  A\right)  $ $=\left[  \nabla F\left(  A\right)
\right]  \left(  B\right)  $ the directional G\^{a}teaux derivative in the
direction $B$ at the point $A.$ The map $F$ is called Fr\'{e}chet
differentiable at a point $A\in\mathcal{B}_{1}$ if there exists a linear
operator
\begin{equation}
dF\left(  A\right)  :\mathcal{B}_{1}\rightarrow\mathcal{B}_{2}%
\end{equation}
such that
\begin{equation}
\lim_{\left\Vert B\right\Vert \rightarrow0}\frac{1}{\left\Vert B\right\Vert
}\left\{  F\left(  A+B\right)  -F\left(  A\right)  -\left[  dF\left(
A\right)  \right]  \left(  B\right)  \right\}  =0,
\end{equation}
$dF\left(  A\right)  $ is called the Fr\'{e}chet derivative of $F$ at $A.$ $F$
is Fr\'{e}chet differentiable in an open subset $U$ of $\mathcal{B}_{1}$if and
only if it is

\begin{itemize}
\item G\^{a}teaux differentiable i.e. $\nabla F\left(  A\right)  $ exists

\item $\left[  \nabla F\left(  A\right)  \right]  \left(  B\right)
\equiv\nabla_{B}F\left(  A\right)  $ is linear in $B$

\item $\sup_{B\neq0}\left\{  \frac{\left\Vert \nabla_{B}F\left(  A\right)
\right\Vert }{\left\Vert B\right\Vert }\right\}  <\infty$
\end{itemize}

then
\begin{equation}
\nabla_{B}F\left(  A\right)  =\left[  dF\left(  A\right)  \right]  \left(
B\right)  =\left[  \nabla F\left(  A\right)  \right]  \left(  B\right)  .
\end{equation}
Fr\'{e}chet differentiation defines a map (generally non linear) from the
space of functions, $\mathcal{F}\left(  \mathcal{B}_{1},\mathcal{B}%
_{2}\right)  ,$ mapping $\mathcal{B}_{1}$ to $\mathcal{B}_{2}$
\begin{align}
d  &  :\mathcal{F}\left(  \mathcal{B}_{1},\mathcal{B}_{2}\right)
\rightarrow\mathcal{F}\left(  \mathcal{B}_{1},\mathcal{B}_{2}\right)
\nonumber\\
d  &  :F\rightarrow dF.
\end{align}
that produces the Fr\'{e}chet derivative, the Fr\'{e}chet derivative at a
point $A$ is $dF\left(  A\right)  \in\mathcal{L}\left(  \mathcal{B}%
_{1},\mathcal{B}_{2}\right)  ,$ the bounded linear maps from $\mathcal{B}_{1}$
to $\mathcal{B}_{2}.$

\subsection{Tangent and Normal Spaces \label{TanNorm}}

\subsubsection{Differentiability of the Constraint}

The constraint functional is a quadratic map%
\begin{equation}
\Xi_{2}:\mathcal{B}_{2sa}\left(  \mathcal{H}^{N}\right)  \rightarrow
\mathfrak{D}\mathcal{\subset}L_{1+}\left(  \mathbb{R}^{3}\right)  \cap
L_{3+}\left(  \mathbb{R}^{3}\right)  \subset L_{1}\left(  \mathbb{R}%
^{3}\right)  ,
\end{equation}
where%
\begin{equation}
\Xi_{2}\left(  A\right)  =\left(  A\left\vert \widehat{\Phi^{\dagger}\Phi
}_{\delta}A\right.  \right)  =\rho.
\end{equation}
The super operator valued distribution $\widehat{\Phi^{\dagger}\Phi}_{\delta}$
is defined by
\begin{align}
\widehat{\Phi^{\dagger}\Phi}_{\delta}  &  :\mathbb{R}^{3}\cong\mathcal{D}%
_{\delta}\subset\mathcal{S}\left(  \mathbb{R}^{3}\right)  ^{d}\rightarrow
\mathcal{B}_{sa}\left(  \mathcal{B}_{2sa}\left(  \mathcal{H}^{N}\right)
\right) \nonumber\\
\widehat{\Phi^{\dagger}\Phi}_{\delta}\left(  \boldsymbol{r}\right)   &
\equiv\widehat{\Phi^{\dagger}\Phi}\left(  \delta_{\mathbf{r}}\right)
\in\mathcal{B}_{sa}\left(  \mathcal{B}_{2sa}\left(  \mathcal{H}^{N}\right)
\right)  ,
\end{align}
(where we have used the identification of $\mathbb{R}^{3}$ with the set of
Dirac delta functions $\mathcal{D}_{\delta},$ $\mathcal{S}\left(
\mathbb{R}^{3}\right)  $ is the space of Schwartz test functions also called
the space of test functions of rapid decrease, and $\mathcal{S}\left(
\mathbb{R}^{3}\right)  ^{d}$ is the space of Tempered distributions), in terms
of the operator valued distribution $\Phi^{\dagger}\Phi_{\delta}$
\begin{align}
\Phi^{\dagger}\Phi_{\delta}  &  :\mathbb{R}^{3}\rightarrow\mathcal{B}%
_{2sa}\left(  \mathcal{H}^{N}\right) \nonumber\\
\Phi^{\dagger}\Phi_{\delta}\left(  \boldsymbol{r}\right)   &  \equiv
\Phi^{\dagger}\left(  \delta_{\mathbf{r}}\right)  \Phi\left(  \delta
_{\mathbf{r}}\right)  \equiv\Phi^{\dagger}\left(  \boldsymbol{r}\right)
\Phi\left(  \boldsymbol{r}\right)  \in\mathcal{B}_{2sa}\left(  \mathcal{H}%
^{N}\right)
\end{align}
as
\begin{equation}
\widehat{\Phi^{\dagger}\Phi}_{\delta}\left(  \boldsymbol{r}\right)  \left\vert
A\right)  =\left\vert \left[  \Phi\left(  \boldsymbol{r}\right)  ^{\dagger
}\Phi\left(  \boldsymbol{r}\right)  A\right]  _{+}\right)  .
\end{equation}
The density $\rho_{A^{2}}$ associated with an operator $A\in\mathcal{B}%
_{2sa}\left(  \mathcal{H}^{N}\right)  $ is given by
\begin{equation}
\left(  \left.  \left[  \Phi\left(  \boldsymbol{r}\right)  ^{\dagger}%
\Phi\left(  \boldsymbol{r}\right)  A\right]  _{+}\right\vert A\right)
=\operatorname{Tr}\left\{  \Phi\left(  \boldsymbol{r}\right)  ^{\dagger}%
\Phi\left(  \boldsymbol{r}\right)  A^{2}\right\}  \text{ }\forall
\boldsymbol{r\in}\mathbb{R}^{3}.
\end{equation}
One does not distinguish between densities that are only different on sets of
measure zero i.e.
\begin{equation}
\rho\sim\rho^{\prime}\text{ if }\int\left\vert \rho\left(  \boldsymbol{r}%
\right)  -\rho^{\prime}\left(  \boldsymbol{r}\right)  \right\vert
d\boldsymbol{r}=0\text{ i.e. }\rho\overset{a.e}{=}\rho^{\prime}.
\end{equation}
It is easy to see that $\Xi_{2}$ is Fr\'{e}chet differentiable as the
directional G\^{a}teaux derivative, $\nabla_{B}\Xi_{2}\left(  A\right)  ,$ of
$\Xi_{2}$ exists at $A$ for all $B$
\begin{align}
&  \lim_{t\rightarrow0}\frac{1}{t}\left\{  \Xi_{2}\left(  A+tB\right)
-\Xi_{2}\left(  A\right)  \right\}  =\lim_{t\rightarrow0}\frac{1}{t}\left\{
\left(  \left(  A+tB\right)  \left\vert \widehat{\Phi^{\dagger}\Phi_{\delta}%
}\left(  A+tB\right)  \right.  \right)  -\left(  A\left\vert \widehat{\Phi
^{\dagger}\Phi_{\delta}}A\right.  \right)  \right\} \nonumber\\
&  \left.  \qquad\qquad=\right.  \lim_{t\rightarrow0}\left\{  2\left(
A\left\vert \widehat{\Phi^{\dagger}\Phi_{\delta}}B\right.  \right)  -t\left(
B\left\vert \widehat{\Phi^{\dagger}\Phi_{\delta}}B\right.  \right)  \right\}
=2\left(  A\left\vert \widehat{\Phi^{\dagger}\Phi_{\delta}}B\right.  \right)
\end{align}
leading to a G\^{a}teaux derivative
\begin{equation}
\nabla\Xi_{2}\left(  A\right)  =2\left(  \left[  \Phi^{\dagger}\Phi_{\delta
}\left(  A\right)  \right]  _{+}\right\vert ,
\end{equation}
which is clearly a bounded linear functional
\begin{equation}
\nabla\Xi_{2}\left(  A\right)  :\mathcal{B}_{2sa}\left(  \mathcal{H}%
^{N}\right)  \rightarrow L_{1}\left(  \mathbb{R}^{3}\right)  .
\end{equation}
The three conditions for Fr\'{e}chet differentialability are satisfied and the
Fr\'{e}chet derivative of $\Xi_{2}$ at $A$ is then also
\begin{equation}
d\Xi_{2}\left(  A\right)  =\nabla\Xi_{2}\left(  A\right)  =2\left(  \left[
\Phi^{\dagger}\Phi_{\delta}\left(  A\right)  \right]  _{+}\right\vert
\end{equation}
and
\begin{equation}
d\Xi_{2}\left(  A\right)  :\mathcal{B}_{2sa}\left(  \mathcal{H}^{N}\right)
\rightarrow L_{1}\left(  \mathbb{R}^{3}\right)  .
\end{equation}


The differential map
\begin{equation}
d\Xi_{2}:\mathcal{B}_{2sa}\left(  \mathcal{H}^{N}\right)  \rightarrow
\mathcal{L}\left(  \mathcal{B}_{2sa}\left(  \mathcal{H}^{N}\right)
,L_{1}\left(  \mathbb{R}^{3}\right)  \right)
\end{equation}
is linear (as $\Xi_{2}$ is quadratic) and one can consider
\begin{equation}
\ker d\Xi_{2}=\left\{  \left.  A\right\vert d\Xi_{2}\left(  A\right)
=0\right\}
\end{equation}
and
\begin{equation}
\operatorname{range}d\Xi_{2}=\left\{  \left.  B\right\vert B=d\Xi_{2}\left(
A\right)  \right\}  .
\end{equation}
$A\in\ker d\Xi_{2}$ if
\begin{equation}
\left(  \left[  \Phi^{\dagger}\left(  \boldsymbol{r}\right)  \Phi\left(
\boldsymbol{r}\right)  A\right]  _{+}\right\vert =0\,\forall\boldsymbol{r\in
}\mathbb{R}^{3}\Longleftrightarrow A=0,
\end{equation}
so is empty. However $\ker d\Xi_{2}\left(  A\right)  $ is not empty as
\begin{align}
\left(  A\left\vert \widehat{\Phi^{\dagger}\left(  \boldsymbol{r}\right)
\Phi\left(  \boldsymbol{r}\right)  }B\right.  \right)   &  =\left(
A\left\vert \left[  \Phi^{\dagger}\left(  \boldsymbol{r}\right)  \Phi\left(
\boldsymbol{r}\right)  B\right]  _{+}\right.  \right)  =\operatorname{Tr}%
\left\{  A\left[  \Phi\left(  \boldsymbol{r}\right)  ^{\dagger}\Phi\left(
\boldsymbol{r}\right)  B\right]  _{+}\right\} \nonumber\\
&  =\operatorname{Tr}\left\{  A\Phi\left(  \boldsymbol{r}\right)  ^{\dagger
}\Phi\left(  \boldsymbol{r}\right)  B\right\}  =0\quad\forall\boldsymbol{r}%
\in\mathbb{R}^{3}%
\end{align}
implies that $A$ is localized on $\boldsymbol{\Delta}_{A}$ and $B$ on
$\boldsymbol{\Delta}_{B}$ and
\begin{equation}
\boldsymbol{\Delta}_{A}\cap\boldsymbol{\Delta}_{B}=\phi,
\end{equation}
where
\begin{equation}
\boldsymbol{\Delta}_{X}=\operatorname*{supp}\eta_{X}=\overline{\left\{
\left.  \mathbf{r}\right\vert \eta_{X}\left(  \mathbf{r}\right)
\neq0\right\}  },
\end{equation}
i.e. the closure of the set of points where $\eta_{X}$ is nonzero. The set
$\boldsymbol{\Delta}_{X}$ might not be bounded thus not compact. If
$\boldsymbol{\Delta}_{A}\cap\boldsymbol{\Delta}_{B}=\phi$ then $A\bot B$ i.e.
\begin{equation}
\left(  \left.  A\right\vert B\right)  =0.
\end{equation}


\begin{theorem}
The operators $\left\{  \left.  \left\vert \left[  \Phi\left(  \boldsymbol{r}%
\right)  ^{\dagger}\Phi\left(  \boldsymbol{r}\right)  A\right]  _{+}\right)
\right\vert \mathbf{r\in}\mathbb{R},A\in\mathcal{B}_{2sa}\left(
\mathcal{H}^{N}\right)  \right\}  $ are linearly independent .

\begin{proof}
This is true if the only solution to
\begin{align}
&  \int_{\boldsymbol{\Delta}_{A}}f\left(  \boldsymbol{r}\right)  \left\vert
\left[  \Phi\left(  \boldsymbol{r}\right)  ^{\dagger}\Phi\left(
\boldsymbol{r}\right)  A\right]  _{+}\right)  d\boldsymbol{r}\nonumber\\
&  =\frac{1}{2}\int_{\boldsymbol{\Delta}_{A}}f\left(  \boldsymbol{r}\right)
\left\vert A\Phi\left(  \boldsymbol{r}\right)  ^{\dagger}\Phi\left(
\boldsymbol{r}\right)  +\Phi\left(  \boldsymbol{r}\right)  ^{\dagger}%
\Phi\left(  \boldsymbol{r}\right)  A\right)  d\boldsymbol{r}=0;\,f\left(
\boldsymbol{r}\right)  \in\mathbb{R} \label{Th1}%
\end{align}
is $f\left(  \boldsymbol{r}\right)  \overset{a.e.}{=}0.$ Eq. $\left(
\ref{Th1}\right)  $ is equivalent to
\begin{equation}
\int_{\boldsymbol{\Delta}_{A}}f\left(  \boldsymbol{r}\right)
\operatorname{Tr}\left\{  YA\Phi\left(  \boldsymbol{r}\right)  ^{\dagger}%
\Phi\left(  \boldsymbol{r}\right)  \right\}  d\boldsymbol{r}=0;\,\forall
Y\in\mathcal{B}_{2sa}\left(  \mathcal{H}^{N}\right)  ,
\end{equation}
thus as%
\begin{equation}
\int_{\boldsymbol{\Delta}_{A}}f\left(  \boldsymbol{r}\right)
\operatorname{Tr}\left\{  YA\Phi\left(  \boldsymbol{r}\right)  ^{\dagger}%
\Phi\left(  \boldsymbol{r}\right)  \right\}  d\boldsymbol{r}=0;\,\forall
Y\in\mathcal{B}_{2}\left(  \mathcal{H}^{N}\right)  , \label{Theo1}%
\end{equation}
where
\begin{equation}
Y=Y_{1}+iY_{2};\,Y_{1},Y_{2}\in\mathcal{B}_{2sa}\left(  \mathcal{H}%
^{N}\right)  .
\end{equation}
Eq. $\left(  \ref{Theo1}\right)  $ can be recast as%
\begin{equation}
\operatorname{Tr}\left\{  YA\mathfrak{f}\right\}  =\left\langle Y\left\vert
A\mathfrak{f}\right.  \right\rangle _{\mathcal{B}_{2}\left(  \mathcal{H}%
^{N}\right)  }=0;\,\forall Y\in\mathcal{B}_{2}\left(  \mathcal{H}^{N}\right)
, \label{p1}%
\end{equation}
where we constrain $\mathfrak{f}$ to be a Hilbert-Schimdt operator i.e.
\begin{equation}
\mathfrak{f}=\int_{\boldsymbol{\Delta}_{A}}f\left(  \boldsymbol{r}\right)
\Phi\left(  \boldsymbol{r}\right)  ^{\dagger}\Phi\left(  \boldsymbol{r}%
\right)  d\boldsymbol{r\,}\in\mathcal{B}_{2sa}\left(  \mathcal{H}\right)  ,
\end{equation}
so that%
\begin{equation}
\left.  \mathfrak{f}\right\vert _{\mathcal{H}^{N}}\in\mathcal{B}_{2sa}\left(
\mathcal{H}^{N}\right)  ,1\leq N\leq\infty.
\end{equation}
Eq. $\left(  \ref{p1}\right)  $ shows that
\begin{equation}
A\mathfrak{f}=\mathfrak{f}A=0.
\end{equation}
The operator $\mathfrak{f}$ can expanded in an alternate form by considering
its restrictions to $\mathcal{H}^{N}$ for $0\leq N\leq\infty,$ as
\begin{equation}
\mathfrak{f}=\sum_{1\leq N\leq\infty}\binom{\boldsymbol{r}_{N}}{N}^{-1}%
\int_{\boldsymbol{\Delta}_{A}}\sum_{1\leq j\leq N}f\left(  \boldsymbol{r}%
_{j}\right)  \left\vert \boldsymbol{r}_{1}\cdots\boldsymbol{r}_{N}%
\right\rangle \left\langle \boldsymbol{r}_{1}\cdots\boldsymbol{r}%
_{N}\right\vert d\boldsymbol{r}_{1}\cdots d\boldsymbol{r}_{N}.
\end{equation}
As $A$ is campact and self adjoint it has the expansion
\begin{equation}
A=\sum_{1\leq j\leq r_{N}}\alpha_{j}\left\vert \Psi_{j}\right\rangle
\left\langle \Psi_{j}\right\vert ;\,\alpha_{j}\in\mathbb{R},\left\vert
\Psi_{j}\right\rangle \in\mathcal{H}^{N},1\leq j\leq r_{N}.
\end{equation}
thus one can express $\mathfrak{f}A=0$ as
\begin{align}
&  \sum_{1\leq j\leq N}\int_{\boldsymbol{\Delta}_{A}}f\left(  \boldsymbol{r}%
_{j}\right)  \left\vert \boldsymbol{r}_{1}\cdots\boldsymbol{r}_{N}%
\right\rangle \left\langle \boldsymbol{r}_{1}\cdots\boldsymbol{r}%
_{N}\right\vert d\boldsymbol{r}_{1}\cdots d\boldsymbol{r}_{N}\sum_{1\leq j\leq
r_{N}}\alpha_{j}\left\vert \Psi_{j}\right\rangle \left\langle \Psi
_{j}\right\vert \nonumber\\
&  =\sum_{1\leq j\leq r_{N}}\alpha_{j}\sum_{1\leq j\leq N}\int%
_{\boldsymbol{\Delta}_{A}}f\left(  \boldsymbol{r}_{j}\right)  \left\langle
\boldsymbol{r}_{1}\cdots\boldsymbol{r}_{N}\left\vert \Psi_{j}\right.
\right\rangle \left\vert \boldsymbol{r}_{1}\cdots\boldsymbol{r}_{N}%
\right\rangle \left\langle \Psi_{j}\right\vert d\boldsymbol{r}_{1}\cdots
d\boldsymbol{r}_{N}\nonumber\\
&  =\sum_{1\leq k\leq r_{N}}\alpha_{k}\sum_{1\leq j\leq N}\int%
_{\boldsymbol{\Delta}_{A}}f\left(  \boldsymbol{r}_{j}\right)  \Psi_{k}\left(
\boldsymbol{r}_{1}\cdots\boldsymbol{r}_{N}\right)  \left\vert \boldsymbol{r}%
_{1}\cdots\boldsymbol{r}_{N}\right\rangle \left\langle \Psi_{k}\right\vert
d\boldsymbol{r}_{1}\cdots d\boldsymbol{r}_{N}=0
\end{align}
giving
\begin{equation}
\alpha_{k}\left\{  \sum_{1\leq j\leq N}f\left(  \boldsymbol{r}_{j}\right)
\right\}  \Psi_{k}\left(  \boldsymbol{r}_{1}\cdots\boldsymbol{r}_{N}\right)
=0\,\forall k\text{ and almost everywhere for }\boldsymbol{r}_{j}%
\in\boldsymbol{\Delta}_{A}.
\end{equation}
If
\begin{equation}
\sum_{1\leq j\leq N}f\left(  \boldsymbol{r}_{j}\right)  \neq0\text{ almost
everywhere for }\boldsymbol{r}_{j}\in\boldsymbol{\Delta}_{A},\,1\leq j\leq N
\end{equation}
then
\begin{equation}
\alpha_{k}\Psi_{k}\left(  \boldsymbol{r}_{1}\cdots\boldsymbol{r}_{N}\right)
=0\text{ }\forall k\text{ and almost everywhere for }\boldsymbol{r}_{j}%
\in\boldsymbol{\Delta}_{A},\,1\leq j\leq N.\text{ }%
\end{equation}
If further $\alpha_{k}\neq0$ then
\begin{equation}
\Psi_{k}\left(  \boldsymbol{r}_{1}\cdots\boldsymbol{r}_{N}\right)
\overset{a.e.}{=}0,
\end{equation}
which cannot be true as
\begin{align}
\Phi\left(  \boldsymbol{r}\right)  ^{\dagger}\Phi\left(  \boldsymbol{r}%
\right)  \left\vert \Psi_{k}\right\rangle  &  =0\forall\boldsymbol{r}%
\in\boldsymbol{\Delta}_{A}\nonumber\\
&  \Longleftrightarrow\Phi\left(  \boldsymbol{r}\right)  ^{\dagger}\Phi\left(
\boldsymbol{r}\right)  \Psi_{k}\left(  \boldsymbol{r}_{1}\cdots\boldsymbol{r}%
_{N}\right)  =0\forall\boldsymbol{r}\in\boldsymbol{\Delta}_{A}\nonumber\\
&  \Longleftrightarrow\left\langle \Psi_{k}\left\vert \Phi\left(
\boldsymbol{r}\right)  ^{\dagger}\Phi\left(  \boldsymbol{r}\right)  \right.
\Psi_{k}\right\rangle =0\forall\boldsymbol{r}\in\boldsymbol{\Delta}_{A}.
\end{align}
Thus%
\begin{equation}
\sum_{1\leq j\leq N}f\left(  \boldsymbol{r}_{j}\right)  \overset{a.e.}{=}%
0;\boldsymbol{r}_{j}\in\boldsymbol{\Delta}_{A},\,1\leq j\leq
N\Longleftrightarrow f\left(  \boldsymbol{r}\right)  \overset{a.e.}{=}0.
\end{equation}

\end{proof}
\end{theorem}

As $\left\{  \left.  \left\vert \left[  \Phi\left(  \boldsymbol{r}\right)
^{\dagger}\Phi\left(  \boldsymbol{r}\right)  A\right]  _{+}\right)
\right\vert \mathbf{r\in}\mathbb{R},A\in\mathcal{B}_{2sa}\left(
\mathcal{H}^{N}\right)  \right\}  $ linearly independent they form generalized
basis of $\mathcal{D}_{A}\subset\mathcal{B}_{2sa}\left(  \mathcal{H}%
^{N}\right)  $, where $\mathcal{D}_{A}\cong$ $L_{1}\left(  \boldsymbol{\Delta
}_{A}\right)  =$ $\operatorname{range}\left\{  d\Xi_{2}\left(  A\right)
\right\}  .$ (Not a Hilbert space isomorpism)

\subsubsection{Tangent and Normal Spaces of the Constraint}

The map $\Xi_{2}$ defines a differentiable manifold in $\mathcal{B}%
_{2sa}\left(  \mathcal{H}^{N}\right)  $ given by
\begin{equation}
\mathfrak{M}=\left\{  \left.  A\right\vert \Xi_{2}\left(  A\right)
=\rho\right\}  .
\end{equation}
A tangent vector $\delta A$ at the point $A$ is defined by curves $x\left(
t\right)  \subset\mathfrak{M,}t\in\left[  -\delta,\delta\right]
\subset\mathbb{R},$ such that
\begin{align*}
\Xi_{2}\left(  x\left(  t\right)  \right)   &  =\rho\,,t\in\left[
-\delta,\delta\right] \\
x\left(  0\right)   &  =A\\
x  &  :\left[  -\delta,\delta\right]  \longrightarrow\mathfrak{M}%
\end{align*}
and is defined to be the equivalence class, $\left[  \frac{dx\left(  0\right)
}{dt}\right]  ,$ of velocity vectors to such curves at the point $A$ that have
$\frac{dx\left(  0\right)  }{dt}$ as a representative. One can show that any
tangent vector $\delta A$ (which equals $\frac{dx\left(  0\right)  }{dt}$ for
some curve $x)$ belongs to $\ker\left\{  d\Xi_{2}\left(  A\right)  \right\}  $
i.e.%
\begin{equation}
\left(  \left.  \left[  \Phi\left(  \boldsymbol{r}\right)  ^{\dagger}%
\Phi\left(  \boldsymbol{r}\right)  A\right]  _{+}\right\vert \delta A\right)
=0\forall\boldsymbol{r\in}\mathbb{R}^{3},
\end{equation}
so that
\begin{equation}
\mathcal{T}_{A}\subseteq\operatorname{Ker}\left\{  d\Xi_{2}\left(  A\right)
\right\}  ,
\end{equation}
where the tangent space $\mathcal{T}_{A}$ $\subset\mathcal{B}_{2sa}\left(
\mathcal{H}^{N}\right)  $ at the point $A$ is defined as
\begin{equation}
\mathcal{T}_{A}=\left\{  \delta A\left\vert \text{there exists a }x\left(
t\right)  \text{ such that }x\left(  0\right)  =A\text{ and }\delta
A=\frac{dx\left(  0\right)  }{dt}\right.  \right\}  .
\end{equation}
However one should note that $\delta A\in\ker\left\{  d\Xi_{2}\left(
A\right)  \right\}  $ does not imply that there exists a $\frac{dx\left(
0\right)  }{dt}$ such that $\delta A=\frac{dx\left(  0\right)  }{dt}$ i.e.
this condition is necessary but not sufficient for $\delta A$ to be a tangent vector.

The normal space, $\mathcal{N}_{A}\subset\mathcal{B}_{2sa}\left(
\mathcal{H}^{N}\right)  ,$ at the point $A$ is defined as
\begin{equation}
\mathcal{N}_{A}=\left\{  B\left\vert \left(  \left.  \delta A\right\vert
B\right)  =0\,\forall\boldsymbol{r\in}\mathbb{R}^{3},\forall\delta
A\in\mathcal{T}_{A}\right.  \right\}  =\left(  \mathcal{T}_{A}\right)  ^{\bot
},
\end{equation}
i.e. the orthogonal complement to the tangent space $\mathcal{T}_{A}.$ Clearly%
\begin{equation}
\mathcal{N}_{A}=\left(  \mathcal{T}_{A}\right)  ^{\bot}\supseteq\ker\left\{
d\Xi_{2}\left(  A\right)  \right\}  ^{\bot}\cong\operatorname{range}\left\{
d\Xi_{2}\left(  A\right)  \right\}  .
\end{equation}


One can show that $\left\{  \left.  \left[  d\Xi_{2}\left[  \boldsymbol{r}%
\right]  \left(  A\right)  \right]  \left(  \cdot\right)  \equiv\left(
\left[  \Phi\left(  \boldsymbol{r}\right)  ^{\dagger}\Phi\left(
\boldsymbol{r}\right)  A\right]  _{+}\right\vert \equiv\operatorname{Tr}%
\left\{  \left[  \Phi\left(  \boldsymbol{r}\right)  ^{\dagger}\Phi\left(
\boldsymbol{r}\right)  A\right]  _{+}\cdot\right\}  \right\vert
\boldsymbol{r\in}\mathbb{R}^{3}\right\}  $ is a linearly independent set that
forms a generalized basis for the space $\ker\left\{  d\Xi_{2}\left(
A\right)  \right\}  ^{\bot}$, so that $A$ is a \emph{normal, }hence a\emph{
regular, }\cite{Hestenes}\emph{ }point of the constraint (see Appendix basis
\ref{Genbasis}). As $d\Xi_{2}\left(  A\right)  $ is $1-1$ for the space
$\ker\left\{  d\Xi_{2}\left(  A\right)  \right\}  ^{\bot}\cong%
\operatorname{range}\left\{  d\Xi_{2}\left(  A\right)  \right\}  $ and
\begin{equation}
\left[  d\Xi_{2}\left(  A\right)  \right]  :\mathcal{B}_{2sa}\left(
\mathcal{H}^{N}\right)  \rightarrow L_{1}\left(  \boldsymbol{\Delta}%
_{A}\right)
\end{equation}
is thus surjective, . *****

\section{Complexification of $\mathcal{B}_{2sa}\left(  \mathcal{H}^{N}\right)
$ \label{Complex}}

The complex Hilbert space $\mathcal{B}_{2}\left(  \mathcal{H}^{N}\right)  $
can be considered as the complexification of the real Hilbert space
$\mathcal{B}_{2sa}\left(  \mathcal{H}^{N}\right)  $ i.e.
\begin{equation}
\mathcal{B}_{2}\left(  \mathcal{H}^{N}\right)  =\mathcal{B}_{2sa}\left(
\mathcal{H}^{N}\right)  \oplus i\mathcal{B}_{2sa}\left(  \mathcal{H}%
^{N}\right)  ,
\end{equation}
such that $Z\in\mathcal{B}_{2}\left(  \mathcal{H}^{N}\right)  ,$ where
\begin{equation}
Z=X+iY;\text{ }X,Y\in\mathcal{B}_{2sa}\left(  \mathcal{H}^{N}\right)
\end{equation}
The inner product in $\mathcal{B}_{2}\left(  \mathcal{H}^{N}\right)  $ is
defined in terms of the inner product in $\mathcal{B}_{2sa}\left(
\mathcal{H}^{N}\right)  $ as
\begin{equation}
\left(  Z_{1}\left\vert Z_{2}\right.  \right)  _{HS}=\left(  X_{1}\left\vert
X_{2}\right.  \right)  +\left(  Y_{1}\left\vert Y_{2}\right.  \right)
+i\left\{  \left(  X_{1}\left\vert Y_{2}\right.  \right)  -\left(
Y_{1}\left\vert X_{2}\right.  \right)  \right\}  ,
\end{equation}
which corresponds to
\begin{equation}
\left(  Z_{1}\left\vert Z_{2}\right.  \right)  _{HS}=\operatorname{Tr}\left\{
Z_{1}^{\dagger}Z_{2}\right\}  ,
\end{equation}
as%
\begin{equation}
\operatorname{Tr}\left\{  \left(  X_{1}-iY_{1}\right)  \left(  X_{2}%
+iY_{2}\right)  \right\}  =\operatorname{Tr}\left\{  X_{1}X_{2}-iY_{1}%
X_{2}+iX_{1}Y_{2}+Y_{1}Y_{2}\right\}  .
\end{equation}


Note that $\mathcal{B}_{2sa}\left(  \mathcal{H}^{N}\right)  $ (self adjoint
operators) and $i\mathcal{B}_{2sa}\left(  \mathcal{H}^{N}\right)  $ (anti self
adjoint operators) are not orthogonal subspaces (they are not even subspaces
under $\mathbb{C}$-linear combination).

\section{State Normalization \label{Normalization}}

If $D^{N}\in\mathcal{B}_{2}\left(  \mathcal{H}^{N}\right)  $ then
\begin{equation}
\operatorname{Tr}\left\{  \mathcal{N}D^{N}\right\}  =N\operatorname{Tr}%
\left\{  D^{N}\right\}  ,
\end{equation}
where
\begin{equation}
\mathcal{N}=\int\Phi\left(  \boldsymbol{r}\right)  ^{\dagger}\Phi\left(
\boldsymbol{r}\right)  d\boldsymbol{r,}%
\end{equation}
as
\begin{equation}
\mathcal{N}\left\vert \Psi\right\rangle =N\left\vert \Psi\right\rangle
\quad\forall\left\vert \Psi\right\rangle \in\mathcal{H}^{N}.
\end{equation}
For any state $D$
\begin{equation}
\operatorname{Tr}\left\{  \mathcal{N}D\right\}  =\int\rho\left(
\boldsymbol{r,\xi}\right)  d\boldsymbol{r}d\boldsymbol{\xi,}%
\end{equation}
thus
\begin{equation}
N\operatorname{Tr}\left\{  D^{N}\right\}  =N\Rightarrow\operatorname{Tr}%
\left\{  D^{N}\right\}  =1.
\end{equation}


\section{Does $\rho_{1}=\rho_{2}$ pointwise or in the mean}

When do we consider two densities to be equal i.e.
\begin{equation}
\rho_{1}=\rho_{2}%
\end{equation}
is it when
\begin{equation}
\rho_{1}\left(  \mathbf{r}\right)  =\rho_{2}\left(  \mathbf{r}\right)  \text{
}\forall\mathbf{r\in}\mathbb{R}^{3}%
\end{equation}
or
\begin{equation}
\int\left\vert \rho_{1}\left(  \mathbf{r}\right)  -\rho_{2}\left(
\mathbf{r}\right)  \right\vert d\mathbf{r}=0\text{ i.e. }\rho_{1}%
\overset{a.e}{=}\rho_{2}?
\end{equation}
Two FORDO's are equal when
\begin{align}
\left(  A\left\vert D^{1}\right.  \right)   &  =\left(  A\left\vert
D^{2}\right.  \right)  \text{ }\forall\mathbf{A\in}\mathcal{B}_{sa}\left(
\mathcal{H}^{1}\right)  \Longleftrightarrow\\
\left(  A\left\vert D_{1}^{1}-D_{2}^{1}\right.  \right)   &  =0\text{ }%
\forall\mathbf{A\in}\mathcal{B}_{sa}\left(  \mathcal{H}^{1}\right)
;\,D_{1}^{1}-D_{2}^{1}\in\mathcal{B}_{1sa}\left(  \mathcal{H}^{1}\right)  ,
\end{align}
thus $D_{1}^{1}-D_{2}^{1}=0$ as an operator in $\mathcal{B}_{1sa}\left(
\mathcal{H}^{1}\right)  $ i.e.
\begin{equation}
\int\left\vert D_{1}^{1}\left(  \mathbf{r}_{1}\mathbf{,r}_{2}\right)
-D_{2}^{1}\left(  \mathbf{r}_{1}\mathbf{,r}_{2}\right)  \right\vert
d\mathbf{r}_{1}d\mathbf{r}_{2}=0\equiv D_{1}^{1}\left(  \mathbf{r}%
_{1}\mathbf{,r}_{2}\right)  \overset{a.e}{=}D_{2}^{1}\left(  \mathbf{r}%
_{1}\mathbf{,r}_{2}\right)
\end{equation}
so in particular
\begin{equation}
\int\left\{  \rho_{1}\left(  \mathbf{r}\right)  -\rho_{2}\left(
\mathbf{r}\right)  \right\}  d\mathbf{r}=0\equiv\rho_{1}\overset{a.e}{=}%
\rho_{2}.
\end{equation}


\section{Spin\label{Spin}}

\section{Symmetric Function Relationships}

\bibliographystyle{apsrev}
\bibliography{agpopt,bw,novel,San2008,coleman,Coleman2,bibrefs,lowdin-per-olov,Janak}

\end{document}